%!TEX root = /Users/kbj/Dropbox/*classes/601book/601master.tex
\chapter{Features}
\label{ch:features}
\setcounter{exx}{0}
% \setcounter{page}{137}

\noindent
One problem our grammar suffers from presently has to do with the fact that the heads of sentences are often not words. Many finite clauses are headed by bound morphology, as in~(\ref{ex4:1}).
\begin{exe}
\ex \label{ex4:1}
\small
\begin{forest}
[IP 
	[DP [she, roof]]
	[\ibar{I} 
		[\obar{I} [pres] ]
		[VP [\ibar{V}
			[\obar{V} [enjoy] ]
			[DP [chocolate, roof] ]
		] ]
	]
]
\end{forest}
\end{exe}
We reached this conclusion because it is morphology of this sort---agreement/tense morphology---that is in the one-to-one relation with the family of strings that we call IP. What we need to understand is how it is that this bound morpheme manages to find itself expressed on the verb that follows. 

There is some evidence that in certain cases of this sort, the verb moves to the position our rules put the inflectional morphology. The rule responsible is, therefore, sometimes called \textsc{Verb Movement}. Its classic description is found in \citet{Chomsky:1957}.

\section{The Classic account of Verbal Inflection}

Chomsky uses a type of reasoning that I will dub a ``Correlation Argument,'' because it's based on the correlation that holds between a verb's syntactic position and its inflection. The generalization that accounts like his hope to derive is the dependence between a verb's syntactic position on its inflectional class. This dependence can be seen from the relative positions of auxiliary verbs in English finite clauses. An auxiliary verb cannot follow \emph{not,} \emph{so}, \emph{too} (which I'll call ``polarity terms'') if it is inflected with the tense/subject-agreement morpheme. We find paradigms like those in~(\ref{ex4:2}).
\begin{exe}
	\ex \label{ex4:2} 
	\begin{xlist}
	\ex[]{Andre has not been eating.}
	\ex[]{Andre is not eating.}
	\ex[*]{Andre not has been eating.}
	\ex[*]{Andre not is eating.}
	\ex[]{Andre has too/so been eating.}
	\ex[]{Andre is too/so eating.}
	\ex[*]{Andre too/so has been eating.}
	\ex[*]{Andre too/so is eating.}
	\end{xlist}
\end{exe}
But when an auxiliary verb is not inflected with tense/subject-agreement morphology it can follow these terms.
\begin{exe}
	\ex \label{ex4:2a}
	\begin{xlist}
	\ex[]{Andre must not have eaten.}
	\ex[]{Andre must not be eating.}
	\ex[]{Andre will too/so have eaten.}
	\ex[]{Andre will too/so be eating.}
	\end{xlist} 
\end{exe}
In these contexts we find that the polarity item must be found immediately before the auxiliary verb; it cannot precede the contents of \obar{I}, the modal, for instance.
\begin{exe}
	\ex \label{ex4:2b}
	\begin{xlist}
	\ex[*]{Andre not must have eaten.}
	\ex[*]{Andre not must be eating.}
	\ex[*]{Andre must have not eaten.}
	\ex[*]{Andre must be not eating.}
	\ex[*]{Andre too/so will have eaten.}
	\ex[*]{Andre too/so will be eating.}
	\ex[*]{Andre will have too/so eaten.}
	\ex[*]{Andre will be too/so eating.}
	\end{xlist} 
\end{exe}
Some English speakers may judge (\ref{ex4:2b}c,d) grammatical. For these speakers, I suspect that \emph{not} is being used differently than it is in situations where \emph{not} appears between the modal and auxiliary, as in~(\ref{ex4:mana}).
\begin{exe}
	\ex \label{ex4:mana}
	\begin{xlist}
	\ex[]{Andre will not have eaten.}
	\ex[]{Andre will not be eating.}
	\end{xlist} 
\end{exe}
In (\ref{ex4:mana}), \emph{not} negates the entire sentence. The speaker of these sentences means to deny the positive versions of these sentences, i.e. the sentences in~(\ref{ex4:manb}).
\begin{exe}
	\ex \label{ex4:manb}
	\begin{xlist}
	\ex[]{Andre will have eaten.}
	\ex[]{Andre will be eating.}
	\end{xlist} 
\end{exe}
By contrast, the sentences in (\ref{ex4:2}c) and (\ref{ex4:2}d) are positive assertions. The speaker of these sentences means to assert that Andre must have done something---namely not eaten---or must be doing something---namely not eating. This difference between (\ref{ex4:2}c,d) and (\ref{ex4:mana}) can be brought out by considering how they combine with so-called question ``tags.'' These tags, illustrated in (\ref{ex4:tag1}), have the opposite polarity as the clause they are appended to.
\begin{exe}
	\ex \label{ex4:tag1}
	\begin{xlist}
	\ex[]{Andre will eat, won't he?}
	\ex[]{Andre will not eat, will he?}
	\end{xlist} 
\end{exe}
If the polarity of the tag is the same as the clause it's appended to, the result has a rhetorical flavor:
\begin{exe}
	\ex \label{ex4:tag2}
	\begin{xlist}
	\ex[]{Andre will eat, will he?}
	\ex[]{Andre will not eat, won't he?}
	\end{xlist}
\end{exe}
When \emph{not} stands between the modal and a following auxiliary, adding a positive question tag gives a question of the sort illustrated by (\ref{ex4:tag1}), as in~(\ref{ex4:tag3}).
\begin{exe}
	\ex \label{ex4:tag3} Andre will not have eaten, will he (have)?
\end{exe}
But when \emph{not} follows both modal and auxiliary, adding a positive question tag yields the rhetorical sorts of questions in (\ref{ex4:tag2}), as in~(\ref{ex4:tag4}).
\begin{exe}
	\ex \label{ex4:tag4} Andre will have not eaten, will he (have)?
\end{exe}
On the other hand, when the auxiliary is finite, and so in the \obar{I} position, a \emph{not} that follows it behaves the same as the one that follows a modal:
\begin{exe}
	\ex \label{ex4:tag5} Andre has not eaten, has he? 
\end{exe}
To the extent that there are real judgments here---it's not easy to be clear about what the difference between the ``real'' and ``rhetorical'' questions is---it suggests that we do not have the same use of \emph{not} when it stands before and after a non-finite auxiliary. If these uses of \emph{not} correspond to the syntactic position it occupies, then we want to conclude that \emph{not} is in the same position in (\ref{ex4:tag6}), but a different position in~(\ref{ex4:tag7}).
\begin{exe}
	\ex \label{ex4:tag6}
	\begin{xlist}
	\ex[]{Andre will not have eaten.}
	\ex[]{Andre has not eaten.}
	\end{xlist}
	\ex \label{ex4:tag7} Andre will have not eaten.
\end{exe}

These data, then, lead us to conclude that polarity items lie between \obar{I} and VP. And from this we can conclude from the contrast between (\ref{ex4:2}) and (\ref{ex4:2a}) that verbs are in \obar{I} when they bear the inflection that resides there. There is, in other words, a correlation between a verb's inflectional class and its syntactic position.

That auxiliary verbs have a different position when they are inflected is confirmed by the action of rules that affect VP. These rules seem to necessarily strand the verb inflected for agr/tense. This is indicated by the following examples of VP deletion and VP topicalization.
\pagebreak
\begin{exe}
	\raggedright
	\ex \label{ex4:10} 
	\begin{xlist}
	\ex[]{Aynat is eating pickles because Emily is $\triangle$.}
	\ex[]{Aynat should be eating pickles because Emily should $\triangle$.}
	\ex[*]{\label{ex4:10vpe}Aynat is eating pickles because Emily $\triangle$.}
	\ex[]{I claimed that Aynat is eating pickles, and [\VP eating pickles] she~is.}
	\ex[]{I claimed that Aynat has to be eating pickles and [\VP be eating pickles] she has to.}
	\ex[*]{\label{ex4:10vpm}I claimed that Aynat is eating pickles, and [is eating pickles],~she.}
	\end{xlist}
\end{exe}
If VP Ellipsis is capable of eliding any VP, then the ungrammaticality of (\ref{ex4:10vpe}) indicates that the finite form of \emph{be} is not within a VP. Similarly, (\ref{ex4:10vpm}) indicates that the movement rule responsible for moving a VP at the beginning of a clause cannot relocate a finite form of \emph{be}. This would follow if finite forms of \emph{be} are not within the VP that is being moved in these constructions.

A Verb Movement account of these phenomena lets phrase structure rules fix the location of the inflectional morphemes involved, and moves the appropriate verb to this location by way of a transformational rule. Thus, simple finite clauses are assigned a D-Structure representation like that in (\ref{ex4:4}), and clauses with a complex tense (i.e., an auxiliary verb), have D-Structure representations like those in~(\ref{ex4:5}).
\begin{exe}
\ex \label{ex4:4} \vspace*{-2pt}
\small
\begin{forest}
[\ibar{I}
	[\obar{I} [\emph{tns/agr}\\ \emph{modal}] ]
	[NegP	[\ibar{Neg}
		[\obar{Neg} \emph{not}\\\emph{so} ]
		[VP [\ibar{V}
			[\obar{V} ]
			[YP ]
		] ]
	] ]
]
\end{forest}
\vspace*{8pt}
\ex \label{ex4:5}\marginnote{In the previous chapter we saw that there is evidence that verbs which have an external θ-role are accompanied with a ``\obar{v}'' that expresses that $\theta$-role. One of our tasks in this chapter is to understand how to capture the correlation between the $\theta$-role in \obar{v} and the verb that follows. We take up this task in 4.6. Until then, I will ignore the possible presence of \obar{v} in the parses, and revert to parses in which VPs involve \obar{V} and not \obar{v}.}
a. 
\small
\begin{forest}
[\ibar{V}
	[\obar{V}  [have]]
	[IP [\ibar{I}
		[\obar{I} [\emph{en}] ]
		[VP [\ibar{V}
			[\obar{V} ]
			[YP ]
		] ]
	]]
]
\end{forest}
\hspace*{12pt}
b. 
\small
\begin{forest}
[\ibar{V}
	[\obar{V} [be]]
	[IP [\ibar{I}
		[\obar{I} [\emph{ing}\\\emph{en}] ]
		[VP [\ibar{V}
			[\obar{V} ]
			[YP ]
		] ]
	]]
]
\end{forest}
\end{exe}
In these parses I've decided to position polarity items as heads that embed the following VP. Moreover, I've renamed these items, and the phrase they head, ``Negation.'' This renaming is common in the literature, as it is typical that only the negative item among the polarity terms is talked about. The representation in which \obar{Neg} embeds VP is argued for in \citet{Pollock:1989}, and we will consider his reasons soon. There are other possibilities which are consistent with what we've seen so far, however. One is to treat polarity items as heading a phrase that is a left-linearized daughter of the highest \ibar{V}, like an adverbial modifier can be.

The representations in (\ref{ex4:5}) claims that auxiliary verbs select IPs as their complements. This gives us a way of capturing the fact that auxiliary verbs determine the inflection that is borne by the VPs they embed. The auxiliary verb \emph{have}, for instance, determines that the VP it embeds will be inflected as what's called a ``perfect participle'' in traditional grammars. These forms frequently look very much like the simple past form of a verb: \emph{heard}, \emph{liked}, and \textit{dreamed} can be either perfect participles or past tense. But not all perfect forms look like simple pasts, for example: \emph{taken} is the perfect of \textit{take}, and \textit{took} is the past; \emph{eaten} is the perfect, and \textit{ate} the past. If we imagine that auxiliary \emph{have} selects an Inflectional Phrase headed by the morphology that is responsible for inflecting a verb as a perfect participle, we can derive that the perfect form of a verb follows auxiliary \textit{have}. Similarly, there is an auxiliary verb \emph{be} that causes the VP it is embeds to be headed by a verb in the so-called ``progressive'' or ``present participle'' inflectional form, for example: \emph{eating}, \emph{running}, etc. Another auxiliary verb \emph{be} embeds the passive participle, which we've encountered in our discussion of the passive rule. For all three of these auxiliary verbs, we might speculate that they select an Inflectional Phrase headed by the morphology responsible for inflecting the embedded verb.

The Verb Movement rule could be formulated with~(\ref{ex4:vm1}).
\begin{exe}
\ex \label{ex4:vm1} Verb Movement
\begin{xlist}
	\ex Remove a verb, $\alpha$, from the position it occupies and
	\ex Adjoin $\alpha$ to \obar{I}.
\end{xlist}
\end{exe}
This will produce from the D-structures in (\ref{ex4:4}) and (\ref{ex4:5}) representations like those in~(\ref{ex4:7}).
\begin{exe}
\ex \label{ex4:7} 
a.
\small
\begin{forest}
[IP
	[XP ]
	[\ibar{I} 
		[\obar{I} 
			[\obar{V},name=binder ]
			[\obar{I} [\emph{tns/agr}] ]
		 ]
		[NegP [\ibar{Neg}
			[\obar{Neg} [\emph{not}\\ \emph{so}] ]
			[VP [\ibar{V}
				[{\ }, name=trace ]
				[YP ]
			] ]
		] ]
	]
]
\draw[->,>=stealth', ultra thin, gray] (trace) to[out=south west,in=south,looseness=1] node[pos=0.75,xshift=-10mm]{} (binder);
\end{forest}
b.
\begin{forest}
[VP
	[\obar{V} [\emph{have}\\\emph{be}]]
	[IP [\ibar{I}
		[\obar{I}
			[\obar{V}, name=binder ]
			[\obar{I} [\emph{en}\\\emph{ing}] ]
		 ]
		[VP [\ibar{V}
			[{\ }, name=trace ]
			[YP ]
		]]
	]]
]
\draw[->,>=stealth', ultra thin, gray] (trace) to[out=-110,in=south,looseness=2.1] node[pos=0.75,xshift=-10mm]{} (binder);
\end{forest}
\end{exe}
This gives a straightforward account of the correlation described above. Verbs stand in the positions determined by the inflections they bear because they are moved to the positions that the phrase structure rules assign to these affixes.

To force Verb Movement in these contexts, we could posit a requirement that an affix must share an \obar{X} with its stem at S-structure. \citet{Lasnik:1981}, where this idea is proposed, calls this condition the ``stray affix filter.''
\begin{exe}
\ex \label{ex4:stray} Stray Affix Filter\\
An affix must co-exist with its stem under a common \obar{X} at S-Structure.
\end{exe}
To prevent Verb Movement in contexts where the \obar{I} contains a word, such as a modal, we could strengthen the Stray Affix Filter to~(\ref{ex4:wcrit}).
\begin{exe}
\ex \label{ex4:wcrit} The Word Criterion\\
A maximal \obar{X} must be parsed as one complete word at S-structure. $\alpha$ a maximal \obar{X} only if it is immediately dominated by \ibar{X}.
\end{exe}
The Word Criterion, then, will force Verb Movement when \obar{I} contains something less than a full word, and block Verb Movement in cases such as (\ref{ex4:toobig}), where the result would put into \obar{I} more than one word.
\begin{exe}
\ex[*]{\label{ex4:toobig} Jill be must eating apples.}
\sn
\small
\begin{forest}
[IP
	[DP [Jill, roof]]
	[\ibar{I} 
		[\obar{I}
			[\obar{V} [be,name=binder] ]
			[\obar{I} [must] ]
		 ]
		[VP [\ibar{V}
			[{\ }, name=trace ]
			[DP [apples, roof] ]
		] ]
	]
]
\draw[->,>=stealth', ultra thin, gray] (trace) to[out=south west,in=south east,looseness=1] node[pos=0.75,xshift=-10mm]{} (binder);
\end{forest}
\end{exe}

What this case shows us is that there is more to providing a syntactic representation with words than populating the \obar{X} positions in a D-structure. It looks like words can be formed by derivations. In the ase of auxiliary verbs in English, for instance, it is Verb Movement that brings the two morphological pieces together. What is needed, then, is to distinguish the morphemes that are inserted at D-structure from the actual words that make up the string that derivations result in. We will continue to let Lexical Insertion be a process that forms D-structures, but we will understand the things that Lexical Insertion inserts to include the pieces that make up a word. Let's call the things that Lexical Insertion matches to \obar{X} positions \emph{morphemes}.
\begin{exe}
	\ex \label{ex4:LI} Lexical Insertion\\
	Insert morphemes into \obar{X} position.
\end{exe}
Lexical Insertion, then, can populate a D-structure with the morphemes that a derivation can assemble into a word. In addition to Lexical Insertion, we'll need a process that matches words to the positions which the syntax has assembled the morphemes in.\marginnote[-20pt]{This method of letting the syntax control how sub-word parts are assembled is the image that \citet{Harris:1946} had and is associated now with what is known as ``Distributed Morphology.'' See \citet{Halle:1993b} and \citet{Marantz:1996}.} This process is sometimes called \emph{vocabulary insertion}.
\begin{exe}
	\ex \label{ex4:vocabI} Vocabulary Insertion\\
	Insert a vocabulary item into an \obar{X} containing the morphemes it expresses.
\end{exe}
We'll start with the idea that Vocabulary Insertion happens at S-structure. We can view the Word Criterion as a result of the conditions on Vocabulary Insertion. It will be responsible for inserting an inflected auxiliary verb into an \obar{I} position that carries both the auxiliary and the inflection that makes it tensed.

These are the cases, then, where it appears that the morphology in \obar{I} and the verb combine by way of movement. They all involve the auxiliary verbs \emph{have} and \emph{be}. But this process doesn't seem to be involved for English ``main'' verbs. In these cases, the verbs do not occupy the positions determined by the inflections they bear, but remain instead in the location given to main verbs by the phrase structure rules. Examples like (\ref{ex4:6}) are ungrammatical.
\begin{exe}
	\ex \label{ex4:6} *Andre likes not/too/so apples.
\end{exe}
And the processes that indicate whether a verb is in its VP, i.e., VP ellipsis and VP fronting, indicate that finite main verbs are indeed positioned within the VPs they head.
\begin{exe}
\ex \label{ex4:mvnom} 
\begin{xlist}
	\ex[*]{Andre put apples on the table because Jill put $\triangle$.}
	\ex[*]{I said that Andre put apples on the table and [\VP apples on the table ] Andre put.}
\end{xlist}
\end{exe}
The process that brings inflection together with the verb on which it is expressed must have a different outcome when main verbs are involved. The challenge is to find a way of conceiving of this process so that both kinds of outcomes are possible. There have been a number of attempts to capture this flexibility, none of them entirely successful.

One account, found in \citet{Chomsky:1957}, posits a transformational rule that moves the inflection down onto the verbal root when main verbs are involved. A sentence like (\ref{ex4:affhop}), for example, would enter into the derivation in~(\ref{ex4:affhopd}).
\pagebreak
\begin{exe}
\ex \label{ex4:affhop} Andre likes apples.
\ex \label{ex4:affhopd}
\small
\begin{forest}
[IP
	[DP [Andre, roof]]
	[\ibar{I} 
		[{\ }, name=trace ]
		[VP [\ibar{V}
			[\obar{V}
				[\obar{V} [like] ]
				[\obar{I},name=binder [pres] ]
			 ]
			[DP [apples, roof] ]
		] ]
	]
]
\draw[->,>=stealth', ultra thin, gray] (trace) to[out=south east,in=north east,looseness=1] node[pos=0.75,xshift=10mm]{} (binder);
\end{forest}
\end{exe}
On this view, then, verbs and their inflections always combine by way of a movement rule. The difference between auxiliary and main verbs concerns what moves: the verb, in cases involving an auxiliary, and the inflection, in all other cases.

A problem with this proposal is that it requires movement transformations to be able to move things downwards. We have seen that for other movement transformations this does not seem to be possible. Outside of these contexts where main verbs get inflected without moving, the same generalization holds with respect to the rule that moves verbs around. It looks like it is worth preserving the generalization that movement rules do not move things to downwards. This can be seen, for example, in~(\ref{ex4:downv}).
\begin{exe}
\ex \label{ex4:downv}
\small
\begin{forest}
[IP
	[CP [\ibar{C}
		[\obar{C} [that] ]
		[IP
			[DP [Anne Michelle, roof]]
			[\ibar{I}
				[\obar{I}
					[\obar{V} [be, name=binder] ]
					[\obar{I} [past] ]
				 ]
				[VP [leave, roof] ]
			 ]
		 ]
	]]
	[\ibar{I} 
		[\obar{I} [must] ]
		[VP [\ibar{V}
			[{\ }, name=trace ]
			[IP [bothering her, roof] ]
		] ]
	]
]
\draw[->,>=stealth', ultra thin, gray] (trace) to[out=south,in=south east,looseness=1.5] node[pos=0.75,xshift=-10mm]{} (binder);
\end{forest}
\vspace*{-12pt}
\sn = That Anne Michelle was leave must bothering her.
\end{exe}
In (\ref{ex4:downv}), a verb, \emph{be}, has moved into an \obar{I} that is not ``higher'' than where that \emph{be} starts out. This is the kind of movement that we saw was blocked for the other examples of movement: Argument Movement, Wh Movement, and Complex NP Shift.
\pagebreak 

To block downwards movement, we formulated the constraint in~(\ref{ex4:upwardold}).
\begin{exe}
	\raggedright
	\ex \label{ex4:upwardold} Upwards\\
	For $\alpha$ to move to position $\beta$, $\beta$ must c-command $\alpha$.
\end{exe}
This condition will correctly block (\ref{ex4:downv}), but it would also block Chomsky's Affix Hopping rule, and this is the problem faced by Chomsky's proposal. Interestingly, it would also block \underline{good} examples of Verb Movement. This can be seen by looking at the structure formed by Verb Movement in~(\ref{ex4:7}a), repeated here.
\begin{exe}
	\raggedright
	\exr{ex4:7} 
	\small
\begin{forest}
[IP
	[XP ]
	[\ibar{I} 
		[\obar{I} 
			[\obar{V},name=binder ]
			[\obar{I} [\emph{tns/agr}] ]
		 ]
		[NegP [\ibar{Neg}
			[\obar{Neg} [\emph{not}\\ \emph{so}] ]
			[VP [\ibar{V}
				[{\ }, name=trace ]
				[YP ]
			] ]
		] ]
	]
]
\draw[->,>=stealth', ultra thin, gray] (trace) to[out=south west,in=south,looseness=1] node[pos=0.75,xshift=-10mm]{} (binder);
\end{forest}
\end{exe}

If we are to use Upwards to prevent (\ref{ex4:downv}), then there is a problem with the way Upwards is conceived. There are several ways of capturing the idea that Verb Movement shares with the other movement rules the property that it is always ``upwards.'' Some of those avenues involve preserving the formulation of Upwards in (\ref{ex4:upwardold}), and changing either how c-command is defined or changing the structures that result from Head Movement.\marginnote[0pt]{See \citealt{Matushansky:2006} for some discussion of these alternatives.} None of these revisions, however, make Affix Hopping possible, and for this reason the proposal in \citet{Chomsky:1957} has been largely abandoned. In these notes, I will preserve the definition of c-command and also keep the (admittedly somewhat \emph{ad hoc}) decision about the structure that Verb Movement creates. To bring Verb Movement in line with Upwards, I shall redefine Upwards.\marginnote[19pt]{By ``projection,'' is meant the relation between heads and their phrases that the \ibar{X} Skeleton defines. \ibar{X} is a projection of its head \obar{X}, and XP is a projection of them both.} 
\begin{exe}
	\raggedright
	\ex \label{ex4:upwards} Upwards\\
	If $\alpha$ moves from position T to position P, then the sister of P must have a projection that dominates~T.
\end{exe}
In (\ref{ex4:7}), the sister of the moved \obar{V} is \obar{I}. The projections of \obar{I} include the \obar{I} immediately dominating it, the \ibar{I} dominating that \obar{I} and, finally, IP. Both IP and \ibar{I} dominate the position that this moved \obar{V} came from. For this reason, Verb Movement in this scenario satisfies the formulation of Upwards in (\ref{ex4:upwards}). By contrast, (\ref{ex4:downv}) violates Upwards. There is no projection of the sister to \emph{be} in the position it occupies after moving that dominates the position \emph{be} moved from. The definition of Upwards in (\ref{ex4:upwards}) applies to the other cases of movement---Argument Movement, Complex NP Shift, and Wh Movement---we've seen in the same way that the earlier formulation of Upwards did. It also prevents Affix Hopping. We cannot use Affix Hopping to bring tense inflection and main verbs together in English.

A second account lets the mechanism that brings a main verb together with its inflection be a modification of how ``lexical insertion'' works. Pre\-sent\-ly, our formulation of lexical insertion matches morphemes to a single \obar{X} positions at D-structure by inserting them into the \obar{X} positions. One word cannot, consequently, span more than one terminal position in a phrase marker. Jonathan Bobaljik\cite{Bobaljik:1994a,Bobaljik:1995a} argues that this should be weakened so that a lexical item can be put into a position that does not map directly onto the one position that a terminal position in a phrase marker occupies. In particular, he proposes to let an inflected form of a main verb be inserted into a position occupied by two \obar{X}s, just in case those two \obar{X}s are linearly adjacent. This account requires a different architecture than the one we have adopted, since ``adjacency'' is not a relation that can be evaluated at D-structure. It is only after the linearization algorithm has interpreted S-structures that it will be known which positions are adjacent to another. An implementation of Bobaljik's proposal, then, requires some changes to our present architecture. Under his proposal, a representation like (\ref{ex4:bob1}) could be directly mapped onto the string \emph{Andre likes apples}.
\begin{exe}
\ex \label{ex4:bob1} 
\small
\begin{forest}
[IP
	[DP [Andre, roof]]
	[\ibar{I} 
		[\obar{I} [pres] ]
		[VP [\ibar{V}
			[\obar{V} [like] ]
			[DP [apples, roof] ]
		] ]
	]
]
\end{forest}
\end{exe}
This proposal has the benefit of explaining why finite sentences containing polarity items in English are ungrammatical if they do not contain an auxiliary verb.
\begin{exe}
\ex \label{ex4:notfin} 
\begin{xlist}
	\ex[*]{Andre not likes apples.}
	\ex[*]{Jill not ran down the street.}
	\ex[*]{Winnie so read that encyclopedia.}
	\ex[*]{Bart too ran into the street.}
\end{xlist}
\end{exe}
Because polarity terms are linearized between \obar{I} and the verb that follows, it will prevent an inflectional morpheme residing in \obar{I} from being adjacent to that verb and block lexical insertion from rendering them as a single inflected verb. Indeed, this proposal predicts that a finite clause with no auxiliary verbs in it should not allow any material to stand between \obar{I} and the main verb. The only grammatical parse for an example like (\ref{ex4:completely}) \vpageref[below]{ex4:completely} should be one in which the adverb is to the left of \obar{I}, as indicated.
\begin{exe}
\ex \label{ex4:completely} Jill completely finished the apple.
\sn
\small
\begin{forest}
[IP
	[DP [Jill, roof]]
	[\ibar{I} 
		[AdvP [completely, roof]]
		[\ibar{I}
			[\obar{I} [past] ]
			[VP [\ibar{V} 
				[\obar{V} [finish] ]
				[DP [the apple, roof] ]
			] ]
		 ]
	]
]
\end{forest}
\end{exe}

This is not an obviously correct prediction, however. Contrasts like those in (\ref{ex4:completely2}) indicate that \emph{completely} must be placed between \obar{I} and the following~VP.
\begin{exe}
\ex \label{ex4:completely2} 
\begin{xlist}
	\ex[]{Jill will completely finish the apple.}
	\ex[*]{Jill completely will finish the apple.}
\end{xlist}
\end{exe}
This account, then, has difficulties making sense of data like (\ref{ex4:completely}). Unlike Chomsky's proposal, it does address the problem of why the mechanism that inflects main verbs is sensitive to the presence of polarity items; but it doesn't explain how these items differ from other material that appears to be able to stand between the main verb and finite inflection.

The final idea we will consider makes the relationship between an \obar{I} that is associated with inflectional morphology and the verb that expresses that morphology to be akin to the relationship we have posited between a DP and the term that assigns that DP Case. We can think of \obar{I} as assigning the relevant inflectional morphology to the verb that follows, as in~(\ref{ex4:agree1}).
\begin{exe}
\ex \label{ex4:agree1} 
\small
\begin{forest}
[IP
	[DP [Andre, roof]]
	[\ibar{I} 
		[\obar{I} [finite, name=trace] ]
		[VP [\ibar{V}
			[\obar{V} [likes, name=binder] ]
			[DP [apples, roof] ]
		] ]
	]
]
\draw[->,>=stealth', ultra thin, gray] (trace) to[out=south west,in=west,looseness=1] node[pos=0.5,xshift=-6mm]{agr/tns} (binder);
\end{forest}
\end{exe}
\citet{Chomsky:1995}, where this proposal can be found, calls the assignment relationship between an inflected verb and the \obar{I} where that inflection is determined: \textsc{agree}. The precise formulation of \textsc{agree} will be taken up later; but it should be clear from (\ref{ex4:agree1}) that it doesn't match perfectly the condition we arrived at for Case assignment in the previous chapter. While the finite \obar{I} does c-command the verb it assigns tense/agreement inflection to in (\ref{ex4:agree1}), it is not klose to that verb: there is more than one phrase that dominates the verb but not \obar{I}. This would block Case assignment. One desirable consequence of the formulation of \textsc{agree} is to make it account for the inability of a polarity item to stand between \obar{I} and a main verb. That would be a natural place for this account to attempt an explanation for this fact.

If finite \obar{I} assigns inflectional morphology to a main verb that follows, then there are a couple ways of characterizing the different outcome with auxiliary verbs. One possibility is that auxiliary verbs are inflected in the same way, and that there is an independent reason that they must move to the \obar{I} that assigns them their inflection. Another possibility is that auxiliary verbs are prevented from getting their inflectional morphology by way of \textsc{agree}, and instead make recourse to Verb Movement. Hopefully, a better understanding of what makes an auxiliary verb different from a main verb will help us understand how to express the difference in how their inflection influences their surface position. We can learn something about this question by considering how these processes arise in other languages. We'll look at other languages in the sections that follow, and then return to the question of how to give a picture of the main verb/auxiliary verb distinction in English.

We should also note here the special status of the ``auxiliary'' \emph{do}. This provides a way of speaking the bound tense morphemes in English alongside whatever mechanism allows these morphemes to be pronounced on the main verb that follows.\marginnote[-20pt]{There is a main verb \emph{do} in English as well. It is found in sentences like \emph{Mary will do this}. The thing I'm calling an auxiliary \emph{do} doesn't have the distribution of a verb, so it is not the same kind of thing that the auxiliary verbs \emph{have} and \emph{be} are. The auxiliary \emph{do}, for instance, is not found in infinitival clauses, like \emph{have} or \emph{be} are. There are contrasts between examples like \emph{*I will try to do eat nattoo.} and \emph{I will try to have eaten nattoo.} It appears that \emph{do} only shows up in \obar{I}.} In general, we find \emph{do} in contexts where the tense morphemes cannot otherwise find a way of being expressed. For instance, \emph{do} appears in contexts where there is no auxiliary verb, but there is the presence of \emph{not}, which blocks the mechanism by which main verbs inflect.
\begin{exe}
	\raggedright
	\ex \label{ex4:do1} 
	\begin{xlist}
	\ex Smith did not eat nattoo.
	\ex Smith did so eat nattoo.
	\ex Smith did too eat nattoo.
	\end{xlist}
\end{exe}
If \emph{not} isn't present, however, \emph{do} cannot be employed:\marginnote[0pt]{These examples are grammatical if there is focus on \emph{did}. We might imagine that this is the case because to focus-mark \obar{I} requires a way of pronouncing \obar{I} independently of the verb.}
\begin{exe}
	\raggedright
	\ex \label{ex4:do2} 
	\begin{xlist}
	\ex[*]{Smith did eat nattoo.}
	\ex[*]{Smith did eat nattoo.}
	\ex[*]{Smith did eat nattoo.}
	\end{xlist}
\end{exe}
And \emph{do} is not present when there is an auxiliary verb that can move to form a word with \obar{I}.
\begin{exe}
	\raggedright
	\ex \label{ex4:do3} 
	\begin{xlist}
		\ex[*]{Smith did have eaten nattoo.} 
		\ex[*]{Smith did be eating nattoo.}
	\end{xlist}
\end{exe}
A way of describing this distribution is to say that \emph{do} appears when no other way of pronouncing the bound morphemes in \obar{I} are possible. Finding a way of expressing that condition isn't possible in the syntax we presently have. There are proposals in the literature about how to achieve that, but to look at these now would take us too far off track. I will ignore the distribution of auxiliary verb \emph{do} in these lectures.

We will, however, look at how to model the competing ways in which auxiliary and main verbs in English inflect. That is best done, as mentioned above, by looking at how these systems of inflecting verbs vary cross-linguistically. But first, there is one last piece to the picture of English verb placement that we should see.

\section{Movement to \obar{C}}
\label{sec:movetoC}

The rules we've looked at that relate verbs to \obar{I} can be embellished with another that plays a role in forming questions. One sort of question involves relocating the finite verb to the beginning of the sentence, as in~(\ref{ex4:12}).
\begin{exe}
	\ex \label{ex4:12}
	\begin{xlist}
	\ex[]{Have you eaten pickles?}
	\ex[]{Should you eat pickles?}
	\end{xlist}
\end{exe}
These are called Yes/No questions. Another sort of question which involves a relocation of the finite verb seeks more specific information and is called a Wh-Question. In these questions, as we saw in the previous chapter, there is also movement of a phrase into Specifier of CP. Some examples are in~(\ref{ex4:13}).
\begin{exe}
	\ex \label{ex4:13} 
	\begin{xlist}
	\ex[]{Which pickles have you eaten?}
	\ex[]{What should you eat?}
	\end{xlist}
\end{exe}
In both of these types of questions, the contents of \obar{I} has moved to some sentence-initial position. As a consequence of this rule's targeting \obar{I}, only material which can stand in \obar{I} can move. Therefore, main verbs in English do not undergo this process, and instead the nearly meaningless auxiliary verb \emph{do} is used.
\begin{exe}
	\ex \label{ex4:14} 
	\begin{xlist}
	\ex[*]{Eat you pickles?}
	\sn[*]{Which pickles eat you?}
	\ex[]{Do you eat pickles?\\ Which pickles do you eat?}
	\end{xlist}
\end{exe}

Where does \obar{I} move to? The common answer to this question, based on work of the Dutch linguists Jan Koster\cite{Koster:1975} and Hans den Besten,\cite{denBesten:1983} is that it is \obar{C}. The reason for this speculation is that the rule involved here seems to be in complementary distribution with complementizers. That is, its effects are blocked in those situations where we believe that complementizers are present in this position. Thus, there is a distinction between~(\ref{ex4:15}a) and~(\ref{ex4:15}b).
\begin{exe}
	\ex \label{ex4:15} 
	\begin{xlist}
	\ex[]{Have you eaten?}
	\ex[*]{I remember (that) have you eaten.}
	\end{xlist}
\end{exe}
The same contrast is found in the position of finite auxiliaries in wh-questions.
\begin{exe}
	\raggedright
	\ex \label{ex4:embIC} 
	\begin{xlist}
		\ex[]{Which pickles have you eaten?}
		\ex[*]{I remember that which pickles have you eaten.}
	\end{xlist}
\end{exe}
Instead, Yes/No questions that are embedded are marked with a special complementizer: \emph{whether}. And wh-questions in embedded contexts involve no evidence of a \obar{C} at all.

Under this view, the rule involved in these questions, then, might be formulated as in~(\ref{ex4:itoc}).%\enlargethispage*{\baselineskip}
\begin{exe}
	\ex \label{ex4:itoc} \textsc{I-to-C Movement}
	\begin{xlist}
		\ex Remove an \obar{I}, $\alpha$, from the position it occupies and
		\ex Adjoin $\alpha$ to \obar{C}.
	\end{xlist} 
\end{exe}
This will produce derivations like~(\ref{ex4:itocd1}).\marginnote[0pt]{Note that this model of questions means that every root sentence that is a question in English must be a CP. By contrast, non-question root sentences in English are~IPs.}
\begin{exe}
	\raggedright
	\ex \label{ex4:itocd1} 
	\small
	\begin{forest}
	[CP[\ibar{C}
		[\obar{C} 
			[\obar{C}
				[\obar{I} [should, name=binder] ]
				[\obar{C} ]
			 ]
		]
		[IP
			[DP [she, roof]]
			[\ibar{I}
				[{\ }, name=trace ]
				[VP [run, roof] ]
			 ]
		 ]
	]]
	\draw[->,>=stealth', ultra thin, gray] (trace) to[out=south west,in=south east,looseness=1] node[pos=0.75,xshift=10mm]{} (binder);
	\end{forest}
\end{exe}
The Word Criterion will correctly block movement of \obar{I} to \obar{C} in those cases where \obar{C} is filled with a complementizer, assuming that a complementizer and the modal or inflected verb in \obar{I} cannot together form one word.

Okay, so this is, roughly, a sketch of the rules that go into making up our knowledge of this fragment of English syntax. We've got three rules, one that moves auxiliary verbs to \obar{I}, another than moves \obar{I} to \obar{C} and a third that inflects main verbs. The first and last are driven by the Stray Affix Filter, or some parallel condition on the view that inflection is assigned under \textsc{agree}, and the second arises in questions. We can make the movement of \obar{I} to \obar{C} obligatory with the Stray Affix Filter if we imagine that there is an inflectional morpheme in the \obar{C} position of questions as well. This morpheme has no phonology associated with it, apparently, but expresses itself instead by forcing movement of \obar{I}. We can think of this morpheme as an affix on words that belong to the part of speech that \obar{I} names, just as tense/agreement inflection is an affix on the words that belong to the part of speech that \obar{V} names. I'll represent this abstract question affix with ``Q.''
\begin{exe}
	\raggedright
	\ex \label{ex4:itocd2} 
	\small
	\begin{forest}
	[CP
		[\obar{C} 
			[\obar{I} [should, name=binder] ]
			[\obar{C} [Q] ]
		]
		[IP 
			[DP [she, roof] ]
			[\ibar{I}
				[{\ }, name=trace ]
				[VP [run, roof] ]
			 ]
		]
	]
	\draw[->,>=stealth', ultra thin, gray] (trace) to[out=south west,in=south east,looseness=1] node[pos=0.75,xshift=10mm]{} (binder);
	\end{forest}
\end{exe}
With this amendment, both \obar{I}-to-\obar{C} movement and Verb Movement are controlled by the Word Criterion, which prevents \obar{I}, \obar{C} or \obar{V} from combining whenever they cannot form a word and forces it when there is an affix that needs a stem to make a complete word.
 
There's one last feature of these rules that we should consider. In all the cases so far examined, the \obar{V} that moves to \obar{I} is always the one that is closest to it. It appears that this is not just an accident of the examples we have chosen. Instead it looks like it is a necessary feature of the verb movement rule. Consider, for example, how these rules might combine to apply to a representation like~(\ref{ex4:16}).
\begin{exe}
	\ex \label{ex4:16} 
	\small
\begin{forest}
[IP
	[\ibar{I}
		[\obar{I} [agr/past]]
		[VP [\ibar{V}
			[\obar{V} [make]]
			[VP
				[DP [it, roof]]
				[\ibar{V}
					[\obar{V} [be]]
					[AP  [obvious that something's wrong, roof]]
				 ]
			 ]
		] ]
	 ]
]
\end{forest}
\end{exe}
Suppose that we move \emph{be} to \emph{agr/past} and form thereby the past tense form: \emph{were}. 
\begin{exe}
\raggedright
	\ex \label{ex4:16a}
	\small
\begin{forest}
[IP
	[\ibar{I}
		[\obar{I}
		[\obar{V} [be, name=binder]]]
		[\obar{I} [agr/past]]
		]
		[VP [\ibar{V}
			[\obar{V} [make]]
			[VP
				[DP [it, roof]]
				[\ibar{V}
					[{\ }, name=trace]
					[AP  [obvious that something's wrong, roof]]
				 ]
			 ]
		] ]
	 ]
\draw[->,>=stealth', thin, gray] (trace) to[out=south west,in=south,looseness=1] node[pos=0.75,xshift=-10mm]{} (binder);
\end{forest}
\end{exe}
But this is an ungrammatical outcome. Instead, the only grammatical outcome is one in which the past tense morphology is expressed on \textit{make}. We need to find a way of preventing derivations that include the movement in (\ref{ex4:16a}). 

Lisa Travis\cite{Travis:1984} suggested that there is a constraint on head movement rules that blocks them from moving past other \obar{X}s. She calls this condition the ``head movement constraint,'' and I will formulate it as (\ref{ex4:hmc}).\marginnote[0pt]{Travis formulates it differently, using the notion of government described in chapter~3. For the cases that we will look at, (\ref{ex4:hmc}) will do.}
\begin{exe}
	\ex \label{ex4:hmc} \textsc{The Head Movement Constraint}\\ 
				If \obar{X} adjoins to \obar{Y}, then there can be no \obar{Z} that c-commands \obar{X} and is c-commanded by~\obar{Y}.
\end{exe}

If the structure of sentences with \emph{not} in them is correct, then the Head Movement Constraint is violated by cases of verb movement to inflect auxiliary verbs, as in~(\ref{ex4:hmcp}).
\begin{exe}
	\raggedright
	\ex \label{ex4:hmcp} 
	\small
	\begin{forest}
	[\ibar{I}
		[\obar{I}
			[\obar{V} [have, name=binder] ]
			[\obar{I} [pres] ]
		 ]
		[NegP [\ibar{Neg} 
			[\obar{Neg} [not] ]
			[VP [\ibar{V}
				[{\ }, name=trace ]
				[IP [eaten cake, roof] ]
			] ]
		] ]
	]
	\draw[->,>=stealth', ultra thin, gray] (trace) to[out=south west,in=south,looseness=1] node[pos=0.75,xshift=-10mm]{} (binder);
	\end{forest}
\end{exe}
This tension between the Head Movement Constraint and this syntax for negation has a payoff, which we will examine later. For the moment, let's bracket this issue and develop a better understanding of the rest of the verb movement system.


\section{Verb Second word order}

Let's now consider verb placement phenomena in languages closely related to English. We'll start with examples from the Germanic family. In German we find that the position of the finite verb depends, roughly speaking, on whether the clause it is in is embedded or not. In embedded clauses, the verbs stack up at the end of the sentence in the mirror image of their arrangement in English: the finite verb comes last.%\enlargethispage*{22pt}
\begin{exe}
	\ex \label{ex4:20} 
	\begin{xlist}
	\ex[]{\gll \ldots da{\ss} Hans das Buch kauft. \\ \ldots that John the book buys\\}
	\ex[]{\gll \ldots da{\ss} Hans das Buch gekauft hat. \\ \ldots that John the book bought has\\}
	\ex[]{\gll \ldots da{\ss} Hans das Buch gekauft haben mu{\ss}. \\ \ldots that John the book bought have must.\\}
	\end{xlist}
\end{exe}
From this we can conclude, perhaps, that VPs are head-final. Further, if there are IPs and a method of combining verbs with the inflection associated with the \obar{I} of these IPs that employs Head Movement, they too are head final in German.

Interestingly, we find a different word-order in root, or independent, clauses. Here, the inflected verb no longer comes finally in the series. Instead it appears immediately following the subject. So unlike (\ref{ex4:20}), we find only the word order in~(\ref{ex4:21}).
\begin{exe}
	\ex \label{ex4:21} 
	\begin{xlist}
	\ex[]{\gll Hans kauft das Buch. \\ John buys the book\\}
	\ex[]{\gll Hans hat das Buch gekauft. \\ John has the book bought\\}
	\ex[]{\gll Hans mu{\ss} das Buch gekauft haben \\ John must the book bought have\\}
	\end{xlist}
\end{exe}
This seems to suggest that there is a movement rule which relocates finite verbs into the position immediately following the subject. Using the logic of the correlation argument, we might imagine that the position where finite inflection is in German immediately follows the subject, and it's into this position that finite verbs are driven in German.

But this would miss the fact that the position of the finite verbs differs for embedded and independent clauses. What we want is some way of forcing verbs to move into the post-subject position in root clauses only. This suggests that it is not the finite distinction that is responsible for verbs' position in root clauses, but something else. Something that distinguishes root from embedded clauses.

We've already seen a similar difference in the grammar of English: recall that \obar{I}-to-\obar{C} movement is restricted to root clauses in English. Perhaps the verbs are moving through \obar{I} into \obar{C} in cases like (\ref{ex4:21}), then. This would credit German and English with two differences. On the one hand there is the difference in headedness that we see most directly in embedded contexts. And then there is something that allows/forces subjects to move past \obar{C} in embedded clauses. We might imagine that this second force, whatever it is, is like the process that moves wh-phrases in English into Specifier of CP. Thus, the examples in (\ref{ex4:21}) might get a parse like that in~(\ref{ex4:22}). The subject moves into Specifier of IP, to satisfy the EPP and get Case, and then into Specifier of~CP.
\begin{exe}
\ex \label{ex4:22}
\small
\begin{forest}
[CP
	[DP [Hans, roof, name=binder]]
	[\ibar{C}
		[\obar{C} ]
		[IP
			[{\ }, name=trace ]
			[\ibar{I} 
				[VP [\ibar{V}
					[vP
						[{\ }, name=trace1 ]
						[\ibar{v} 
							[VP [\ibar{V}
								[DP [das Buch, roof]]
								[\obar{V} [gekauft] ]
							]]
							[\obar{v} [agent] ]
						]
					 ]
					[\obar{V}[haben] ]
				]]
				[\obar{I} [pres]]
			 ]
		 ]
	]
]
\draw[->,>=stealth', ultra thin, gray] (trace1) to[out=west,in=-105,looseness=1.4] node[pos=0.75,xshift=-10mm]{} (trace);
\draw[->,>=stealth', ultra thin, gray] (trace) to[out=south west,in=south,looseness=1] node[pos=0.75,xshift=-10mm]{} (binder);
\end{forest}
\end{exe}
\marginnote[-280pt]{I have put the \obar{I} in this parse to the right of the VP it embeds. We will see evidence that \obar{C} is linearized so that it is to the left of the IP it embeds, and we have just seen evidence that \obar{V} is linearized so that it follows whatever phrase it embeds. German appears to be a language, then, that allows the linearization algorithm to choose different orderings for its heads: CPs (and DPs) are head initial, but VPs are head final. It is not easy to determine how the linearization algorithm linearizes IPs. I have decided, as an initial hypothesis, that \obar{I}s are linearized to the right of their complement. We should be alert to the possibility that the other ordering is correct and seek evidence that helps decide between these two.

I have also made the complement of the auxiliary verb, \emph{haben}, a VP rather than an IP headed by the inflection that causes the embedded verb to be a perfect participle. The relationship between auxiliary verbs in German and the VPs they embed is the same as is found in English, so a first guess should be that the syntax of both languages involves making the complement to an auxiliary verb an IP. I will ignore this, hopefully irrelevant, detail in my parses of German.}
The finite verb follows the subject because it moves through \obar{I} into \obar{C}, as indicated in~(\ref{ex4:22a}).
\begin{exe}
\ex \label{ex4:22a}
\small
\begin{forest}
[CP
	[DP [Hans, roof, name=binder]]
	[\ibar{C} 
		[\obar{C} 
			[\obar{I} ,name=binder2
				[\obar{V} [haben] ]
				[\obar{I} [pres] ]
			 ]
			[\obar{C} ]
		 ]
		[IP 
			[{\ }, name=trace ]
			[\ibar{I} 
				[VP [\ibar{V} 
					[vP
						[{\ }, name=traces ]
						[\ibar{v}
							[VP [\ibar{V}
								[DP [das Buch, roof] ]
								[\obar{V} [gelesen] ]
							] ]
							[\obar{v} [agent] ]
						 ]
					 ]
					[{\ }, name=traceV ]
				] ]
				[{\ }, name=traceI ]
			]
		]
	]	
]
\draw[->,>=stealth', ultra thin, gray] (traceI) to[out=south west,in=south,looseness=1] node[pos=0.75,xshift=-10mm]{} (binder2);
\draw[->,>=stealth', ultra thin, gray] (traceV) to[out=south east,in=south,looseness=1] node[pos=0.75,xshift=10mm]{} (traceI);
\draw[->,>=stealth', ultra thin, gray] (trace) to[out=south,in=south west,looseness=3.6] node[pos=0.75,xshift=-10mm]{} (binder);
\end{forest} 
\end{exe}

This leads to the expectation that we might find other phrases preceding the finite verb; and we do. It is not just the subject that may immediately precede the finite verb in root contexts, any phrase can. When some other phrase comes before the verb, the subject (typically) immediately follows the finite verb and the phrase that shows up at the front is understood to be ``topicalized.'' We can capture these word-orders by letting the finite verb move into \obar{C} in all these clauses, but letting the Specifier of CP be filled by any phrase through a movement rule. Thus, alongside (\ref{ex4:22}) we also find~(\ref{ex4:23}).
\begin{exe}
\ex \label{ex4:23}
\small
\begin{forest}
[CP
	[DP [das Buch, roof, name=binder]]
	[\ibar{C} 
		[\obar{C}
			[\obar{I} , name=binderI
				[\obar{V} [haben] ]
				[\obar{I} [pres] ]
			 ]
			[\obar{C} ]
		 ]
		[IP
			[DP [Hans, roof, name=binder2]]
			[\ibar{I}
				[VP [\ibar{V}
					[vP
						[{\ }, name=trace ]
						[\ibar{v}
							[VP [\ibar{V}
								[{\ }, name=tracet ]
								[\obar{V} [gekauft] ]
							]]
							[\obar{v} [agent] ]
						 ]
					 ]
					[{\ }, name=traceV ]
				]]
				[{\ }, name=traceV2 ]
			 ]
		 ]
	]
]
\draw[->,>=stealth', ultra thin, gray] (traceV) to[out=south east,in=south,looseness=1] node[pos=0.75,xshift=-10mm]{} (traceV2);
\draw[->,>=stealth', ultra thin, gray] (trace) to[out=west,in=south,looseness=1] node[pos=0.75,xshift=10mm]{} (binder2);
\draw[->,>=stealth', ultra thin, gray] (tracet) to[out=west,in=-105,looseness=1.4] node[pos=0.75,xshift=10mm]{} (binder);
\draw[->,>=stealth', ultra thin, gray] (traceV2) to[out=south west,in=south,looseness=1] node[pos=0.75,xshift=-10mm]{} (binderI);
\end{forest}
\end{exe}
The generalization about German word-order can be described this way: any phrase may be initial, but exactly one must be. German is sometimes described as having ``Verb Second'' word-order for this reason. The account we've just sketched of this captures Verb Second by way of a rule that moves the verbs that have moved into finite \obar{I} into \obar{C}, and by moving anything, but always something, into the single position that exists to the left of \obar{C}.

That the verb in these cases has in fact moved into \obar{C} is further substantiated by cases where Verb Second word-order is found in embedded clauses. Though there is considerable dialectal variation here, in the standard dialects, Verb Second is possible in embedded clauses just in those cases where a complementizer may go missing. As in English, it is possible in German to unexpress the complementizer when the CP it heads is the complement to a certain class of verbs. The verb \emph{say}, for example, can go without a complementizer---as in (\ref{ex4:24})---and when it does so in German we find Verb Second word-order, see (\ref{ex4:25}).
\begin{exe}
	\ex \label{ex4:24} Kathryn said (that) Martin has a dime.
	\pagebreak
	\ex \label{ex4:25}
	\begin{xlist}
	\ex[]{\gll Er sagt, da{\ss} die Kinder diesen Film gesehen haben.\\ He says that the kids this film seen have\\ `He says that the kids have seen the film.'}
	\ex[]{\gll Er sagt, diesen Film haben die Kinder gesehen. \\ he says this film have the kids seen\\ `He says that, this film, have the kids seen.'}
	\end{xlist}
\end{exe}

Despite the similarities in verb movement that German and English have, note that one place where they differ is whether main verbs fall under the scope of verb movement. Unlike English, German main verbs can undergo movement. This is seen by their ability to move into \obar{C} in main clauses, as in~(\ref{ex4:21}a).

There are a host of mysteries about the grammar of verb placement in German that we will not attempt to solve here. What, for instance, is responsible for forcing movement of something into the Specifier of CP in German root clauses. And why isn't the movement of a phrase into Specifier of CP allowed in embedded clauses, as in (\ref{ex4:26})?
\begin{exe}
	\ex[*]{\label{ex4:26} Ich sagte [\CP\ das Buch [ da\ss\ [\IP\ Hans [\VP\ gekauft] hat]]].}
\end{exe}
\citet{Vikner:1995} discusses some of the ideas that have been offered for answers to these questions. One thing we can model, however, is the force that requires \obar{I} to move to \obar{C} in root declarative clauses in German. We see a very similar thing happening in English questions: \obar{I} must move to \obar{C} in questions of English, but only when they are root sentences. We captured that fact by imagining that root questions in English were CPs headed by a silent question morpheme. That morpheme, represented by ``Q,'' we can imagine is a suffix on words of category \obar{I}, and the Word Criterion will therefore require it to share an \obar{X} with its root \obar{I} at S-structure. This is what forces movement of \obar{I} to \obar{C} in English root questions. We can give a parallel model to the obligatory movement of \obar{I} to \obar{C} in root declarative sentences in German. Imagine that root sentences of German must be CPs, and that when they are declaratives, the head of that CP is a silent suffix that requires an \obar{I} stem. I'll call this silent suffix ``top.''%
%
\marginnote[0pt]{A common description of the conditions that accompany movement of a phrase into Specifier of CP is that the phrase is ``topicalized.'' I choose ``top'' as the name for the complementizer to evoke this description.}%
%
On this view, the structure in~(\ref{ex4:germwcv}) violates the Word Criterion, because its \obar{C} is not occupied by a complete word.
\begin{exe}
	\raggedright
	\ex \label{ex4:germwcv} 
	\small
	\begin{forest}
	[CP
		[DP [Martin, roof, name=topic]]
		[\ibar{C} 
			[\obar{C} [top] ]
			[IP
				[{\ }, name=trace ]
				[\ibar{I}
					[VP [\ibar{V}
						[vP [das Buch gekauft, roof] ]
						[{\ }, name=tracev ]
					] ]
					[\obar{I}
						[\obar{V} [have, name=binderv] ]
						[\obar{I} [pres] ]
					 ]
				 ]
			 ]
		]
	]
	\draw[->,>=stealth', ultra thin, gray] (tracev) to[out=south east,in=south,looseness=1] node[pos=0.75,xshift=10mm]{} (binderv);
	\draw[->,>=stealth', ultra thin, gray] (trace) to[out=south west,in=south,looseness=1] node[pos=0.75,xshift=-10mm]{} (topic);
	\end{forest}
\end{exe}
This is remedied by moving \obar{I} to \obar{C} to form~(\ref{ex4:germwc}).
\begin{exe}
	\raggedright
	\ex \label{ex4:germwc} 
	\small
	\begin{forest}
	[CP
		[DP [Dan, roof]]
		[\ibar{C} 
			[\obar{C}
				[\obar{I}, name=binder
					[\obar{V} [haben] ]
					[\obar{I} [pres] ]
				 ]
				[\obar{C} [top] ]
			 ]
			[IP
				[{\ } ]
				[\ibar{I}
					[VP [das Buch gekauft, roof]]
					[{\ }, name=trace ]
				 ]
			 ]
		]
	]
	\draw[->,>=stealth', ultra thin, gray] (trace) to[out=south west,in=south,looseness=2.5] node[pos=0.75,xshift=-10mm]{} (binder);
	\end{forest}
\end{exe}

There are some things that we might notice about the verb movement processes in German which are similar to the parallel processes in English. For example, the Head Movement Constraint is preserved in the grammar of verb placement in German. Sentences like the following, where the verb from an embedded VP moves into the higher clause, are ungrammatical.
\begin{exe}
	\ex[*]{\label{ex4:27} \gll Martin schlief den Stubentiger auf dem Bett lassen.\\ Martin slept the cat on the bed let\\ `Martin let the cat sleep on the bed.'}
\end{exe}
The grammatical way of expressing the sentence that underlies (\ref{ex4:27}) is (\ref{ex4:27a}), whose syntax is in (\ref{ex4:27b}). The ungrammatical string in (\ref{ex4:27}) is produced by letting the embedded verb move in violation of the Head Movement Constraint in the way that (\ref{ex4:27c}) indicates.
\begin{exe}
\raggedright
	\ex \label{ex4:27a} \gll Martin l\"ast den Stubentiger auf dem Bett schlafen.\\ Martin let the cat on the bed sleep\\ `Martin let the cat sleep on the bed.'
	\ex \label{ex4:27b}
	\small
	\begin{forest}
	[CP
		[DP [Martin, name=binder, roof]]
		[\ibar{C}
			[\obar{C}
				[\obar{I}, name=binderi
					[\obar{V} [lassen\\ let]]
					[\obar{I} [past]]
				 ]
				[\obar{C} [top]]
			 ]
			[IP
				[{\ }, name=trace ]
				[\ibar{I}
					[VP [\ibar{V}
						[VP
							[DP [den Stubentiger\\ the cat, roof]]
							[\ibar{V}
								[PP [auf dem Bett\\ on the bed, roof]]
								[\ibar{V} [\obar{V} [schlafen\\ sleep]] ]
							 ]
						 ]
						[{\ }, name=tracev ]
					 ]
					[{\ }, name=tracei ]
				 ]
			 ]
		 ]
	]]
	\draw[->,>=stealth', thin, gray] (trace) to[out=south west,in=south,looseness=1] node[pos=0.75,xshift=-10mm]{} (binder);
	\draw[->,>=stealth', thin, gray] (tracev) to[out=south east,in=south,looseness=1] node[pos=0.75,xshift=-10mm]{} (tracei);
	\draw[->,>=stealth', thin, gray] (tracei) to[out=south west,in=south,looseness=1] node[pos=0.75,xshift=-10mm]{} (binderi);
	\end{forest}
	\ex \label{ex4:27c}
	\small
	\begin{forest}
	[CP
		[DP [Martin, name=binder, roof]]
		[\ibar{C}
			[\obar{C}
				[\obar{I}, name=binderi
					[\obar{V} [lassen\\ let]]
					[\obar{I} [past]]
				 ]
				[\obar{C} [top]]
			 ]
			[IP
				[{\ }, name=trace ]
				[\ibar{I}
					[VP [\ibar{V}
						[DP [den Stubentiger\\ the cat, roof]]
						[\ibar{V}
							[PP [auf dem Bett\\ on the bed, roof]]
							[\ibar{V} [{\ }, name=tracev ]]
						 ]
					 ]]
				[{\ }, name=tracei ]
			 ]
		 ]
	]]
	\draw[->,>=stealth', thin, gray] (trace) to[out=south west,in=south,looseness=1] node[pos=0.75,xshift=-10mm]{} (binder);
	\draw[->,>=stealth', thin, gray] (tracev) to[out=east,in=south east,looseness=1] node[pos=0.75,xshift=-10mm]{} (tracei);
	\draw[->,>=stealth', thin, gray] (tracei) to[out=south west,in=south,looseness=1] node[pos=0.75,xshift=-10mm]{} (binderi);
	\end{forest}	
\end{exe}
The system that moves heads around in German seems to be constrained by the Head Movement Constraint, in the same way that the system that moves heads around in English is.

Moreover, we can observe that there is a pattern to all these rules in German and English: the moved term always adjoins to another head. Because this appears to be something that is the same across German and English, it is a candidate language universal. Indeed, over a large range of cases examined across languages, this generalization appears to hold up. This suggests a constraint on movement rules that is sensitive to the phrase/head distinction. One of the first proposals for such a constraint is found in \citet{Baltin:1982}, who argues that this is part of a more general constraint on movement rules. He proposes that the position to which a phrase or head is moved always has the same phrasal status as the phrase or head being moved. So a head adjoins to another head position only, and a maximal projection can only move to another maximal projection position, and so on. Baltin calls this the \textsc{Like-attracts-Likes} Condition:
\begin{exe}
	\ex \label{ex4:28} \textsc{Likes Attracts Likes}\\ 
		An \obar{X} may only adjoin to another \obar{X}, and an XP may only move to a Specifier position or adjoin to another~XP.\\
		\hfill \citep[roughly (1), p.~2]{Baltin:1982}\hfill\hspace*{.15in}
\end{exe}
One can view the Likes Attracts Likes constraint as a special instance of constraint that Joe Emonds\cite{Emonds:1970} defended in his dissertation. Emonds suggested that no parse in a derivation can have a shape that couldn't have been made by phrase structure rules. %
%
\marginnote[20pt]{Emonds' proposal was slightly more complex. He suggested that there was a distinction between transformational rules that applied to root nodes only, and all the others. The root transformations had special properties. This distinction remains in contemporary theories---\obar{I}-to-\obar{C} movement is restricted to root clauses (mostly)---but there isn't a well-developed explanation of it.}%
%
In our terms, this amounts to saying that every parse in a derivation must meet the requirements of the \ibar{X} Skeleton. The \ibar{X} Skeleton, recall, sanctions the following shapes.
\begin{exe}
\raggedright
	\ex \label{ex3:skel1} \ibar{X} Skeleton
	\begin{xlist}
		\ex XP $\rightarrow$ \ibar{X}
		\ex XP $\rightarrow$ $\alpha$P \ibar{X}
		\ex \ibar{X} $\rightarrow$ \{ $\alpha$P, \ibar{X}\}
		\ex \ibar{X} $\rightarrow$ \{ \obar{X}, $\alpha$P\}
		\ex \ibar{X} $\rightarrow$ \obar{X}
	\end{xlist}
\end{exe}
For Emonds' condition to be correct, we'd have to add to (\ref{ex3:skel1}) the shapes in~(\ref{ex3:skel2}).
\begin{exe}
\raggedright
	\ex \label{ex3:skel2}
	\begin{xlist}
		\ex XP $\rightarrow$ \{ XP, $\alpha$P\}
		\ex \ibar{X} $\rightarrow$ \{ \obar{X}, \obar{$\alpha$}\}
	\end{xlist}
\end{exe}
Evidence that we need to add to our stock of licit shapes for a D-structure the one described by (\ref{ex3:skel2}a) comes, perhaps, from examples like~(\ref{ex3:advp}).
\begin{exe}
\raggedright
	\ex \label{ex3:advp}
	She said that probably Jill will leave.
	\sn
	\small
	\begin{forest}
	[\ibar{C}
		[\obar{C} [that]]
		[IP
			[AdvP [probably, roof] ]
			[IP
				[DP [Jill, roof] ]
				[\ibar{I} [will leave, roof]]
			 ]
		 ]
	]
	\end{forest}
\end{exe}
We'll see evidence for adding to our stock of possible shapes for a D-structure the form in (\ref{ex3:skel2}b) below (see page~\pageref{ex4:newX}). If we add (\ref{ex3:skel2}) to the \ibar{X} Skeleton, then Likes Attract Likes is a special instance of Emonds' condition. Emonds called his condition ``Structure Preservation,'' and I'll formulate it in terms appropriate for our grammar.
\begin{exe}
\raggedright
	\ex \label{ex3:strucpre} Structure Preservation\\
	Every element of a derivation must conform to the \ibar{X} Skeleton. 
\end{exe}
In what follows, I'll refer to this idea about how derivations are constrained variously as Likes Attracts Likes or Structure Preservation.

The derivations that involve moving \obar{V} or \obar{I} in both German and English conform to the Head Movement constraint and Structure Preservation. Finally, note that these movement rules both languages all have the feature that the verb or \obar{I} that has moved has moved up. We don't find cases where Verb Movement has relocated the verb downwards, as in examples like~(\ref{ex4:29}).
\begin{exe}
	\ex \label{ex4:29}
	\begin{xlist}
	\ex[*]{\gll \ldots da{\ss} Hans das Buch gehabt kauft.\\ \ldots that Hans the book had buys\\}
	\ex[*]{John must had buy the book.}
	\end{xlist}
\end{exe}
This is just a consequence of Upwards, which we've seen holds of all the movement rules we've seen so far. This too looks like a good candidate for a universal condition.

Let's take a look now at what we find with respect to verb placement in some of the other Germanic languages. In the Scandinavian languages we find a situation similar to German, as \citet{Vikner:1995} reviews. I will use Danish as a guide; much of what we see for Danish is found in Norwegian and Swedish as well. The same sensitivity to embedding is found in the placement of the finite verb in these languages. As the contrast in (\ref{ex4:32}) indicates, Danish is like German in allowing any constituent, but only one, to precede the finite verb in independent clauses. That is, it shares with German the trait of being ``Verb Second.''
\begin{exe}
	\ex \label{ex4:32} 
	\begin{xlist}
	\ex[]{\gll B{\o}rnen har set denne film \\ kids-the have seen this film\\ `The kids have seen this film.'}
	\ex[]{\gll Denne film bar b{\o}rnen set. \\ this film have kids-the seen\\ `The kids have seen this film.'}
	\ex[*]{\gll Denne film b{\o}rnen har set. \\ this film kids-the have seen\\ `The kids have seen this film.'}
	\end{xlist}
\end{exe}
But, as (\ref{ex4:33}) shows, these traits are not shared by embedded clauses, where, instead, the subject must precede the finite verb.
\pagebreak
\begin{exe}
	\ex \label{ex4:33} 
	\begin{xlist}
	\ex[*]{\gll Jeg ved ikke hvor {i g{\aa}r} har koen st{\aa}et.\\ I know not where yesterday has cow-the stood\\ `I don't know where the cow stood yesterday.' }
	\ex[]{\gll Jeg ved ikke hvor koen har {i g{\aa}r} st{\aa}et.\\ I know not where cow-the has yesterday stood\\ `I don't know where the cow stood yesterday.'}
	\end{xlist}
\end{exe}
Moreover, as the contrasts in (\ref{ex4:34}) show, the placement of the finite verb relative to negation is sensitive to the embedded/non-embedded context.
\begin{exe}
	\ex \label{ex4:34} 
	\begin{xlist}
	\ex[]{\gll B{\o}rnen har ikke set denne film.\\ kids-the have not seen this film\\ `The kids haven't seen this movie.'}
	\ex[*]{\gll B{\o}rnen ikke har set denne film \\ kids-the not have seen this film\\ `The kids haven't seen this movie.'}
	\ex[]{\gll Jeg ved hvor b{\o}rnen ikke har set denne film \\ I know where kids-the not have seen this film\\ `I know where the kids haven't seen this movie.'}
	\ex[*]{\gll Jeg ved hvor b{\o}rnen har ikke set denne film \\ I know where kids-the have not seen this film\\ `I know where the kids haven't seen this movie.'}
	\end{xlist}
\end{exe}
This indicates that the finite verb has moved out of VP, past negation, into \obar{C} in independent clauses, just as it does in German.

It looks, therefore, like Danish has the word-order of English---placing verbs before their objects---but the syntax of verb movement that German has. Note in particular that main verbs are able to move as well as auxiliary verbs, as we see partly in (\ref{ex4:33}) (look at the root clause word-order). If we were to look further, we would see that Danish (and the remainder of Germanic) also have the range of constraints we have seen on verb movement operations.

There are a couple differences, however. One at least potential difference between Danish and German is the existence of \obar{V}-to-\obar{I}. Whereas we cannot easily discern whether such an operation exists in German, it apparently does not in Danish (nor does it in standard Norwegian and Swedish), since when the verb has not relocated into \obar{C}, it remains to the right of negation. (Of course, we have made the perhaps incorrect assumption that negation occupies the same place in Danish as it does in English.)

Another difference concerns the range of contexts where verb second is found in embedded contexts. Recall that in German the verb moves into \obar{C} in embedded clauses only in those contexts where a complementizer is not required to fill that \obar{C}. But in Danish (as in the other Scandinavian languages), V2 word-order is possible even in those embedded clauses which have a complementizer associated with them, as in (\ref{ex4:35}).
\enlargethispage*{20pt}
\begin{exe}
	\ex \label{ex4:35} \gll Vi ved at denne bog har Bo ikke l{\ae}st. \\ we know that this book has Bo not read\\ `We know that Bo hasn't read this book.
\end{exe}
A variety of hypotheses about this situation have been offered. Let us adopt for concreteness the hypothesis that there is a CP ``shell'' that can be embedded within a regular CP. This CP shell provides the \obar{C} into which verbs move in embedded clauses in Scandinavian. This would give to the embedded CP in (\ref{ex4:35}), the representation in~(\ref{ex4:35t}).
\begin{exe}
	\raggedright
	\ex \label{ex4:35t} 
	\small
	\begin{forest}
	[CP
		[DP [denne bog, roof]]
		[\ibar{C} 
			[\obar{C}
				[\obar{I}, name=binder
					[\obar{V} [har] ]
					[\obar{I} [pres] ]
				 ]
				[\obar{C} [top] ]
			 ]
			[IP
				[DP [Bo, roof]]
				[\ibar{I}
					[{\ }, name=trace ]
					[NegP [\ibar{Neg}
						[\obar{Neg} [ikke] ]
						[VP [\ibar{V}
							[{\ }, name=tracev ]
							[vP [l{\ae}st, roof] ]
						] ]
					] ]
				 ]
			 ]
		]
	]
	\draw[->,>=stealth', ultra thin, gray] (tracev) to[out=-130,in=-109,looseness=1.5] node[pos=0.75,xshift=-10mm]{} (trace);
	\draw[->,>=stealth', ultra thin, gray] (trace) to[out=-115,in=south,looseness=1.5] node[pos=0.75,xshift=-10mm]{} (binder);
	\end{forest}
\end{exe}
\marginnote[-100pt]{In this derivation, the verb moves to \obar{I} across \obar{Neg} in violation of the Head Movement Constraint. We noted earlier that a similar situation arises in English. We will return to this feature of the derivation.}

So, let's see what we've got so far. If we factor out the parts to the set of rules in German, Danish and English that are common to the syntax of verb placement, we have the following:
\begin{exe}
	\ex \label{ex4:36} \textsc{Universals}\\
								a. Structure Preservation\\ 
								b. Upwards Constraint\\ 
								c. The Head Movement Constraint (with some counterexamples)\\ 
								e. Word Criterion
\end{exe}
The differences that are left can be boiled down to the following four parts.
\enlargethispage*{22pt}
\begin{exe}
	\ex \label{ex4:37} 
	\begin{xlist}
	\ex[]{German VPs (and IPs?) are head final.}
	\ex[]{English and Danish VPs (and IPs?) are head initial.}
	\ex[]{German and Danish: verbs (and \obar{I}?) are movable.}
	\ex[]{English: auxiliary verbs and \obar{I} are movable.}
	\ex[]{German and Danish: Root \obar{C} and Specifier of CP must be filled. (i.e., Verb Second)}
	\ex[]{English: Only wh-phrases move to Specifier of CP, and only questions trigger \obar{I}-to-\obar{C} movement.}
	\end{xlist}
\end{exe}
Putting this in the terms of language acquisition which we started with, the differences in (\ref{ex4:37}) constitute what must be learned in order to determine whether the grammar he or she is acquiring is German, Danish or English. By contrast, the Universals in (\ref{ex4:36}) are candidates for how the acquisition trajectory is limited. The guesses that children make about the syntax they are acquiring don't seem to include ones that violate one of the constraints in Universals.

\section{The Pollockian revolution}

Let us now add to our family of languages: French. In French we also find evidence of rules that move verbs. Just as in English, French verbs distribute themselves relative to polarity items, like sentential negation, on the basis of their inflectional class. Thus, for example, we find that finite verbs in French, just like their counterparts in English, only appear before negation.
\begin{exe}
	\ex \label{ex4:51.1}
	\begin{xlist}
	\ex[]{\gll Jean n'a pas lu livres.\\ John n'have not read books\\ `John hasn't read books.'}
	\ex[*]{Jean ne pas a lu livres.}
	\end{xlist}
\end{exe}
A word about negation in French is needed before the contrast in (\ref{ex4:51.1}) will be interpretable. Unlike English, sentential negation in French is expressed by way of two words, one (\emph{ne}) appears to be in a position like that we have put inflectional morphemes in English. We might imagine that like the contracted form of \emph{not} in English (\emph{n't}) it is contracted onto whatever word has moved into \obar{I} (as in (\ref{ex4:51.1}a)). The second word (\emph{pas}) is the one that has a distribution more like that of English not. And, as (\ref{ex4:51.1}) illustrates, the finite form of have (\emph{ha}) must precede this part of sentential negation.

French differs from English (but is like German and Danish) in moving main verbs. This can be seen in French by the fact that they are placed to the left of sentential negation when they are finite, as (\ref{ex4:51.2}) demonstrates.
\begin{exe}
	\ex \label{ex4:51.2}
	\begin{xlist}
	\ex[]{\gll Jean n'aime pas Marie. \\ John ne'love not Mary\\ `John doesn't love Mary.'}
	\ex[*]{Jean ne pas aime Marie.}
	\end{xlist}
\end{exe}
As the contrast between the position of main verbs in (\ref{ex4:51.2}a) and (\ref{ex4:51.1}a) indicates, it is the inflectional class of the verb that determines its position. That is, just as in English, there is a correlation between the inflectional class of the verb and its syntactic position---a correlation that is captured by fixing the position of inflectional morphemes with the phrase structure rules and driving the verbs to these positions with the verb movement operation.

That main verbs may move is also responsible for another difference between French and English. This derives from the ordering that the Projection Principle places on arguments and non-arguments. As we have had occasion to discuss several times, now, the Projection Principle forces arguments to be closer to their \mbox{$\theta$-marking} head than non-arguments. In verb initial languages, like French and English, this means that complements should come before non-complements. This is roughly true for English, as we have seen; but it is not the case for French.
\begin{exe}
	\ex \label{ex4:51.3}
	\begin{xlist}
	\ex[]{Jean embrasse souvent Marie.}
	\ex[*]{John kisses often Mary.}
	\end{xlist}
\end{exe}
This apparent difference in the function of the Projection Principle can be credited to verb movement. Since we already know that main verbs in French but not English move into \obar{I}, it makes sense that the main verbs in French, but not English, should be able to be separated from their objects by all sorts of material, including adverbs.

But now consider what happens in French non-finite clauses.%
%
\marginnote[30pt]{These data all from \citet{Pollock:1989}.}
%
\begin{exe}
	\ex \label{ex4:51.4}
	\begin{xlist}
	\ex[]{\gll \textbf{Comprendre} \textit{{\`a} peine} l'italien apres cinq ans d'{\'e}tude d{\'e}note un manque de don pour les langues. \\ understand barely the-Italian after five years {of study} shows a lack of gift for the languages\\ `To barely understand Italian after five years of study shows a lack of talent for languages.'}
	\ex[]{\gll \textbf{Perdre} \textit{compl{\`e}tement} la t{\^e}te pour les belles \'etudiantes c'est dangereux. \\ lose completely the head for the pretty students {it is} dangerous\\ `To completely lose your head for pretty students is dangerous.'}
	\end{xlist}
	\ex \label{ex4:51.4a}
	\begin{xlist}
	\ex[]{\gll ne \textit{pas} \textbf{sembler} heureux est une condition pour \'ecrire {des romans}\\ ne not seem happy is a condition for writing novels\\ `To not seem happy is a (pre?)condition for writing novels.'}
	\ex[*]{\gll ne \textbf{sembler} \textit{pas} heureux est une condition pour \'ecrire {des romans}. \\ ne seem not happy is a condition for writing novels\\ `To not seem happy is a (pre?)condition for writing.'}
	\end{xlist}
\end{exe}
Here it appears that main verbs may move past adverbs but not negation. (The infinitival verbs are in bold face, in these examples, and the term that separates them from their objects is in italics.) In (\ref{ex4:51.4}), we see examples of main verbs with an adverb between it and its object, and this word order is grammatical. In (\ref{ex4:51.4a}), by contrast, we see that the main verb \emph{sembler} (`seem') cannot be separated from its object (here the adjective \emph{heureux} (`happy')) by negation.
\pagebreak

Auxiliary verbs behave differently.
\begin{exe}
	\ex \label{ex4:frenchauxinf}
	\begin{xlist}
	\ex[]{\gll ne \textit{pas} \textbf{{\^e}tre} heureux est une condition pour {\'e}crire {des romans}\\ ne not be happy is a condition for writing novels\\ `To not be happy is a condition for writing novels.'}
	\ex[]{\gll N'\textbf{\^etre} \textit{pas} heureux est une condition pour \'ecrire {des romans}.\\ ne'be not happy is a condition for writing novels\\ `To not be happy is a condition for writing novels.'}
	\end{xlist}
\end{exe}
They may optionally appear to the left of negation.

We learn two important things from these facts. First, that the contrast between auxiliary and main verbs that seems to distinguish English from German/Danish is actually not a distinction in languages. Instead, it is a difference which is found within a single language: French. That is, this is not a parameter along which languages vary, then, but rather a parameter along which clause-types vary. We need to express the distinction between English, and these other languages, in terms of the clause types these languages host, not in terms of the targets for their verb movement rules. So, we're looking for something that distinguishes English finite clauses and French non-finite clauses on the one hand from French finite clauses on the other.

The second lesson of this paradigm is the main point of Pollock's paper. He suggests that the two different positions that main verbs may occupy across finite and non-finite clauses warrants giving a structure like that in (\ref{ex4:51.7}) to clauses, where the higher \obar{I} is equated with Tense and the lower one with Agreement. This answers to the correlation that appears to hold for main verbs between whether they bear tense morphology or not: in finite clauses they do, and in non-finite clauses they do not.
\begin{exe}
\ex \label{ex4:51.7}
\small
\begin{forest}
[TP[\ibar{T}
	[\obar{T} ]
	[NegP [\ibar{Neg}
		[\obar{Neg} ]
		[AgrP [\ibar{Agr}
			[\obar{Agr} ]
			[VP ]
		] ]
	] ]
]]
\end{forest}
\end{exe}
Pollock also introduces the idea that Neg heads a phrase which stands between the two other Inflectional phrases, as indicated in this pares. There are several reasons for wanting to treat negation differently from other adverbials. One is that the syntactic position of negation is comparatively rigid when compared to other adverbs. Another is that only negation blocks the process that inflects main verbs in English, as we saw earlier; other adverbs don't. Pollock suggests distinguishing negation from adverbs structurally, and then making reference to this structural difference to capture these ways in which negation behaves uniquely. I won't examine in detail how this idea works, partly because we will eventually go in a slightly different direction than does Pollock. But let us nonetheless adopt---at least as an intermediary hypothesis---the thesis that negation does head a phrase as shown in (\ref{ex4:51.7}). It should be noted, however, that this introduces a problem: how is it that \obar{Agr} can move to \obar{T} past negation without violating the Head Movement Constraint. Pollock offers a solution to this problem that I will come to soon.

Note that though the difference between main and auxiliary verbs in French non-finite clauses that we've just reviewed speaks on behalf of two head positions to which verbs may move, it doesn't really indicate what the identity of these two head positions might be. While most of Pollock's proposal involves examining how the hypothesis that there are two inflectional positions within a sentence can be used to explain the differences in verb position across clause types, he also assigns values to these inflectional phrases. But, in fact, it is extremely difficult to tell what the value of these heads is. Pollock decides in favor of giving the higher \obar{I} the value of Tense and the lower one the value of Agreement. This involves two decisions: first that these positions should be equated with inflectional categories, and second that they should be equated with these particular ones. The first decision, note, is a bit of a leap here as there is no convergent evidence of the headedness sort to anchor this conclusion. That is, part of the reason we decided that the landing site of verb is \obar{I} came in response to the evidence that \obar{I} was a head.

There is, however, a sort of reasoning which suggests that the first decision is correct, but that the second isn't. Consider the following verbs, drawn from a variety of IndoEuropean languages.%
%
\marginnote[0pt]{The argument I'm reporting here comes from \citet{Belletti:1990}.}%
%
\begin{exe}
	\ex \label{ex4:51.8}
	\begin{xlist}
	\ex[]{\gll Legg - eva - no\\ read - Imperfect - 3plur\\}
	\ex[]{\gll Parl - er - \`o \\ speak - Future - 1sing\\}
	\ex[]{\gll Buhl - t - en \\ woo - past - plur\\}
	\end{xlist}
\end{exe}
In each of these cases, the morpheme which encodes tense information precedes that which encodes agreement. Is there a way of determine the syntactic arrangement of inflectional terms from their morphological arrangement? Mark Baker has made famous an argument for thinking that the answer to this question is yes. Following proposals by Pieter Muysken and Donna Gerdts, Baker argues that there are correlations between the order that inflectional affixes have relative to the verbal stem and the order of the syntactic operations that these affixes encode. Thus, for example, in Chamorro, the passive morpheme necessarily comes closer to the verbal stem than does the subject agreement affix, as (\ref{ex4:51.9}) shows.
% //todo Put in typological fact from Marit about aspect and tense.
\begin{exe}
	\ex \label{ex4:51.9} \gll Para.u.fan - s - in - aolak l famgu'un {gi as} tat-n-niha\\ irr.3pS - pl - pass - spank the children obl father.their\\ `The children are going to be spanked by their father.'%
	%
	\marginnote[20pt]{(\ref{ex4:51.9}) comes from \citet{Gibson:1980}. The principle formulated in \citet{Baker:1985a} has its roots in many earlier works, among which are: \citet{Muysken:1979, Muysken:1981} and \citet{Gerdts:1981}.}%
	%
\end{exe}
This ordering on affixation mimics the order in which the syntactic operations of Passive and subject agreement take place---the passive must bring the underlying object into the subject relation before the agreement process can apply.
% \pagebreak

Baker calls correlations of this type, which he suggests are widespread enough to be considered law-like, the ``Mirror Principle.''
\begin{exe}
	\ex \label{ex4:51.10} \textsc{Mirror Principle}\\
				Morphological derivations must directly reflect syntactic derivations (and vice versa).\\
\hspace*{\fill} \citep[p.~375]{Baker:1985a} \hspace*{.5in}
\end{exe}
Baker argues that the Mirror Principle will emerge once the correct theory of the syntax-morphology interface is found. In particular, (\ref{ex4:51.10}) can be derived on models that posit for each of the agreement, passive, and other such operations one rule underlying both the morphological and syntactic effects. Baker tentatively concludes that these meta-rules apply in the Syntax---that is, after the verbal roots are inserted into phrase-markers. The resulting picture is sketched in (\ref{ex4:51.11}).
\begin{exe}
	\ex \label{ex4:51.11}
	\[ \begin{array}{lcr}
		\mbox{D-structure:} & V_{\text{root}} & \\
		& \downarrow & \leftarrow \: \text{Metarule} \; \alpha\\
		& V_{\text{root}}+Af_{\alpha} & \\
		& \downarrow & \leftarrow \: \text{Metarule} \; \beta\\
		\mbox{S-structure:} & V_{\text{root}}+Af_{\alpha}+Af_{\beta} &
		\end{array}
	\]
\end{exe}
Each metarule brings about the relevant syntactic operation and adds to the verbal root the corresponding affix. Note in particular that this picture correlates syntactic operations with affixes; the correspondences that Baker summarizes with the Mirror Principle only concern the relation between affixal orderings and the orderings of syntactic operations. Indeed, Baker's argument leads to the conclusion that the metarules in (\ref{ex4:51.11}) necessarily involve affix-syntactic operation pairs. It is the order of affixes that correlates with the relevant syntactic operations, and not some more abstract information, such as morphological features or the like. We shall have an opportunity to revisit this point.

Now if Baker's conclusions from the Mirror Principle are imported into the domain of Verb Movement, then the arrangement of morphemes in (\ref{ex4:51.8}) suggests that the higher of Pollock's positions should be associated with agreement morphology, and the lower with tense morphology. Then, the fact that tense morphology comes closer to the verbal stem will follow from the Head Movement Constraint. So we should adopt, perhaps, a picture like that in (\ref{ex4:51.12}).
\begin{exe}
\ex \label{ex4:51.12}
\small
\begin{forest}
[AgrP[\ibar{Agr}
	[\obar{Agr} ]
	[NegP [\ibar{Neg}
		[\obar{Neg} ]
		[TP [\ibar{T}
			[\obar{T} ]
			[VP ]
		] ]
	] ]
]]
\end{forest}
\end{exe}

We can reach this conclusion, perhaps, through a different route. If we compare Icelandic with Danish, we see a difference in verb placement that suggests that only Icelandic has movement of verbs into the \obar{I} that stands above negation. In embedded Icelandic clauses the finite verb must precede negation, as in~(\ref{ex4:51.13}).
\begin{exe}
	\ex \label{ex4:51.13}
	\begin{xlist}
	\ex[*]{\gll {\'Eg} spur{\dh}i {af hverju} Helgi ekki hef{\dh}i {lesi\dh\ } \textipa{\textthorn}essa b\'ok\\ I asked whether Helgi not had read this book\\ `I asked whether Helgi hadn't read this book.'}
	\ex[]{\gll \'Eg spur{\dh}i {af hverju} Helgi hef{\dh}i ekki {lesi{\dh}\ } \textipa{\textthorn}essa b\'ok \\ I asked whether Helgi had not read this book\\ `I asked whether Helgi hadn't read the book.'}
	\end{xlist}
\end{exe}
But as we have seen, Danish finite verbs cannot precede negation. Another example illustrating this fact is~(\ref{ex4:51.14}).
\begin{exe}
	\ex \label{ex4:51.14}
	\begin{xlist}
	\ex[]{\gll Jeg spurgte hvorfor Peter ikke havde l{\ae}st den.\\ I asked why Peter not had read it\\ `I asked why Peter hadn't read it.'}
	\ex[*]{\gll Jeg spurgte hvorfor Peter havde ikke l{\ae}st den.\\ I asked why Peter had not read it\\  `I asked why Peter hadn't read it.'}
	\end{xlist}
\end{exe}
If we assume that the position of negation is constant across these languages, then this indicates that the verb has moved past this position in Icelandic, but not Danish. Now, interestingly, Icelandic has a full paradigm of subject agreement, but Danish has no subject agreement. Comparing Icelandic with Danish, then, what we find is a correlation between subject agreement and position to the left of \obar{Neg}. Indeed, this correlation holds throughout the Scandinavian languages.%
%
\marginnote[-20pt]{This correlation holds only partially. There are some cases in which agreement morphology exists, though in a very impoverished form, and verb movement past negation nonetheless appears to be blocked. (This happens in certain dialects of Faroese, in a register of Norwegian, and in a Swedish dialect heavily influenced by Finnish.) Some of these facts are reported in \citet{Vikner:1995}, where references to the relevant literature can also be found.}% 
%
Standard Swedish and Norwegian are like Danish in lacking subject agreement morphology, and they are like Danish as well in placing their finite verbs to the right of negation. This is just the correlation that (\ref{ex4:51.12}) predicts.

Adopting this modified Pollockian parse for clauses, we are left with a series of questions about how verbs distribute themselves over the two head positions. Why can main verbs not move to the ``agreement'' \obar{X} in non-finite clauses in French, for instance; and why are they unable to occupy this position even in finite clauses in English? We don't have satisfactory answers to these questions. Pollock relies on an observation due originally to Ian Roberts: The \obar{X}s in which main verbs cannot surface are devoid or close to devoid of morphological content. English agreement is much less ``robust,'' somehow, than is French agreement. And subject agreement in French non-finite clauses is considerably less pronounced than it is in French finite clauses. (In fact, it is thoroughly absent.) Now, this difference in agreement morphology Pollock exploits to control whether or not a verb can move into the lower of the two inflectional \obar{X}s---this is because Pollock assigned agreement to this lower head. But we can copy his idea into the assignment we have given to these positions. Tense morphology is also richer in French than it is in English, in both finite and non-finite clauses; and agreement is richer in finite clauses than it is in non-finite clauses too. In fact, there is no evidence of agreement at all in infinitival clauses.

Even if this ``robustness'' in the morphology associated with the \obar{T} and \obar{Agr} positions is not what is relevant (and further cases suggests that it isn't), let's continue to exploit the language that Pollock adopts for whatever the relevant difference turns out to be. Those heads that are associated with morphology that allows main verbs to occupy them, let's call ``strong.'' And those head positions that are not associated with morphology that allows main verbs to sit in them, let's call ``weak.'' We can now describe in these terms what we've seen so far. In particular, we will ascribe to the Tense and Agreement positions different ``strengths,'' and in this way describe how they control the Move \obar{X} operation. For English and French, we might assign the values as shown in (\ref{ex4:Eng}) and~(\ref{ex4:French}).
\begin{exe}
	\ex \label{ex4:Eng} \textbf{English:}
	\begin{xlist}
	\ex[]{\emph{Finite:}\\ Tense: weak\\ Agreement: weak}
	\ex[]{\emph{Infinitive:}\\ Tense: \emph{to}\\ Agreement: empty?}
	\end{xlist}
	\ex \label{ex4:French} \textbf{French:}
	\begin{xlist}
	\ex[]{\emph{Finite:}\\ Tense: can't tell\\ Agreement: strong}
	\ex[]{\emph{Infinitive:}\\ Tense: strong\\ Agreement: empty}
	\end{xlist}
\end{exe}

In these value assignments, I've decided to treat the infinitival marker in English, i.e., \emph{to}, as something that can optionally move from \obar{T} to \obar{Agr}. On this view it is rather like an auxiliary verb in French, and should behave similarly, therefore, in infinitival clauses. For some speakers of English there is, in fact, an alternation that resembles that we've seen for French auxiliary verbs:
\begin{exe}
	\ex \label{ex4:51.17}
	\begin{xlist}
	\ex[]{Jerry tried [ not to eat candy ].}
	\ex[]{Jerry tried [ to not eat candy ].}
	\end{xlist}
\end{exe}
Note that we cannot determine whether finite \obar{T} is strong or weak in French because verbs are always driven up to \obar{Agr} in this context.

Turning now to Scandinavian, we have something like:%\enlargethispage*{22pt}
\begin{exe}
	\ex \label{ex4:51.18} \textbf{Danish, Norwegian and Swedish:}
	\begin{xlist}
	\ex[]{\emph{Finite:}\\ Tense: weak\\ Agreement: not there}
	\ex[]{\emph{Infinitive:}\\ haven't looked yet (it varies across the group, as it turns out)}
	\end{xlist}
	\ex \label{ex4:ice} \textbf{Icelandic:}
	\begin{xlist}
	\ex[]{\emph{Finite:}\\ Tense: can't tell\\ Agreement: strong}
	\ex[]{\emph{Infinitive:}\\ haven't looked yet}
	\end{xlist}
\end{exe}
My decision to classify the Mainland Scandinavian group as having weak tense morphology is based on the fact that this language group behaves like English, and not like French, with respect to the positioning of main verbs relative to their complements. That is, in this group, non-complements cannot intervene between main verb and DP objects (except in contexts where the verb has relocated into \obar{C}).

Okay, this is one part of the story. The morals are: (i) there are two inflectional heads between \obar{C} and VP in some languages, and (ii) it is properties of those heads that determine what must and can move into them. Let's now address the problem that treating negation as a head poses. Recall that we should not expect the Head Movement Constraint to allow movement of \obar{T} to \obar{Agr} past negation, if negation is a head that lies between them. Pollock's suggestion is that \emph{not} and \emph{pas} are in the Specifier of NegP, and that \emph{ne} is what heads NegP in French. In English, we might imagine that this head is silent; I'll represent this silent morpheme with \emph{neg}. Thus, the correct way of parsing sentences with sentential negation in them is as in~(\ref{ex4:51.20}).
\begin{exe}
\ex \label{ex4:51.20}
\small
\begin{forest}
[XP
	[XP ]
	[\ibar{Agr}
		[\obar{Agr} ]
		[NegP 
			[AdvP [pas\\ not, roof]]
			[\ibar{Neg}
				[\obar{Neg} [ne] ]
				[TP [\ibar{T}
					[\obar{T} ]
					[VP ]
				] ]
			 ]
		 ]
	 ]
]
\end{forest}
\end{exe}
On this view, then, when \obar{T} moves to \obar{Agr} it doesn't have to skip \obar{Neg}, it can move through this position and pick up the head (if there is one) on the way. In a finite French clause, then, the surface parse might look roughly as in (\ref{ex4:nehmc}) \vpageref[below]{ex4:nehmc}.
\begin{exe}
\ex \label{ex4:nehmc}
\small
\begin{forest}
[AgrP
	[DP [Jean, roof]]
	[\ibar{Agr} 
		[\obar{Agr}
			[\obar{Neg} , name=binder
				[\obar{T}
					[\obar{V} [etre] ]
					[\ibar{T} [\emph{inf}] ]
				 ]
				[\obar{Neg} [ne] ]
			 ]
			 [\obar{Agr}]
		]
			[NegP
				[AdvP [pas, roof]]
				[\ibar{Neg}
					[{\ }, name=traceneg ]
					[TP [\ibar{T}
						[{\ }, name=tracetense ]
						[VP [\ibar{V}
							[{\ }, name=traceverb ]
							[AP [heureux, roof] ]
						] ]
					] ]
				 ]
			 ]
		]
	]
]
\draw[->,>=stealth', ultra thin, gray] (traceverb) to[out=south west,in=south,looseness=1] node[pos=0.75,xshift=-10mm]{} (tracetense);
\draw[->,>=stealth', ultra thin, gray] (tracetense) to[out=south west,in=south,looseness=1] node[pos=0.75,xshift=-10mm]{} (traceneg);
\draw[->,>=stealth', ultra thin, gray] (traceneg) to[out=south west,in=south,looseness=2] node[pos=0.75,xshift=-10mm]{} (binder);
\end{forest}
\end{exe}
As this parse indicates, \emph{ne} gets carried to the position before \emph{pas} by the tensed verb as it moves into \obar{Agr}.%
%
\marginnote[-2in]{The {S-structure} in (\ref{ex4:nehmc}) is derived as follows (ignoring the movement of the subject). The main verb, \emph{aime}, is moved into \obar{T}, where it combines with the tense morpheme. The resulting tensed verb (that is, \obar{T}) moves and adjoins to \emph{ne}. the resulting negated verb (that is, \obar{Neg}) moves into \obar{Agr}, where the agreement morpheme combines with the verb. This surface parse, then, claims that negation---\emph{ne}---is closer to the verbal stem than is the agreement morpheme, and this is somewhat surprising.}%
%
\ Note that we have to let something force the bound morphemes, that is, the tense and agreement suffixes, to linearize to the right of the term that has adjoined to them, but force \emph{ne} to linearize to the left of the term that has adjoined to it.

This can't be the only way that \emph{ne} manages to get to the left of \emph{pas}, because \emph{ne} shows up before \emph{pas} in infinitival clauses as well, even when no verb has moved into \obar{Agr}. In these situations, we might imagine that \emph{ne} moves on its own into \obar{Agr}. This, in fact, is what Pollock suggests. He hypothesizes that \emph{ne} adjoins to \obar{Agr} on its own in much the same way that so-called ``clitics'' do. ``Clitics'' are prosodically weak pronouns that, in the Romance languages, surface in construction with a verb; more particularly, they form a prosodic word with a verb and are therefore adjacent to that verb (or adjacent to something else that is part of the verb). In finite clauses, it is the finite verb that clitics appear with, and in infinitival clauses, it is typically the infinitival verb that the clitic appears on. One way of capturing this distribution is let clitics adjoin to \obar{Agr}, where they will come close enough to the finite or infinitival verb to be parsed with that verb prosodically.

Clitics, then, instantiate another instance of movement. Object clitics, for instance, satisfy the Theta Criterion and Projection Principle by way of occupying a position that is sister to the verb they are arguments of, but surface in \obar{Agr} position. This type of movement is called ``Cliticization.'' Cliticization doesn't have the same properties that Head Movement has---it isn't subject to the Head Movement Constraint, for instance---and therefore we can treat it is as being independent of the rules we are investigating here. With this assumption about \emph{ne}, then, an infinitive that doesn't have verb movement to \obar{Agr}, but is negated, might have an S-structure like~(\ref{ex4:51.22}).
\begin{exe}
\ex \label{ex4:51.22}
\small
\begin{forest}
[AgrP
	[PRO ]
	[\ibar{Agr}
		[\obar{Agr}
			[ne, name=binder ]
			[\obar{Agr}  ]
		 ]
		[NegP
			[AdvP [pas, roof]]
			[\ibar{Neg}
				[{\ }, name=traceneg ]
				[TP [\ibar{T}
					[\obar{T}
						[\obar{V} [sembl, name=binderv] ]
						[\obar{T} [er] ]
					 ]
					[VP [\ibar{V}
						[{\ }, name=traceverb ]
						[AP [heureux]]
					] ]
				] ]
			 ]
		 ]
	 ]
]
\draw[->,>=stealth', ultra thin, gray] (traceneg) to[out=south west,in=south,looseness=1] node[pos=0.75,xshift=-10mm]{} (binder);
\draw[->,>=stealth', ultra thin, gray] (traceverb) to[out=south west,in=-80,looseness=1] node[pos=0.75,xshift=-10mm]{} (binderv);
\end{forest}
\vspace*{15pt}
\end{exe}

If we adopt Pollock's suggestion, that \emph{ne} cliticizes to \obar{Agr}, we might be better poised to understand why it is linearized to the left of the verb+ tense+\-agreement complex. Indeed, we could impose a general constraint on adjunction that would linearize the moved item to the left of the phrase it adjoins to. Something along the lines of (\ref{ex4:adjlin}), perhaps.
\begin{exe}
	\ex \label{ex4:adjlin} \textsc{Adjunction Linearization}\\
				If \obar{$\alpha$} adjoins to \obar{$\beta$}, then \obar{$\alpha$} precedes \obar{$\beta$}.
\end{exe}
That tense morphemes are linearized to the right of their verbal stem is the predominant order in the world's languages.%
%
\marginnote{Julien equated tense and aspect morphemes in her survey. These are often difficult to distinguish, and in some languages aspect is the only means of expressing temporal information.}%
%
\ \citet{Julien:2002b}, who surveyed 530 languages, found tense suffixes far outnumbered tense prefixes. She divided her sample of languages into 280 genera—groups of languages that are historically related, and counted how many of those genera used tense prefixes or tense suffixes. (If languages in a genera had either order, the genera was counted as illustrating both types.) What she found was that 202 genera had tense suffixes but only 60 genera had tense prefixes. Most of the languages that had tense prefixes in her survey are head initial, and Julien argues that they have a different morphological status than tense suffixes.

With these preliminaries, let's reconsider how the derivation in (\ref{ex4:nehmc}) might go for French. The D-structure representation would be as in~(\ref{ex4:nehmcd}).
\begin{exe}
	\raggedright
	\ex \label{ex4:nehmcd}
	\small
	\begin{forest}
	[AgrP
		[DP [Jean, roof]]
		[\ibar{Agr} 
			[\obar{Agr} [\nth{3}\\ sing] ]
			[NegP
				[AdvP  [pas]]
				[\ibar{Neg}
					[\obar{Neg} [ne] ]
					[TP [\ibar{T}
						[\obar{T} [pres] ]
						[VP [\ibar{V}
							[\obar{V} [\^etre] ]
							[AP [heureux, roof] ]
						] ]
					] ]
				 ]
			 ]
		]
	]
	\end{forest}
\end{exe}
The first two steps of the derivation could involve Verb Movement to \obar{T} and cliticization of \emph{ne}, giving~(\ref{ex4:nehmc1}).
\begin{exe}
	\raggedright
	\ex \label{ex4:nehmc1}
	\small
	\begin{forest}
	[AgrP
		[DP [Jean, roof]]
		[\ibar{Agr} 
			[\obar{Agr}
				[\obar{Neg} [ne, name=binder] ]
				[\obar{Agr} [\nth{3}\\ sing] ]
			 ]
			[ NegP
				[AdvP  [pas]]
				[ \ibar{Neg}
					[{\ }, name=trace ]
					[TP [\ibar{T}
						[\obar{T}
							[\obar{V} [\^etre, name=binderV] ]
							[\obar{T} [pres] ]
						 ]
						[VP [\ibar{V}
							[{\ }, name=traceV ]
							[AP [heureux, roof] ]
						] ]
					] ]
				]
			]
		]
	]
	\draw[->,>=stealth', ultra thin, gray] (traceV) to[out=south west,in=-80,looseness=1] node[pos=0.75,xshift=-10mm]{} (binderV);
	\draw[->,>=stealth', ultra thin, gray] (trace) to[out=south west,in=south,looseness=1] node[pos=0.75,xshift=-10mm]{} (binder);
	\end{forest}
\end{exe}
At this step in the derivation, \obar{Neg} no longer stands between \obar{Agr} and \obar{T}, and this means that using movement to bring them together no longer invokes a violation of the Head Movement Constraint. To bring \obar{T} together with the \obar{Agr} that contains the agreement morphology, but not \emph{ne}, could be achieved by using Head Movement to produce~(\ref{ex4:nehmc2}).
\begin{exe}
	\raggedright
	\ex \label{ex4:nehmc2}
	\small
	\begin{forest}
	[AgrP
		[DP [Jean, roof]]
		[\ibar{Agr} 
			[\obar{Agr}
				[\obar{Neg} [ne, name=binderne] ]
				[\obar{Agr}
					[\obar{T}, name=binder
						[\obar{V} [\^etre] ]
						[\obar{T} [pres] ]
					 ]
					[\obar{Agr} [\nth{3}\\ sing] ]
				 ]
			 ]
			[NegP
				[AdvP [pas, roof]]
				[\ibar{Neg}
					[{\ } ]
					[TP [\ibar{T}
						[{\ }, name=trace ]
						[VP [\ibar{V}
							[{\ } ]
							[AP [heureux, roof] ]
						] ]
					] ]
				 ]
			 ]
		]
	]
	\draw[->,>=stealth', ultra thin, gray] (trace) to[out=south west,in=south,looseness=1] node[pos=0.75,xshift=-10mm]{} (binder);
	\end{forest}
\end{exe}
Let us assume that these are the derivations that French uses to avoid violation of the Head Movement Constraint.%
%
\marginnote[-180pt]{Note that Upwards permits this movement. The sister to the moved \obar{T} in the position it moves to is \obar{Agr}. And \obar{Agr} has a projection, namely \ibar{Agr}, that dominates the position that \obar{T} moved from.} We will return to look at how English manages this.%
%

So, now what we would like to do is see if we can't get a better understanding of what it is about this ``strong''/``weak'' contrast that makes main verbs susceptible to it, but not auxiliary verbs. Pollock makes a proposal in this direction. He suggests that we think of main verbs as differing from auxiliary verbs with respect to their ability to assign a \mbox{$\theta$-role}. In particular, he suggests that we define main verbs as ones that have a \mbox{$\theta$-role} to assign, and auxiliary verbs as ones that don't. He then crafts a proposal that would prevent verbs from assigning a \mbox{$\theta$-role} if they move to a position associated with a ``weak'' inflectional head. And this will essentially prevent main verbs from both moving to and moving through the positions we have assigned the ``weak'' value.

There are some immediate problems with this idea that Pollock himself addresses. It wouldn't seem able to account for the British (\ref{ex4:51.24}), nor for the contrast in (\ref{ex4:51.25}).
\begin{exe}
	\ex \label{ex4:51.25} Mary hasn't a dime.
	\ex \label{ex4:51.24}
	\begin{xlist}
	\ex[]{Mary isn't unhappy.}
	\ex[*]{Mary seemsn't unhappy.}
	\end{xlist}
\end{exe}
Why doesn't the \emph{be} in (\ref{ex4:51.25}a) have the same number of \mbox{$\theta$-role}s as (\ref{ex4:51.25}b)? And it sure looks like \emph{has} in (\ref{ex4:51.24}) is capable of moving past negation, though it appears equally like it assigns \mbox{$\theta$-role}. Pollock suggests that the \mbox{$\theta$-role} assigned in (\ref{ex4:51.24}) comes not from the verb, but from some hidden term. Something similar might be attempted in distinguishing the cases in (\ref{ex4:51.25}).

But there is another kind of problem, specifically for the prediction that movement through a weak \obar{X} is unavailable to main verbs,  which seems to me insurmountable. Danish, as we have seen, and the other Mainland Scandinavian languages, show Verb Second word-order; (\ref{ex4:51.26}) is an example, for instance, where the main verb has relocated to \obar{C} position.
\begin{exe}
	\ex \label{ex4:51.26} \gll K\o bte han bogen\\ bought he book-the\\ `Did he buy the book?'
\end{exe}
But we have also just discovered that Danish has weak tense morphology. Thus, if we preserve the Head Movement Constraint, this entails that Danish main verbs are capable of moving into \obar{T} and from there moving into \obar{C}. If Danish \obar{T} is weak, this means, in turn, that moving through a weak \obar{X} is possible for main verbs.

For this reason, I will interpret the weak/strong typology to determine only whether a main verb may surface in the relevant \obar{X}. This seems to be the way these values are interpreted most commonly in the literature. It is, in fact, on this basis that I gave the assignments of ``weak'' and ``strong'' to \obar{Agr} and \obar{T} in English, French and the Scandinavian languages in (\ref{ex4:Eng}), (\ref{ex4:French}), (\ref{ex4:51.18}) and (\ref{ex4:ice}). In these assignments, then, I have given a ``weak'' value to those \obar{X} positions for which there is evidence that main verbs cannot surface, and to those \obar{X} positions where there is evidence that main verbs can surface I have given the ``strong'' value. I have also indicated certain positions---namely, French and English infinitival \obar{Agr}---as empty positions. This is my, somewhat idiosyncratic, way of expressing the optionality of movement to these spots. Because they have no morpheme in them, these infinitival \obar{Agr}s do not cause the Word Criterion force a verb stem to move into this position. But, because it is present, it is still a possible landing site for \obar{X} movement. It's in this respect, then, that the infinitival \obar{Agr}s of these languages are different from the finite \obar{Agr} of, say, Danish. Movement of a verb to the position we've identified for finite \obar{Agr} is not possible in Danish---neither auxiliary nor main verbs may surface in this position---and I've coded that here by making \obar{Agr} completely absent in Danish.

Under this typology, \obar{X} positions identified with inflectional morphemes come in four strengths. In their weakest manifestation, they are completely absent, and nothing, as a consequence, may surface in their position. In their next strongest manifestation, they are present but are associated with no morpheme. In this guise we find optional movement of the infinitival marker \emph{to} and auxiliary verbs, but main verbs are still not capable of surfacing in them. (This is what we see for the infinitival \obar{Agr} in French.) In their next strongest manifestation, they are associated with a ``weak'' morpheme. At this level of strength, the Word Criterion kicks in and requires that something express the morphology they are associated with. The ``weak'' property, however, prevents main verbs from doing this by adjoining to the morpheme. Finally, there are the \obar{X} positions occupied by a ``strong'' morpheme. These positions allow either auxiliary or main verb to surface in them, and the Word Criterion forces one or the other to. Because we want to prevent main verbs from surfacing in all but the last of these positions, the constraint that controls the auxiliary/main verb distinction should be stated as~(\ref{ex4:51.27}).%\enlargethispage*{11pt}
\begin{exe}
	\ex \label{ex4:51.27} A main verb merged to a morpheme, $\alpha$, at {S-structure} creates ungrammaticality unless $\alpha$ is ``strong.''%
	%
	\marginnote[-25pt]{(\ref{ex4:51.27}) is going to have certain consequences for the syntax of a sentence in which the verb has moved into \obar{C}. In these cases---questions and, in the Germanic languages other than English, root finite clauses---we must imagine that the \emph{top} morpheme that occupies~\obar{C} is strong.}%
	%
\end{exe}

\section{\textsc{agree}}

We've developed a rough description of  how main and auxiliary verbs are driven into \obar{I}s that uses the ``strong''/``weak'' vocabulary. We should now look at how this matter interacts with the way in which inflectional heads assign their morphology to the relevant verbs when verb movement is not involved. The conception we have adopted of how verbs inflect is that it parallels the kind of relation that we have seen in connection with Case. The tense and agreement heads assign the relevant inflection to the verb in configurations like~(\ref{ex4:agree1a}).
\begin{exe}
\ex \label{ex4:agree1a}
\small
\begin{forest}
[AgrP
	[DP [they, roof]]
	[\ibar{Agr} 
		[\obar{Agr} [\nth{3}\\ plur] ]
		[TP [\ibar{T}
			[\obar{T} [past] ]
			[VP [\ibar{V}
				[\obar{V} [eat] ]
				[DP [nattoo, roof] ]
			] ]
		] ]
	]
]
\end{forest}
\end{exe}
When auxiliary verbs are involved, the verb assigned inflection moves to the heads responsible for that inflection, as in~(\ref{ex4:agree2}).
\begin{exe}
\ex \label{ex4:agree2}
\small
\begin{forest}
[AgrP
	[DP [they, roof]]
	[\ibar{Agr} 
		[\obar{Agr}
			[\obar{T} , name=binderT
				[\obar{V} [have] ]
				[\obar{T} [pres] ]
			]
			[\obar{Agr} [\nth{3}\\ plur] ]
		 ]
		[TP [\ibar{T}
			[{\ }, name=traceT ]
			[VP [\ibar{V}
				[{\ }, name=traceV ]
				[IP [eaten nattoo, roof] ]
			] ]
		] ]
	]
]
\draw[->,>=stealth', ultra thin, gray] (traceV) to[out=south west,in=south,looseness=1] node[pos=0.75,xshift=-10mm]{} (traceT);
\draw[->,>=stealth', ultra thin, gray] (traceT) to[out=south west,in=south,looseness=1] node[pos=0.75,xshift=-10mm]{} (binderT);
\end{forest}
\end{exe}
The syntactic results of inflecting the verbs in these two cases is obviously different, and so a reasonable first question is: could the method of inflecting auxiliary and main verbs be different? Could it, for instance, be that main verbs get inflected by assignment while auxiliary verbs get inflected by moving?

There are some things that weigh against this direction. The conditions we have placed on verb movement are mirrored by conditions on the assignment relationship. Movement of \obar{V} and \obar{I} can move the term to just c-commanding \obar{I}, and never past a c-commanding \obar{X}. And, similarly, the assignment relation always goes from an \obar{I} (or \obar{C}) to a term that it c-commands, and it never crosses a c-commanding \obar{X}. This is indicated by the fact that the parse in (\ref{ex4:agreelow}) cannot produce the outcome in~(\ref{ex4:agreelowa}).
\begin{exe}
\ex \label{ex4:agreelow}
\small
\begin{forest}
[AgrP
	[DP [Tanja, roof]]
	[\ibar{Agr} 
		[\obar{Agr} [\nth{3}\\ sing] ]
		[TP [\ibar{T}
			[\obar{T} [pres] ]
			[VP [\ibar{V}
				[\obar{V} [make] ]
				[VP
					[DP [Jan, roof] ]
					[\ibar{V} [\obar{V} [eat]] ]
				 ]
			] ]
		] ]
	]
]
\end{forest}
\ex[*]{\label{ex4:agreelowa} Jerry make Mary eats.} %\vspace*{-15pt}
\end{exe}
This would be possible if the Agr and T heads of the root clause were able to assign their inflection to the verb that heads the embedded VP. This would be prevented if these heads are blocked from assigning their inflection past the ``closer'' head: \emph{make}. This is quite general: cases where an \obar{I} assigns the inflection it determines on a head that is embedded below another possible target for its inflection do not seem to arise. This looks completely parallel to the Head Movement Constraint, which proscribes movement of a head past a c-commanding head.

That inflectional heads assign the inflection they determine only to verbs they c-command can be appreciated by considering examples like (\ref{ex4:agreenoccom}) \vpageref[below]{ex4:agreenoccom}.
\begin{exe}
\ex \label{ex4:agreenoccom}
\small
\begin{forest}
[AgrP
	[DP [Bart, roof]]
	[\ibar{Agr} 
		[\obar{Agr} [should] ]
		[TP [\ibar{T}
			[\obar{T} [pres] ]
			[VP [\ibar{V}
				[\obar{V} [remember] ]
				[IP 
					[DP [Jill, roof]]
					[\ibar{I}
						[\obar{I} [ing] ]
						[VP [ski, roof] ]
					 ]
				]
			] ]
		] ]
	]
]
\end{forest}
\end{exe}
The grammatical outcome of this structure is one in which the inflectional head determining the suffix \emph{ing} assigns that inflection to the verb it c-commands: \emph{ski}. What can't happen is for this \obar{I} to assign this inflection to \emph{remember}, a verb it does not c-command, but is otherwise a good candidate for being inflected. (\ref{ex4:agreenoccom}) does not lead to the outcome in~(\ref{ex4:agreenoccomo}).
\begin{exe}
\ex[*]{\label{ex4:agreenoccomo} Bart should remembering Jill ski.}
\end{exe}
In general, examples of this sort are not found. If a verb bears inflection that is associated with an \obar{I}, that \obar{I} c-commands the verb. We don't want to design our theory so that the method by which main verbs in English get their inflection is completely divorced from the movement operation that inflects auxiliary verbs. If we did that, we would lose the fact that both Upwards and the Head Movement Constraint seem to govern how both processes work. We should design the system that inflects main verbs in English so that it is, somehow, the same process that is found governing verb movement, the process relevant for inflecting auxiliary verbs.

Note, by the way, that (\ref{ex4:agreenoccom}) confronts us with a question that our move to the Pollockian framework has produced. The root clause in (\ref{ex4:agreenoccom}) contains a modal, and this is a word that we have mapped onto the position that heads sentences. But we now have two heads, and two consequent phrases, that sentences are made from: which of these two should we map modals onto? Modals in English necessarily precede negation, and so this is evidence that modals are in the highest of projections that make up tense/agreement inflection. The evidence we have seen that distinguishes \obar{Agr} from \obar{T} in French doesn't exist in English. We have seen no cases where a verb surfaces in a position higher than its D-structure place but lower than negation in English. When modals or the infinitival \emph{to} appear, we see that \underline{both} tense and agreement morphology are absent. This might be evidence that tense and agreement morphology in English in fact are mapped to the same \obar{X}. This would be consistent with the fact that English verb placement doesn't have the pattern we've seen in French, and that tense and agreement morphology necessarily co-occur. H\"oskuldur Thra\'insson\cite{Thrainsson:1996} has argued that a difference between French and English is how the languages have decided to distribute their morphology across syntactic positions. In French, as we've seen, tense morphology and agreement morphology are inserted into different D-structure positions and brought together in the syntax. Maybe English decides to put tense morphology and agreement morphology together into one \obar{X} position at D-structure, and let the syntax only determine on which verb that morphology gets pronounced. I'll take this view from now on. We can represent English finite clauses as in~(\ref{ex4:TP}).%
%
\marginnote[0pt]{\citet{Bobaljik:1998} argues that whether a language maps tense and agreement onto one, or two, syntactic positions correlates with a number of other differences in their clausal syntax in a way that confrims the idea that this can be the locus of language variation. That languages can differ depending on whether they put two features in different \obar{X} positions at D-structure, or ``bundle'' them together in one \obar{X} has found application in a variety of places. See, for examples, \citet{Giorgi:1996}, \citet{Bianchi:1999}, \citet{Pylkkanen:2008}, \citet{Harley:1989}, \citet{Legate:2014}, and \citet{Erlewine:2018}.}%
%
\begin{exe}
	\raggedright
	\ex \label{ex4:TP}
	\small
	\begin{forest}
	[TP
		[XP ]
		[\ibar{T}
			[\obar{T}
				[\obar{Agr} ]
				[\obar{T} ]
			 ]
			[VP ]
		 ]
	]
	\end{forest}
\end{exe}
This is just like the model of English sentences that we started with, except that I've renamed ``\obar{I}'' ``\obar{T}.'' And to keep the structures of English \obar{T}s parallel to the inflectional head (= \obar{Agr}) that holds tense and agreement morphology in French, I've shown \obar{T} has having these two parts.

For structures like (\ref{ex4:TP}) to be possible D-structures, we will have to add to our inventory of \ibar{X} rules one that allows for an \obar{X} position to have two daughter \obar{X} positions. Let's therefore embellish the rules that generate D-structures so that they include~(\ref{ex4:newXn}).\marginnote[0pt]{Recall that ``\{ $\alpha$, $\beta$ \}'' stands for a linear-order-free relation of sisterhood.}
\pagebreak
\begin{exe}
	\raggedright
	\ex \label{ex4:newX} 
	\begin{xlist}
		\ex XP $\rightarrow$ \{ $\alpha$P, XP\}
		\ex XP $\rightarrow$ \ibar{X}
		\ex XP $\rightarrow$ $\alpha$P \ibar{X}
		\ex \ibar{X} $\rightarrow$ \{ $\alpha$P, \ibar{X} \}
		\ex \ibar{X} $\rightarrow$ \{ $\alpha$P, \obar{X} \}
		\ex \ibar{X} $\rightarrow$ \obar{X}
		\ex \label{ex4:newXn} \obar{X} $\rightarrow$ \{ \obar{$\alpha$}, \obar{X} \}
	\end{xlist}
\end{exe}

With this amendment behind us, let's return to the task of understanding the difference in how main and auxiliary verbs inflect in English. We should endeavor to capture the parallel constraints on the ways that main and auxiliary verbs are inflected in English. It is not clear how that could be done if the means by which they get their inflection is independent. If, on the other hand, the process is the same for both cases, then we can place these constraints on that process and in this way explain why they are shared. This is the direction we will take.

The c-command requirement and the Head Movement Constraint are presently formulated as constraints on the movement relation. That is where they were originally discovered and that is how they were initially framed, as a result. But we can now see that there is reason to rethink this decision. Main verbs do not move, and yet the Head Movement Constraint and the c-command condition nonetheless hold between a main verb and the \obar{I} its inflection derives from. We want the Head Movement Constraint and the c-command conditions to hold, therefore, on the inflection process itself. We then want those constraints to be reflected in how the inflected verbs moves to the \obar{I} that determines their inflection. I will sketch a proposal here that puts together some of the ideas in chapters 3 and 4 of \citet{Chomsky:1995}. It is not a perfectly faithful a reproduction of those ideas, however.%
%
\marginnote[0pt]{The system I will lay out here is modeled after \citet{Adger:2003a}.}%
%

First, let us formulate the condition under which inflectional heads assign verbal inflection: the relation known as \textsc{agree}. We begin by expressing the inflectional morphology involved in terms of features. Let verbs come with features that will determine how their morphological form is determined, and let \obar{Agr} and \obar{T} also come with those features. This allows us to express the assignment of inflectional morphology in terms of matching the relevant features on a verb with the features on \obar{Agr} and \obar{T}. The \textsc{agree} relation specifies the conditions required for this matching of features to occur.\enlargethispage*{\baselineskip}%
%
\marginnote[0pt]{The head with the unvalued features is called the \underline{probe} and the head with the valued features is called the \underline{goal}. \citet{Chomsky:1995,Chomsky:2000}, where this system is laid out, requires that the probe be the c-commanding term, not the c-commanded one. When the probe is the lower of the two in the \textsc{agree} relation, we say that we have ``upwards'' \textsc{agree}. The unvalued features look upwards for their values. See \citet{Zeijlstra:2012} and \citet{Bjorkman:2019} who make the case that all \textsc{agree} relations involve at least an upwards component. In later sections we will see the need for downwards agree as well.}%
%
\begin{exe}
\ex \label{ex4:agree} \textsc{agree} (provisional)\\
Let \obar{X} and \obar{Y} dominate features, and let a feature in \obar{X} have a value and a feature in \obar{Y} be unvalued. A feature in \obar{Y} can get the value of a matching feature in \obar{X} just in case \obar{Y} \textsc{agree}s with \obar{X}. \obar{X} and \obar{Y} \textsc{agree} iff:
\begin{xlist}
	\ex \obar{X} c-commands \obar{Y}, and
	\ex there is no \obar{Z} such that \obar{X} c-commands \obar{Z} and \obar{Z}\\ c-commands~\obar{Y}.
\end{xlist}
\end{exe}
A simple transitive clause would get a representation like that in~(\ref{ex4:agreedemo1}).
\begin{exe}
\ex \label{ex4:agreedemo1}%
%
\marginnote[20pt]{I have collapsed the internal structure of \obar{T} in this representation, and simply listed the two feature bundles---those associated with \obar{T} and those associated with \obar{Agr}---inside the highest \obar{T} node. I will do this when this space-saving doesn't create problems.}%
%
\small
\begin{forest}
[TP
	[DP [they, roof]]
	[\ibar{T} 
		[\obar{T} [agr: \emph{\nth{3}}\\ tense: \emph{past}] ]
		[VP [\ibar{V}
			[\obar{V} [$\surd$eat\\ agr:\\ tense:] ]
			[DP [nattoo, roof] ]
		] ]
	]
]
\end{forest}
\end{exe}
The verb in this example has features for agreement and tense which, when they are given values, will let the process of lexical insertion know which lexical item to put into this position. Similarly, the \obar{T} position is occupied by something that has features for agreement and tense and these are each given a value.

Because \obar{T} c-commands \obar{V}, and there is no head between them, the values for the agreement and tense features can value the agreement and tense features on \obar{V}. This produces~(\ref{ex4:agreedemo1a}).
\begin{exe}
\ex \label{ex4:agreedemo1a}
\small
\begin{forest}
[TP
	[DP [they, roof]]
	[\ibar{T} 
		[\obar{T} [agr: \emph{\nth{3},plur}\\ tense: \emph{past}] ]
		[VP [\ibar{V}
			[\obar{V} [$\surd$eat\\ agr: \emph{\nth{3},plur}\\ tense: \emph{past}] ]
			[DP [nattoo, roof] ]
		] ]
	]
]
\end{forest}
\end{exe}
The values of inflectional features are assigned down to the verb, which then expresses those feature values morphologically.

We can now model the disruptive effect that NegP has on inflecting main verbs in terms of the ``Head Movement Constraint,'' which is now expressed as a condition on \textsc{agree}. If we assume that the head of NegP does not come with the tense or agreement features, a sentence with NegP would get a representation like~(\ref{ex4:negpagree}). %\enlargethispage*{20pt}
\begin{exe}
\ex \label{ex4:negpagree}
\small
\begin{forest}
[TP
	[DP [they, roof]]
	[\ibar{T} 
		[\obar{T} [agr: \emph{\nth{3},plur}\\ tense: \emph{past}] ]
		[NegP
			[AdvP [not, roof] ]
			[\ibar{Neg}
				[\obar{Neg} [\emph{neg}] ]
				[VP [\ibar{V}
					[\obar{V} [$\surd$eat\\ agr: \\tense:] ]
					[DP [nattoo, roof] ]
				] ]
			 ]
		 ]
	]
]
\end{forest}
\end{exe}
The values on the tense and agreement features that \obar{T} has cannot be given to the unvalued agreement and tense features on \obar{V} because \obar{Neg} is a head that intervenes. As a consequence, the agreement feature on the verb will also not be valued. The ungrammaticality of (\ref{ex4:negpagree}), then, can be credited to the inability of this system to inflect the verb with both tense and agreement when NegP is present.

To make explicit the requirement that values of the features on \obar{T} get expressed on a verb, let's adopt~(\ref{ex4:sff}).
\begin{exe}
\ex \label{ex4:sff} The Stray Feature Filter\\
The values of a feature must be morphologically expressed. The morphologically interpreted features must be valued.
\end{exe}
The Stray Feature Filter requires, first, that features with values must have a morphological expression. In general terms, this can be seen as the requirement that every morpheme in a sentence have the exponent which expresses that morpheme in the string of words that results. And, second, the Stray Feature Filter requires that morphemes that come with features must have the information in them that determines how those features should be expressed. That information is the values that features control. 

The ungrammaticality of (\ref{ex4:negpagree}), then, comes from the fact that the value for the agreement feature is not found on a term that can express that value: it is not assigned to a verb, and only verbs in English are equipped with the morphology to express agreement. Moreover, the features on the verb, the ones that will express tense and agreement information, are not valued. The condition in (\ref{ex4:sff}) is just a restatement of \citet{Lasnik:1981}'s Stray Affix filter. His condition said that the bound morphology placed in an inflectional head could not be expressed without the aid of a verbal stem in that very same inflectional head. The Stray Feature Filter does nothing more than translate that into a framework which expresses inflectional morphology with features; (\ref{ex4:sff}) says that the inflectional feature values carried by some inflectional head must find a way of being assigned to a verbal stem that can express them. I've named it the Stray Feature Filter to remind us of this connection.
% //todo smooth out transition here

Let's now consider how this system will work with auxiliary verbs. A simple clause will have the structure in~(\ref{ex4:agreeaux1}).  %\enlargethispage*{2pt}
\begin{exe}
\ex \label{ex4:agreeaux1}
\small
\begin{forest}
[TP
	[DP [they, roof]]
	[\ibar{T} 
		[\obar{T} [agr: \emph{\nth{3},plur}\\ tense: \emph{past}] ]
		[VP [\ibar{V}
			[\obar{V} [have\\agr: \\tense:] ]
			[IP [eaten nattoo, roof] ]
		] ]
	]
]
\end{forest}
\end{exe}
The \textsc{agree}ment relation that we saw in the main verb case will apply here too, producing the representation in~(\ref{ex4:agreeaux2}).
\begin{exe}
\ex \label{ex4:agreeaux2}
\small
\begin{forest}
[TP
	[DP [they, roof]]
	[\ibar{T} 
		[\obar{T} [agr: \emph{\nth{3},plur}\\ tense: \emph{past}] ]
		[VP [\ibar{V}
			[\obar{V} [have\\ agr: \emph{\nth{3},plur}\\ tense: \emph{past}] ]
			[IP [eaten nattoo, roof] ]
		] ]
	]
]
\end{forest}
\end{exe}
This representation satisfies all of our present conditions and so should lead to a grammatical string. As we have seen, however, this isn't what we want. Auxiliary verbs surface in \obar{T} in these contexts. There must be some property that auxiliary verbs have that prevents them from surfacing in their D-structure position when they are inflected. I don't know what this property is, but it is possible, nonetheless, to characterize the consequences it has. It requires the auxiliary to adjoin to the inflectional heads which are responsible for giving values to the features the auxiliary contains. If we call this mysterious property: $\bullet$, then we can state this requirement with~(\ref{ex4:eppstar}).
\begin{exe}
\ex \label{ex4:eppstar} If a morpheme or feature has $\bullet$, then the \obar{X} containing that morpheme or feature must be pronounced inside a \obar{Y} that also contains whatever \obar{X} \textsc{agree}s with.%
%
\marginnote[-105pt]{(\ref{ex4:eppstar}) says the same kind of thing that the Word Criterion expresses. The Word Criterion combines the Stray Feature Filter with the condition on $\bullet$s expressed here. The Word criterion requires that all of the parts of a word that have been put in separate places at D-structure find themselves in one \obar{X} by S-structure. The $\bullet$ expresses the same requirement: the two features that are in an \textsc{AGREE} relation find themselves within one \obar{X} by S-structure. What we learn from the inflection of main verbs in English, is that not all affix/stem pairs are subject to this requirement. The $\bullet$ expresses which pieces of a word are subject to this condition. A feature that has $\bullet$ is what was called a ``strong'' feature in the previous section.}%
%
\end{exe}
Auxiliary verbs have $\bullet$, and so representations like (\ref{ex4:agreeaux2}) will not be well-formed S-structures. Instead, the head movement operation must kick in to produce~(\ref{ex4:agreeaux3}).
\begin{exe}
\ex \label{ex4:agreeaux3}
\small
\begin{forest}
[TP
	[DP [they, roof]]
	[\ibar{T} 
		[\obar{T}
			[\obar{V}, name=binder [$\bullet$have\\ agr: \emph{\nth{3},plur}\\ tense: \emph{past}] ]
			[\obar{T} [agr: \emph{\nth{3},plur}\\ tense: \emph{past}] ]
		 ]
		[VP [\ibar{V}
			[{\ }, name=trace ]
			[IP [eaten nattoo, roof] ]
		] ]
	]
]
\draw[->,>=stealth', ultra thin, gray] (trace) to[out=south west,in=east,looseness=1] node[pos=0.75,xshift=-10mm]{} (binder);
\end{forest}
\end{exe}
\bigskip
%
\marginnote[-40pt]{To my knowledge, this method of forcing auxiliary verb movement in English is not found in the literature. It seems to me, however, a natural way of expressing this fact about auxiliary verbs with a technology that involves $\bullet$s.}%
%
The fact that auxiliary verbs differ from main verbs in needing to move is now coded in the $\bullet$ property. We should hope to find what the $\bullet$ property is, and it seems likely that it has something to do with how the inflectional morphology is expressed since it is a requirement that makes reference to the features expressed by that morphology. One possibility would be to let auxiliary verbs differ from main verbs in the kinds of features they come with.%
%
\marginnote[0pt]{See \citet[]{Lasnik:2003e} for this idea.}%
%
If, for instance, auxiliary verbs are not equipped with the tense and agreement features, but the morphological forms they have nonetheless express those features, then the Stray Feature Filter could not be satisfied by way of \textsc{agree}ment. There will be no way of passing the values of tense and agreement onto the auxiliary verbs on this view because auxiliary verbs do not come with those features. If the Stray Feature Filter can be satisfied by letting the syntax form a word in which the tense and agreement feature values are expressed, then we could see head movement as the sole agent responsible for forming that word in the case of auxiliary verbs. The \obar{T} in (\ref{ex4:agreeaux3}), on this view, would satisfy the Stray Feature Filter. This interpretation of the $\bullet$ property would amount to letting the features and the morphology that expresses those features be packaged in two different ways. In one, the features and the morphology are separated in D-structures and are brought together by syntactic means. This is what happens with auxiliary verbs in English. In the other, the morphology and the features are packaged together into one lexical item at D-structure, but the values those features acquire are determined by \textsc{agree} with another morphological item that resides elsewhere in the syntactic representation. This is what happens with main verbs in English.

While this seems to me a natural way of expressing what the $\bullet$ property is in these scenarios, I'll not commit to that account in these notes. Instead, I'll continue to leave this property unaccounted for, and rely just on the descriptive consequences it has.

The $\bullet$ property can be used to describe what it is that makes auxiliary verbs, unlike main verbs, move. But we must still explain why main verbs are unable to move. One idea that has had found some favor is to posit a force that disfavors movement.%
%
\marginnote[-20pt]{Chomsky assumes that main verbs inflect by moving ``covertly'' to the inflectional heads that hold the relevant morphology. He then assigns different weights to covert and overt movement so that covert movement is favored. A difficulty with his proposals that has not yet been overcome is defining a notion of ``covert'' movement that can do what is required of it in other contexts. Covert movement is thought to arise in other contexts, but in those contexts it always has semantic consequences. Covert main verb movement, by contrast, does not. It appears that covert movement should be designed in such a way that it must have semantic consequences, and it is therefore not a suitable mechanism for inflecting main verbs. See \citet{Stechow:1996} and \citet{Fox:2000} for some discussion.}%
%
\citet{Chomsky:1995} suggests that something of this sort is responsible for blocking movement of main verbs. He observes that we could relate the fact that main verbs are not required to move, that is, they don't have the $\bullet$ property, to the fact that they cannot move if we could build into the grammar a force that allows movement only when it is required.%
%
\marginnote[30pt]{See also \citet{Lasnik:1995} and \citet{Stroik:1999a} for an exploration of the idea that movement cannot be optional.}%
%
Chomsky's own proposal is built upon a different account of how main verbs get inflected, and so will not work for us. But David Pesetsky\cite{Pesetsky:1989} has proposed a way of expressing this notion that is very close to Chomsky's and fits our account of verbal inflection.

To see Pesetsky's idea, it'll be useful to recall what our definition of a sentence is. When we first encountered movement operations in the last chapter, I laid out how we can define the notion of a sentence that involves the collection of parses that we've been using ever since. We allowed for a sentence to involve a derivation, which we think of as a series of parses created by the transformational rules we've been developing. The definition of a sentence is the set of such derivations for a particular D-structure/S-structure pair. The definition is repeated here.
\begin{exe}
	\raggedright
	\ex \label{ex4:deriv4}
	$\mathscr{D}$ = (P_1, P_2\ldots,P_n) is a derivation iff:
	\begin{xlist}
		\ex P_1 is formed by the \ibar{X} Skeleton and Lexical Insertion, and it is semantically interpreted,
		\ex Each P_{i+1} is formed from P_{i} by a rule,
		\ex P_n is subject to the linearization algorithm.
	\end{xlist}
	\ex \label{ex4:sendef} A sentence is the set of derivations which have the same D-structure (P_1) and S-structure~(P_n).
\end{exe}
When more than one rule is used to convert a D-structure parse into an S-Structure parse there are, in principle, more than one way of applying those rules. For instance, if a D-structure requires Rule 1 and Rule 2 to be applied to create a licit S-Structure, then that sentence has potentially two derivations: one that applies Rule~1 before Rule~2, and another that applies Rule~2 before Rule~1. What (\ref{ex4:deriv4}b) allows is for us to think of the sentence with that D-structure as including both of these derivations. A sentence, according to (\ref{ex4:deriv4}) is all the ways the rules of syntax can be applied to a given D-structure so that a particular S-structure is achieved.

Pesetsky's idea is to put a constraint on this set of derivations. Because the aim of this constraint is to control which S-structures are possible, it requires a slight modification of~(\ref{ex4:deriv4}b). Rather than require a sentence with a certain D-structure to have a specific S-structure, Pesetsky's condition on derivations will determine which S-structures are possible from a given D-structure. So first we change (\ref{ex4:sendef}) to~(\ref{ex4:newsendef}).
\begin{exe}
	\raggedright
	\ex \label{ex4:newsendef}
	A sentence is the set of derivations which have the same D-structure and some well-formed S-Structure.
\end{exe}
Now we can give Pesetsky's condition, which will select from the various possible S-structures some D-structure can produce, ones that are allowed. Pesetsky calls his condition ``Earliness,'' and I will formulate it with~(\ref{ex4:earliness}).\marginnote[20pt]{Earliness comes from the unpublished \citet{Pesetsky:1989}, which can be found at his website. This is not his formulation of the condition, but I think it expresses his idea. See also \citet{Pesetsky:1998} and \citet{Pesetsky:2001} for other formulations. A derivation that is not yet blocked but should be involves moving a main verb to \obar{Neg}, where it will be close enough to AGREE with \obar{T}. That should allow a main verb to move past negation into \obar{T} in the same way that auxiliary verbs can (see below). One way of blocking that derivation would be to adopt a condition that requires that every step in a derivation to repair a violation. The movement of a main verb to \obar{Neg} would not repair a violation; it would be a step that puts the verb into a position in which a violation (of the Stray Feature Filter) could be repaired. See \citet{Takahashi:2006b} for a version of Earliness that has this consequence.}
\begin{exe}
\ex \label{ex4:earliness} Earliness\\
Let the set of derivations, $\mathscr{D} = \{ D, D_1,\ldots,D_n \}$, be a sentence. There can be no $D_i \in \mathscr{D}$ if there is a $D_j$ in $\mathscr{D}$ such that $|D_i|\ >\ |D_j|$, where ``$|D|$'' is the number of elements in~$D$.
\end{exe}
What Earliness requires is that derivations from D-structure to S-structure minimize the number of parses in the derivation. If it is possible to satisfy the constraints on S-structure without moving, for instance, then movement is blocked. This is precisely what happens in the inflection of main verbs. They are able to satisfy the Stray Feature Filter by way of \textsc{agree}, and this makes movement of main verbs to \obar{T} superfluous.

This is a popular account of the failure of main verb movement, but it is not entirely successful. This becomes apparent when the cases in which an auxiliary verb combines with NegP are considered. In such cases, we have the D-structure representation in~(\ref{ex4:negauxin}).
\begin{exe}
\ex \label{ex4:negauxin}
\small
\begin{forest}
[TP
	[DP [they, roof]]
	[\ibar{T} 
		[\obar{T} [agr: \emph{\nth{3},plur}\\ tense: \emph{past}] ]
		[NegP
			[AdvP [not, roof]]
			[\ibar{Neg}
				[\obar{Neg} [neg] ]
				[VP [\ibar{V}
					[\obar{V} [$\bullet$$\surd$have\\ agr: \\ tense:] ]
					[IP [eaten nattoo, roof] ]
				] ]
			 ]
		 ]
	]
]
\end{forest}
\end{exe}
As we have seen in the case of main verbs, \textsc{agree} is prevented from giving the value that \obar{T} has for the agreement feature to the verbs beneath \obar{T}. Just as in the case with main verbs, then, this representation does not satisfy the Stray Feature Filter. Instead, as we know, the auxiliary verb moves to the \obar{T} position. It seems reasonable to assume that it's the ability of the auxiliary verb to move that allows it to get inflected. We want movement and \textsc{agree} to have different outcomes in this circumstance.

I don't know of a working solution to this question that has gained currency in the literature. I will make one up here, just so that we have something that works. We are looking for something that distinguishes auxiliary verbs from other verbs that allows just auxiliary verbs to move to \obar{T} when \obar{Neg} is present. One thing that distinguishes auxiliary verbs in English from others that might be relevant is that auxiliary verbs have negative forms, and others don't.
\begin{exe}
	\raggedright
	\ex \label{ex4:negaux}
	\begin{xlist}
		\ex Tanja hadn't run today.
		\ex Karen isn't running today.
	\end{xlist}
	\ex \label{ex4:negv}
	\begin{xlist}
		\ex[*]{Tanja ran't today.}
		\ex[*]{Karen runsn't today.}
	\end{xlist}
\end{exe}
One might be tempted to think that \textipa{[\textsyllabic{n}t]} that is found on these auxiliaries is mere \emph{not} in some prosodically reduced form. If that were the case, then one could explain why only auxiliary verbs in English have these forms with the fact that only auxiliary verbs in English are capable of showing up to the left of \emph{not}, thereby feeding the reduction process. But as Stephen Anderson\cite{Anderson:2008a} notes, the \textipa{[\textsyllabic{n}t]} forms of auxiliary verbs move as heads in contexts where \obar{T} moves to~\obar{C}.
\begin{exe}
	\raggedright
	\ex \label{ext4:ntmove} 
	Haven't you finished this chapter yet?\\
	Wasn't she happy with the solution?
\end{exe}
And this distinguishes these cases from other cases of reduction, which don't have this property of being able to feed movement. For example, auxiliary verbs themselves can be reduced onto certain DPs in Specifier of TP position; this is illustrated in~(\ref{ex4:contrac}).
\begin{exe}
	\raggedright
	\ex \label{ex4:contrac} 
	\begin{xlist}
		\ex She said that you've been tense.
		\ex Who thought that who's pleasant?
		\ex She thinks that he's pleasant. 
	\end{xlist}
\end{exe}
The forms \emph{you've}, \emph{he's} and \emph{who's} prosodically behave like a single unit. But the syntax does not recognize them as a constituent, as can be seen by noting that Wh Movement isn't able to move them.
\begin{exe}
	\raggedright
	\ex \label{ex4:contracm} 
	\begin{xlist}
		\ex[*]{\linkto who've \under did you think \from{\ \ } been tense?}
		\ex[*]{\linkto who's \under did she say \from{\ \ } pleasant?}
	\end{xlist}
\end{exe}
If we treat forms like \emph{isn't} and \emph{haven't} as single \obar{X}s, this is explained. Under our present set of assumptions, for something to be an \obar{X} requires that they be a word. 

Imagine, then, that auxiliary verbs can come with a feature that expresses its polarity, as indicated in~(\ref{ex4:auxlex}).%
%
\marginnote[0pt]{I'll use the ``$\Rightarrow$'' symbol to represent Vocabulary Insertion, the S-structure process that matches a word to clusters of morphemes and features in an~\obar{X}.}%
%
\begin{exe}
	\raggedright
	\ex \label{ex4:auxlex} 
	\small
	\begin{forest}
	[\obar{V} [$\surd$have\\ tense: \emph{pres}\\ agree: \emph{\nth{3},sing}\\ pol: \emph{pos}]
	]
	\end{forest}
	$\Rightarrow$ \emph{has}
	\begin{forest}
	[\obar{V}
		[$\surd$have\\ tense: \emph{pres}\\ agree: \emph{\nth{3},sing}\\ pol: \emph{neg} ]
	]
	\end{forest}
	$\Rightarrow$ \emph{hasn't}
	\begin{forest}
	[\obar{V}
		[$\surd$be\\ tense: \emph{pres}\\ agree: \emph{\nth{3},sing}\\ pol: \emph{neg} ]
	]
	\end{forest}
	$\Rightarrow$ \emph{isn't}
\end{exe}

Consider now how the syntax for negated auxiliaries proceeds.
\bigskip
\begin{fullwidth}
\begin{exe}
	\raggedright
	\ex \label{ex4:negauxt}
	\footnotesize
	\begin{forest}
	[TP
		[DP [Tanja, roof]]
		[\ibar{T} 
			[\obar{T}\\ agr: \emph{\nth{3},sing}\\ tense: \emph{past} ]
			[NegP [\ibar{Neg} 
				[\obar{Neg}\\ pol: \emph{neg} ]
				[VP [\ibar{V}
					[\obar{V} [$\bullet$$\surd$have\\ agr:\\ tense:\\ pol:] ]
					[VP [\ibar{V} [run, roof]] ]
				] ]
			]
		]
	]]
	\end{forest}
	$\longrightarrow$
	\begin{forest}
		[TP
		[DP [Tanja, roof]]
		[\ibar{T} 
			[\obar{T}\\ agr: \emph{\nth{3},sing}\\ tense: \emph{past} ]
			[NegP [\ibar{Neg} 
				[\obar{Neg}\\ pol: \emph{neg} ]
				[VP [\ibar{V}
					[\obar{V} [$\bullet$$\surd$have\\ agr:\\ tense:\\ pol: \emph{neg}] ]
					[IP [run, roof] ]
				] ]
			]
		]
	]]
	\end{forest}
	$\longrightarrow$
	\begin{forest}
	[TP
		[DP [Tanja, roof]]
		[\ibar{T} 
			[\obar{T}\\agr: \emph{\nth{3},sing}\\ tense: \emph{past} ]
			[NegP [\ibar{Neg}
				[\obar{Neg}
					[\obar{V} [$\bullet$$\surd$have\\ agr:\\ tense:\\ pol: \emph{neg}, name=binder] ]
					[\obar{Neg}\\ pol: \emph{neg} ]
				 ]
				[VP [\ibar{V}
					[{\ }, name=trace ]
					[IP [run, roof] ]
				] ]
			] ]
		]
	]
	\draw[->,>=stealth', ultra thin, gray] (trace) to[out=south,in=-80,looseness=1] node[pos=0.75,xshift=-10mm]{} (binder);
	\end{forest}

	$\longrightarrow$
	\begin{forest}
		[TP
			[DP [Tanja, roof]]
			[\ibar{T} 
				[\obar{T}\\agr: \emph{\nth{3},sing}\\ tense: \emph{past} ]
				[NegP [\ibar{Neg}
					[\obar{Neg}
						[\obar{V} [$\bullet$$\surd$have\\ agr: \emph{\nth{3},sing}\\ tense: \emph{past}\\ pol: \emph{neg}] ]
						[\obar{Neg}\\ pol: \emph{neg} ]
					 ]
					[VP [\ibar{V}
						[{\ }]
						[IP [run, roof] ]
					] ]
				] ]
			]
		]
	\end{forest}
	$\longrightarrow$
	\begin{forest}
	[TP
		[DP [Tanja, roof]]
		[\ibar{T} 
			[\obar{T}
				[\obar{Neg}, name=binder
					[\obar{V} [$\bullet$$\surd$have\\ agr: \emph{\nth{3},sing}\\ tense: \emph{past}\\ pol: \emph{neg}] ]
					[\obar{Neg}\\ pol: \emph{neg} ]
				 ]
				[\obar{T}\\agr: \emph{\nth{3},sing}\\ tense: \emph{past}  ]
			 ]
			[NegP [\ibar{Neg}
				[{\ }, name=trace ]
				[VP [\ibar{V}
					[{\ } ]
					[IP [run, roof] ]
				] ]
			] ]
		]
	]
	\draw[->,>=stealth', ultra thin, gray] (trace) to[out=south west,in=south,looseness=1] node[pos=0.75,xshift=-10mm]{} (binder);
	\end{forest}
\end{exe}
\end{fullwidth}
\bigskip
In the first step, \textsc{agree} values the pol feature on \emph{have}, and this allows for the second step in which the \obar{V} containing \emph{have} moves to \obar{Neg}. This movement is required by the $\bullet$ feature on \emph{have}. In \obar{Neg} position, \emph{have} is now close enough to \obar{T} for \textsc{agree} to hold, and this what allows for the fourth parse in this derivation, where the tense and agreement features on \emph{have} are valued. This is again recapitulated by movement of \obar{Neg} into \obar{T}, which is again required by the $\bullet$ on~\emph{have}.

This syntax will also underlie~(\ref{ex4:yup}).
\begin{exe}
	\raggedright
	\ex \label{ex4:yup}
	\begin{xlist}
		\ex Smith had not run.
		\ex Jones is not running
	\end{xlist}
\end{exe}
What we see here is that the polarity feature on \emph{have} and \emph{be} can be valued to \emph{neg} and not result in Vocabulary Insertion producing \emph{haven't} and \emph{isn't}. Instead, \emph{not} expones negation. The morphological machinery developed here doesn't provide a way of capturing this. For this proposal to be successful, then, it requires some process that is capable of capturing the dependency between pronouncing \emph{not} and pronouncing~\textipa{\textsyllabic{n}t}. I gamble that such a process exists, and suggest that we adopt this account of the ability auxiliary verbs have in English to move past~\obar{Neg}.

There are problems with Earliness from examples in which head movement appears to be optional. We've seen something with that appearance in two contexts. In French infinitival clauses, we've seen that auxiliary verbs are able to surface in one of two positions: to the left of \emph{pas} or to the right of \emph{pas}. This is illustrated by (\ref{ex4:frenchauxinf}), repeated here.%\enlargethispage*{14pt}
\begin{exe}
	\exr{ex4:frenchauxinf}
	\begin{xlist}
	\ex[]{\gll ne \textit{pas} \textbf{{\^e}tre} heureux est une condition pour {\'e}crire {des romans}\\ ne not be happy is a condition for writing novels\\ `To not be happy is a condition for writing novels.'}
	\ex[]{\gll N'\textbf{\^etre} \textit{pas} heureux est une condition pour \'ecrire {des romans}.\\ ne'be not happy is a condition for writing novels\\ `To not be happy is a condition for writing novels.'}
	\end{xlist}
\end{exe}
Consider the representation given to (\ref{ex4:frenchauxinf}a), in~(\ref{ex4:frenchauxinfr}).
\begin{exe}
\ex \label{ex4:frenchauxinfr}
\small
\begin{forest}
[AgrP
	[DP [\textsc{pro}, roof]]
	[\ibar{Agr} 
		[\obar{Agr} 
			[\obar{Neg} [ne, name=binder]]
			[\obar{Agr}\\agr: ?? ]
		]
		[NegP
			[AdvP [pas, roof] ]
			[\ibar{Neg}
				[{\ }, name=trace ]
				[TP [\ibar{T}
					[\obar{T}
						[\obar{V} [$\surd$\^etre\\ agr:\\ tense: \emph{ir}, name=binderv] ]
						[\obar{T}\\ $\bullet$tense: \emph{ir} ]
					 ]
					[VP [\ibar{V}
						[{\ }, name=tracev ]
						[AP [heureux, roof] ]
					 ]
				] ]
			 ]
		 ]
	]
]]
\draw[->,>=stealth', ultra thin, gray] (trace) to[out=south west,in=-80,looseness=1] node[pos=0.75,xshift=-10mm]{} (binder);
\draw[->,>=stealth', ultra thin, gray] (tracev) to[out=-100,in=-80,looseness=1] node[pos=0.75,xshift=-10mm]{} (binderv);
\end{forest}
\end{exe}
I have given the infinitival inflectional feature inside \obar{T} the $\bullet$ property, as this head forces head movement of both auxiliary and main verbs to it. (Recall: we saw that main verbs in French infinitival clauses move to \obar{T} and are thereby separated from their DP complements by adverbs.) In general, we can characterize those inflectional heads that allow main and auxiliary verbs to move to them, i.e., those heads that are described as being ``strong,'' in Pollock's paper, as having the $\bullet$ property.

Consider now the status of \obar{Agr} in this representation. The inflectional material this head is associated with is obscure; there is nothing in infinitival clauses that looks like the agreement phenomenon we see in finite clauses. I will indicate the value of the agreement morphology in non-finite \obar{Agr} with ``??,'' as a consequence. In (\ref{ex4:frenchauxinf}a), we don't want this morphology to force movement, so I have placed no $\bullet$ on it or the auxiliary verb. Note that \emph{ne} has moved to \obar{Agr}, and this now allows \obar{Agr} to \textsc{agree} with \obar{T}. The conditions on \textsc{agree} do not allow it to hold between two heads if there is another head intervening. The position \obar{Neg} has at D-structure causes it to intervene between \obar{Agr} and \obar{T}, so at D-structure \textsc{agree} would be blocked between \obar{Agr} and \obar{T}. But once \obar{Neg} moves to \obar{Agr} in the way indicated in (\ref{ex4:frenchauxinf}), this blockage is alleviated and \textsc{agree} is possible. From (\ref{ex4:frenchauxinf}) we can create an S-structure that satisfies the Stray Affix filter by \textsc{agree}.
\begin{exe}
\ex \label{ex4:frenchauxinfra}
\small
\begin{forest}
	[AgrP
	[DP [\textsc{pro}, roof]]
	[\ibar{Agr} 
		[\obar{Agr} 
			[\obar{Neg}, name=probe [ne]]
			[\obar{Agr}\\agr: ?? ]
		]
		[NegP
			[AdvP [pas, roof] ]
			[\ibar{Neg}
				[{\ }, name=trace ]
				[TP [\ibar{T}
					[\obar{T}
						[\obar{V}, name=goal [$\surd$\^etre\\ agr: ??\\ tense: \emph{ir}, name=binderv] ]
						[\obar{T}\\ $\bullet$tense: \emph{ir} ]
					 ]
					[VP [\ibar{V}
						[{\ }]
						[AP [heureux, roof] ]
					 ]
				] ]
			 ]
		 ]
	]
]]
\draw[|-|, dotted] (goal) to[out=south west,in=south west] node[pos=0.5,xshift=-5mm]{\emph{agree}} (probe);
\end{forest}
\end{exe}
This is how (\ref{ex4:frenchauxinf}a) is derived.

But the outcome in (\ref{ex4:frenchauxinf}b) requires a different set of assumptions. In that circumstance, an auxiliary verb that \obar{Agr} \textsc{agree}s with moves to \obar{Agr}. We want this option to only be available to auxiliary verbs. The technology we have allows that to be expressed by giving to auxiliary verbs the $\bullet$. For (\ref{ex4:frenchauxinf}b), then, our representation should be~(\ref{ex4:bulletfraux}).
\begin{exe}
\ex \label{ex4:bulletfraux}
\small
\begin{forest}
[AgrP
	[DP [\textsc{pro}, roof]]
	[\ibar{Agr} 
		[\obar{Agr} 
			[\obar{Neg} [ne, name=binder] ]
			[\obar{Agr} [agr: ??] ]
		 ]
		[NegP
			[AdvP [pas, roof]]
			[\ibar{Neg}
				[{\ }, name=trace ]
				[TP [\ibar{T}
					[\obar{T}
						[\obar{V}, name=binderv [$\bullet$$\surd$\^etre\\ tense: \emph{ir}\\ agr:] ]
						[\obar{T}\\ tense: \emph{$\bullet$ir} ]
					 ]
					[VP [\ibar{V}
						[{\ }, name=tracev ]
						[AP [heureux, roof] ]
					] ]
				] ]
			 ]
		 ]
	]
]
\draw[->,>=stealth', ultra thin, gray] (trace) to[out=south west,in=-70,looseness=1] node[pos=0.75,xshift=-10mm]{} (binder);
\draw[->,>=stealth', ultra thin, gray] (tracev) to[out=south west,in=-50,looseness=1] node[pos=0.75,xshift=-10mm]{} (binderv);
\end{forest}
\end{exe}
Once \textsc{agree} holds between \obar{Agr} and \obar{T} in this configuration, the $\bullet$ value on \emph{\^etre} will require movement, forming the representation in~(\ref{ex4:bulletauxfrm}).
\begin{exe}
	\raggedright
	\ex \label{ex4:bulletauxfrm}
	\small
	\begin{forest}
	[AgrP
		[DP [\textsc{pro}, roof]]
		[\ibar{Agr} 
			[\obar{Agr}
				[\obar{Neg} [ne] ]
				[\obar{Agr} 
					[\obar{T}, name=binder 
					[\obar{V} [$\bullet$$\surd$\^etre\\ tense: \emph{ir}\\ agr:??] ]
						[\obar{T}\\ tense: \emph{$\bullet$ir} ]
					 ]
				 ]
			 ]
			[NegP 
				[AdvP [\textsc{pro}, roof]]
				[\ibar{Neg}
					[{\ } ]
					[TP [\ibar{T}
						[{\ }, name=trace ]
						[VP [\ibar{V}
							[{\ } ]
							[AP [heureux, roof] ]
						] ]
					] ]
				 ]
			 ]
		]
	]
	\draw[->,>=stealth', ultra thin, gray] (trace) to[out=south west,in=south,looseness=1] node[pos=0.75,xshift=-10mm]{} (binder);
	\end{forest}
\end{exe}

The $\bullet$ property requires that a verb expressing the value the agr feature be adjoined to \obar{Agr}, and this results in the auxiliary verb moving to \obar{Agr}. Since this is the only derivation in which the $\bullet$ property gets its requirements satisfied, Earliness permits it. To characterize ``optional movement,'' then, a system which embraces Earliness requires that there be two possible representations: one that forces movement and another in which it is blocked. Under the particulars of the account I am sketching here, that would mean that French speakers who allow the alternation in (\ref{ex4:frenchauxinf}) allow for two kinds of auxiliary verbs: one with $\bullet$ and the other without. Some French speakers report that these two word orders correspond to two registers of French. Perhaps this difference in registers, if it is indeed systematically associated with these two word orders, could be expressed as a kind of difference in vocabulary. In one register, auxiliary verbs have $\bullet$ and in the other, auxiliary verbs do not have $\bullet$.

If these methods for characterizing the flexible position of certain heads can be confirmed, then we can continue to entertain the hypothesis that Earliness exists, and rely on its consequences to block main verbs from moving in English (non-negated) sentences. Earliness has the interesting property of relating the absence of main verb movement in English to the presence of obligatory auxiliary verb movement in English. It, along with the $\bullet$ attribute, expresses the idea that the controls on movement will allow for just two categories: obligatory movement or no movement. With the possible exceptions found in infinitives, this does seem to be the case for all the instances of head movement we've investigated. An open question is whether this holds more generally. All the cases of Argument Movement that we've looked at are obligatory, so those cases seem to fit the prediction. We've also briefly encountered Wh Movement, which is responsible for producing constituent questions of the sort that~(\ref{ex4:whmove}) exemplifies.
\begin{exe}
	\ex \label{ex4:whmove} She asked \linkto what \under Sam gave \from{\ } to him.
\end{exe}
In simple questions, this seems to be a movement that is also obligatory. (\ref{ex4:whmove}) would be ungrammatical if \emph{who} did not move. But in questions where there is more than one moveable term, as in~(\ref{ex4:whmove2}), it is not obligatory for all those terms to move.
\begin{exe}
	\ex \label{ex4:whmove2} She asked what Sam gave to whom.
\end{exe}
In these kinds of questions, the terms that are not required to move are also prevented from moving; (\ref{ex4:whmove3}) is ungrammatical.
\begin{exe}
	\ex[*]{\label{ex4:whmove3} She asked what \linkto whom \under Sam gave to \from{\ }\ \ .}
\end{exe}
So, these cases fit the prediction as well. We need only say that the CPs which are interpreted as constituent questions invoke a requirement that a wh-phrase move to their Specifier position, and Earliness will permit only one wh-phrase to move, since any additional movements will not be needed to satisfy that requirement.

For Argument Movement, Wh Movement and Head Movement in English, then, a grammar that embraces Earliness seems to be consistent with the facts. But it's not hard to find cases that don't conform to Earliness. In many languages, for instance, Wh Movement appears to be optional and there are also languages in which Argument Movement appears to be optional. And Heavy NP Shift in English is not transparently consistent with Earliness either. It is difficult to classify either of the examples in~(\ref{ex4:hnpse}) as ungrammatical, for instance.
\begin{exe}
	\ex \label{ex4:hnpse}
	\begin{xlist}
		\ex She gave a lot of books to me.
		\ex She gave to me a lot of books.
	\end{xlist}
\end{exe}
But if (\ref{ex4:hnpse}b) is derived from (\ref{ex4:hnpse}a) by movement, then Earliness should block it. There are analyses of all these counter-examples that make them consistent with Earliness, or with some modified version of Earliness, but no consensus has formed about whether these analyses are successful. The validity of Earliness, then, is an open question, under active investigation.

In any case, let's adopt this system. Let's assume that verbs inflect by way of the \textsc{agree} relation, and that the Head Movement Constraint is a constraint on this relation and not movement. As we've seen, in order to spread this constraint on \textsc{agree} to cases where head movement arises, we need a method of explaining how head movement is tied to \textsc{agree}. Head movement needs to be seen as one outcome of the \textsc{agree} relation, and therefore an outcome that inherits the constraints on \textsc{agree}. I've sketched a way of doing that in this section, but it's clear that this is not yet a full account. It relies on finding a source for the $\bullet$ property, and for understanding what prevents main verbs from moving through \obar{Neg} into higher heads in English. And it relies on Earliness, which has effects across the grammar that have yet to be confirmed.

\section{\textsc{agree} and $\theta$-roles} % (fold)
\label{sec:theta}

Now that we've developed a vocabulary for describing how morphological information is spread across the heads of a syntactic representation, let's revisit the issue we left at the end of the previous chapter concerning $\theta$-roles. We discovered there that some things fell into place, Burzio's Generalization for instance, if we took the external $\theta$-role associated with the meaning of a verb to be syntactically separated from rest of the ingredients in the meaning of a verb. The syntactic representation we assigned to verbs that have external $\theta$-roles is like that in~(\ref{ex4:tr1}).
\begin{exe}
	\raggedright
	\ex \label{ex4:tr1} 
	\small
	\begin{forest}
	[TP [\ibar{T}
		[\obar{T} [will]]
		[vP
			[DP [Noah, roof]]
			[\ibar{v}
				[\obar{v} [agent] ]
				[VP [\ibar{V}
					[\obar{V} [$\surd$eat] ]
					[DP [nattoo, roof] ]
				] ]
			 ]
		 ]
	]]
	\end{forest}
\end{exe}
The choice of verb that gets inserted into \obar{V} position is tied to the choice of \mbox{$\theta$-role} that is inserted into \obar{v} position, and this connection seems to be arbitrary. There has not been a pattern established between the nature of the verb in \obar{V} and the nature of the $\theta$-role in \obar{v}. So far as I know, this is simply learned: \emph{eat} is obligatorily accompanied by an Agent $\theta$-role in \obar{v}, whereas, \emph{like} is obligatorily accompanied by an Experiencer $\theta$-role in \obar{v}, and so on. To the extent that this connection is arbitrary, it suggests that there is a morphological connection between the $\theta$-role in \obar{v} and the verb in \obar{V}. We want the connection between the morphology in \obar{v} and the morphology in \obar{V} to be the one that a word makes. Now that we have technology for that, let's see how it might be applied to this case.

When we first looked at the denotation of verbs, we embraced the idea that verbs have two pieces to their meaning. One of those pieces describes the activity, state of affairs, or what have you. The other piece describes the $\theta$-roles that verb uses to relate its arguments to that state of affairs. We've since spelled out this idea in terms of an event semantics. Verbs include information that determines what sort of event they describe and also the $\theta$-roles that relate the arguments the syntax will supply to that event. Imagine that these are independent pieces to the vocabulary item that correspond to the words \emph{eat}, \emph{like}, etc. Imagine, that is, that lexical insertion populates D-structures that have a position for a verb with these pieces. Setting aside the features that represent tense/agreement inflection for a moment, the proposal is that \obar{V} positions like~(\ref{ex4:thetfeat1}) can be produced at D-structure.
\begin{exe}
	\raggedright
	\ex \label{ex4:thetfeat1} 
	\small
	\begin{forest}
	[VP[\ibar{V}
		[\obar{V} 
			[\obar{V} [$\surd$visit] ]
			[\obar{v} [patient] ]
		]
		[DP [Adam, roof] ]
	]]
	\end{forest}
\end{exe}
I've called the \obar{X} position that \textsc{patient} is inserted into \obar{v}, because that's what we've called the head position that external $\theta$-roles appear in. I've represented the portion of the verb that corresponds the event description with $\surd$visit. This is often called the ``root'' of the verb. If we give to the Patient $\theta$-role the same kind of denotation we adopted by external $\theta$-roles, then the semantic contribution of each of the pieces in the maximal \obar{V} position of (\ref{ex4:thetfeat1}) is~(\ref{ex4:tfeatden1}).%
%
\marginnote[0pt]{If we take the meanings of verbs to be derived syntactically in the way that is illustrated here, we will have to change our conception of the rule of Existential Closure introduced in the previous chapter. That rule was offered as a rule that changed lexical items---verbs---by binding off their internal argument with an existential operator. The view introduced here is that verbs are the way in which morphologically complex \obar{V}s are pronounced. They are not themselves morphemes with idiosyncratic denotations, and so they cannot be the input of a rule that builds new words from old ones. Instead, we'll have to take Existential Closure to be a syntactic rule that applies to \obar{V}s.}%
%
\begin{exe}
	\raggedright
	\ex \label{ex4:tfeatden1} 
	\begin{xlist}
		\ex \denotes{$\surd$visit} = $\lambda \epsilon$ \textsc{visit}($\epsilon$)
		\ex \denotes{patient} = $\lambda P$ $\lambda x$ $\lambda \epsilon$ $P(\epsilon) \wedge$ \textsc{patient}($x$)($\epsilon$)
	\end{xlist}
\end{exe}
Notice that these denotations, along with our rules of semantic composition, produce the correct meaning with this syntax.
\begin{exe}
	\raggedright
	\ex \label{ex4:tfeatden2} 
	\small
	\begin{forest}
	[$\lambda \epsilon$ \textsc{visit}($\epsilon$) $\wedge$ \textsc{patient}(Adam)($\epsilon$)\\ VP [$\lambda \epsilon$ \textsc{visit} $\wedge$ \textsc{patient}(Adam)($\epsilon$)\\\ibar{V}
		[$\lambda x \lambda \epsilon$ \textsc{visit}($\epsilon$) $\wedge$ \textsc{patient}($x$)($\epsilon$)\\\obar{V} 
			[$\lambda \epsilon$ \textsc{visit}($\epsilon$)\\ \obar{V} [$\surd$visit]]
			[$\lambda P \lambda x \lambda \epsilon$ $P(\epsilon)$ $\wedge$ \textsc{patient}($x$)($\epsilon$)\\\obar{v} [patient]]
		 ]
		[DP [Adam, roof] ]
	]]
	\end{forest}
\end{exe}
It's possible to see the meanings of verbs, then, as being constructed in the syntax and matched to the vocabulary items we recognize as verbs. The idea I'll sketch here is that external $\theta$-roles are simply a special case of the situation illustrated in~(\ref{ex4:thetfeat1}).

Notice that when there are two $\theta$-roles, they cannot semantically combine with the rules we adopted, if they do so within one \obar{X}. (\ref{ex4:ppfl1}) is illicit.
\begin{exe}
	\raggedright
	\ex \label{ex4:ppfl1}
	\hspace*{-25pt}
	\small
	\begin{forest}
	[VP[\ibar{V}
		[\ibar{V}
			[??\\ \obar{V}
				[$\lambda x$ $\lambda \epsilon$ \textsc{visit}($\epsilon$) $\wedge$ \textsc{patient}($x$)($\epsilon$)\\ \obar{V} 
					[$\lambda \epsilon$ \textsc{visit}($\epsilon$)\\ \obar{V} [$\surd$visit] ]
					[$\lambda P$ $\lambda x$ $\lambda \epsilon$ $P(\epsilon)$ $\wedge$ \textsc{patient}($x$)($\epsilon$)\\ \obar{v} [patient] ]
				]
				[$\lambda P$ $\lambda x$ $\lambda \epsilon$ $P(\epsilon)$ $\wedge$ \textsc{agent}($x$)($\epsilon$) \\ \obar{v} [agent] ]
			 ]
			[DP [Adam, roof]]
		 ]
		[DP [Taka, roof] ]
	]]
	\end{forest}
\end{exe}
The denotations for both Patient and Agent requires that they combine with something whose denotation is a function whose domain is the set of eventualities. That requirement is satisfied for the Patient, but not for Agent. The \obar{V} Agent combines with is a function whose domain is the set of things, not the set of eventualities. These two denotations cannot combine. When a $\theta$-role combines with a root, it creates something whose meaning can combine with the argument of the $\theta$-role -- it creates a verb whose meaning is a function from things to descriptions of eventualities. That isn't the right kind of meaning to combine with another $\theta$-role.

What this demonstrates is that if we take verbs to be matched to a \obar{X} position in which the $\theta$-roles and eventuality descriptions have been assembled by Lexical Insertion, we can use the rules of semantic composition to constrain the meanings that verbs have. If we have the denotation for $\theta$-roles right, and there are no other relevant rules of semantic composition, one prediction is that a verb can assign no more than one $\theta$-role. We have taken $\theta$-roles to mediate the connection between verbs and their DP and PP arguments, so the fact that verbs can combine with more than one DP and PP argument shows that something about these assumptions is incorrect. If $\theta$-roles only mediate the combination of verbs with DPs, however, we are closer to the mark. Setting aside specialized constructions, English verbs allow for only one internal DP argument.%
%
\marginnote[0pt]{One of the specialized constructions is the so-called double object construction: ``Helen gave Matt a book.'' We return to look at these cases in chapter \ref{ch:compstructure}.}

Let us provisionally adopt the view, then, that PPs, along with CPs and other clause-like objects, do not semantically combine with their verbs by way of $\theta$-roles. If that is correct, the syntax-semantics given to $\theta$-roles derives the fact that for a verb to combine with more than one DP argument, one of those $\theta$-roles will have to be external. In the case of \emph{visit}, that $\theta$-role is Agent.

What is required, then, is to find a syntax in which the denotation of the Agent $\theta$-role is put in a place different than where the other other pieces of a verb's meaning are assembled. This is precisely the same situation that arises with tense inflection. The denotation of tense needs to be in a place different than where its exponent is, and \textsc{agree} is the way this is achieved. We'll allow $\theta$-roles to be features, then. A verb that comes with two $\theta$-roles will allow one of those $\theta$-roles to be inserted into the \obar{V} position valued. The other $\theta$-role will be placed in a position that gives it the right kind of semantic argument, but distances it from the verb---its exponent. Thus, in the \obar{V} position will be an unvalued $\theta$-role feature that gets into the \textsc{agree} relation with the external $\theta$-role, thereby allowing the external $\theta$-role to be exponed, and satisfy the Stray Feature Filter. 

Reverting back to the truncated way I've represented \obar{X} positions with complex contents, we can represent the D-structure scenario just described with~(\ref{ex4:transdecom}).
\begin{exe}
	\raggedright
	\ex \label{ex4:transdecom} 
	\small
	\begin{forest}
	[TP
		[DP [Taka, roof, name=binder]]
		[\ibar{T} 
			[\obar{T} [can] ]
			[vP
				[{\ }, name=trace ]
				[\ibar{v}
					[\obar{v} [$\theta$-role: \emph{agent}] ]
					[VP [\ibar{V}
						[\obar{V} [$\surd$visit\\ $\theta$-role: \emph{patient}\\ $\theta$-role:] ]
						[DP [Adam, roof] ]
					] ]
				 ]
			 ]
		]
	]
	\draw[->,>=stealth', ultra thin, gray] (trace) to[out=south west,in=south,looseness=1] node[pos=0.75,xshift=-10mm]{} (binder);
	\end{forest}
\end{exe}
Imagine that the Spell Out rules for \emph{visit} are~(\ref{ex:eat1}).
\begin{exe}
	\raggedright
	\ex \label{ex:eat1} 
	{} [_{\text{\obar{V}}} $\surd$visit+\emph{patient}+\emph{agent} ] $\Rightarrow$ \textipa{/vIzIt/}
\end{exe}
This will require \textsc{agree} to produce from (\ref{ex4:transdecom}),~(\ref{ex4:transdecom2}).
\begin{exe}
	\raggedright
	\ex \label{ex4:transdecom2} 
	\small
		\begin{forest}
			[TP
				[DP [Taka, roof, name=binder]]
				[\ibar{T} 
					[\obar{T} [can] ]
					[vP
						[{\ }]
						[\ibar{v}
							[\obar{v}, name=probe [$\theta$-role: \emph{agent}] ]
							[VP [\ibar{V}
								[\obar{V}, name=goal [$\surd$visit\\ $\theta$-role: \emph{patient}\\ $\theta$-role: \emph{agent}] ]
								[DP [Adam, roof] ]
							] ]
						 ]
					 ]
				]
			]
			\draw[|-|, dotted] (goal) to[out=south west,in=south west,looseness=2.3] node[pos=0.5,xshift=-5mm]{\emph{agree}} (probe);
	\end{forest}
\end{exe}
This representation satisfies the constraints on S-structures -- in particular the Stray Feature Filter -- and connects the presence of the verb's external $\theta$-role in a position outside of the VP headed by that verb with the verb. And, as we saw in the previous chapter, our semantics correctly combines the denotations of the morphemes in~(\ref{ex4:transdecom2}). We can see this by considering what the D-structure of the vP in (\ref{ex4:transdecom2}) looks like:
\begin{exe}
	\raggedright
	\ex \label{ex4:transdecom2t} 
	\small
	\begin{forest}
	[$\lambda \epsilon$ \textsc{visit}($\epsilon$) $\wedge$ \textsc{patient}(Adam)($\epsilon$) $\wedge$ \textsc{agent}(Taka)($\epsilon$)\\ vP
		[DP [Taka, name=binder, roof] ]
		[$\lambda x$ $\lambda \epsilon$ \textsc{visit}($\epsilon$) $\wedge$ \textsc{patient}(Adam)($\epsilon$) $\wedge$ \textsc{agent}($x$)($\epsilon$)\\ \ibar{v}
			[$\lambda P$ $\lambda x$ $\lambda \epsilon$ $P(\epsilon)$ $\wedge$ \textsc{agent}($x$)($\epsilon$)\\ \obar{v} [agent]]
			[$\lambda \epsilon$ \textsc{visit}($\epsilon$) $\wedge$ \textsc{patient}(Adam)($\epsilon$) \\ VP [$\lambda \epsilon$ \textsc{visit}($\epsilon$) $\wedge$ \textsc{patient}(Adam)($\epsilon$)\\ \ibar{V}
				[$\lambda x$ $\lambda \epsilon$ \textsc{visit}($\epsilon$) $\wedge$ \textsc{patient}($x$)($\epsilon$)\\ \obar{V}
					[$\lambda \epsilon$ \textsc{visit}($\epsilon$)\\ \obar{V} [$\surd$visit] ]
					[$\lambda P$ $\lambda x$ $\lambda \epsilon$ $P(\epsilon)$ $\wedge$ \textsc{patient}($x$)($\epsilon$)\\ \obar{v} [patient]]
				 ]
				[DP [Adam, roof] ]
			]]
		 ]
	]
	\end{forest}
\end{exe}

We can also use the rules for Vocabulary Insertion to account for optional external $\theta$-roles. Consider, by way of illustration, the verb \emph{melt}. We've seen that this verb has both a transitive and intransitive guise. This can be expressed by letting the vocabulary item \textipa{mElt} be able to spell out a \obar{V} that contains $\surd$melt, whether or not it also contains Agent.
\begin{exe}
	\raggedright
	\ex \label{ex4:melt}
	\begin{xlist}
		\ex {} [_{\text{\obar{V}}} $\surd$melt+\emph{patient}+\emph{agent} ] $\Rightarrow$ /\textipa{mElt}/
		\ex {} [_{\text{\obar{V}}} $\surd$melt+\emph{patient} ] $\Rightarrow$ /\textipa{mElt}/
	\end{xlist}
	\end{exe}
This permits either of the representations in (\ref{ex4:mspello}) to contain~\emph{melt}.
\begin{exe}	
	\ex \label{ex4:mspello}
		\small
		\begin{forest}
		[TP
			[DP [the ice, name=binder, roof]]
			[\ibar{T} 
				[\obar{T} [will] ]
				[VP [\ibar{V}
					[\obar{V} [$\surd$melt\\ $\theta$-role: \emph{patient}] ]
					[{\ }, name=trace ]
				] ]
			]
		]
		\draw[->,>=stealth', ultra thin, gray] (trace) to[out=south west,in=south,looseness=1] node[pos=0.75,xshift=-10mm]{} (binder);
		\end{forest}
	\hspace*{12pt}
	\begin{forest}
	[TP
		[DP  [Pius, name=binder, roof]]
		[\ibar{T} 
			[\obar{T} [will] ]
			[vP
				[{\ }, name=trace ]
				[\ibar{v}
					[\obar{v} [agent] ]
					[VP [\ibar{V}
						[\obar{V} [$\surd$melt\\ $\theta$-role: \emph{patient}\\ $\theta$-role: \emph{agent}] ]
						[DP [the ice, roof] ]
					] ]
				 ]
			 ]
		]
	]
	\draw[->,>=stealth', ultra thin, gray] (trace) to[out=south west,in=south,looseness=1] node[pos=0.75,xshift=-10mm]{} (binder);
	\end{forest}
\end{exe}

\section{\textsc{agree} and Case}

We have now encountered two systems in which inflectional morphology is assigned by one term to another. In addition to the cases examined in this chapter, where an inflectional head assigns feature values to verbs that they c-command, we made use of a system of Case assignment in the previous chapter to determine the position of argument DPs. These two systems do something very similar: they both determine the inflectional morphology on a term by way of the proximity that term has to an item whose presence seems to be responsible for giving a value to that inflectional morphology. In the scenarios involving Case, the values range over nominative, accusative and genitive; and the terms that seem to be responsible for fixing those values are finite (and maybe some infinitival) \obar{T} (for nominative), a determiner for genitive Case, and the external $\theta$-role assigner, v, and prepositions (for accusative). Given the similarities in what these two systems are doing, it is reasonable to hypothesize that they are the same system. We'll consider in this section how the two might be collapsed.

I'll set aside how genitive Case is assigned, and restrict attention to just nominative and accusative. Presently, the two systems are as defined in~(\ref{ex4:vsys}) and~(\ref{ex4:csys}).
\begin{exe}
\ex \label{ex4:vsys}  Verbal Inflection
	\begin{xlist}
	\ex  \textsc{agree} (provisional)\\
	Let \obar{X} and \obar{Y} be terms that dominate features, and let a feature in \obar{X} have a value and a feature in \obar{Y} be unvalued. A feature in \obar{Y} can get the value of a matching feature in \obar{X} just in case \obar{Y} \textsc{agree}s with \obar{X}. \obar{X} and \obar{Y} \textsc{agree} iff:
\begin{xlist}
	\ex \obar{X} c-commands \obar{Y}, and
	\ex there is no \obar{Z} such that \obar{X} c-commands \obar{Z} and \obar{Z}\\ c-commands~\obar{Y}.
\end{xlist}

	\ex Stray Feature Filter\\
	A valued feature must be morphologically expressed and an unvalued feature must be valued at S-structure.

	\ex If a term has $\bullet$, then the \obar{X} containing that term must be inside a \obar{Z} that \obar{X} \textsc{agree}s with at S-structure.

	\ex Feature assignments\marginnote[-15pt]{The statement that unvalued Case features are a part of argument DPs is cryptic. Our model requires that all of the features and morphemes come into a sentence through lexical insertion. For a feature to be a part of a phrase amounts to saying, then, that there is a head in that phrase whose lexical item contains that feature. What that head is, and why it is found in argument DPs, but not non-argument DPs, is unknown.}
	\begin{xlist}
	\ex An agreement feature is part of certain \obar{Agr}s and verbs. It comes with a value in~\obar{Agr}/\obar{T}.
	\ex A tense feature is part of certain \obar{T}s and verbs. It comes valued in~\obar{T}s.
	\end{xlist}
	\pagebreak
\end{xlist}
	\ex \label{ex4:csys} Case Assignment
	\begin{xlist}
	\ex Nominative Case is assigned by finite \obar{T}/\obar{Agr} to the Specifier of the phrase it heads.
	\ex Accusative Case is assigned by \obar{v} and some prepositions to positions they:
	\begin{xlist}
	\ex c-command, and
	\ex are klose to.
\end{xlist}

	\ex The Case Filter\\
	An argument DP must be in a Case-marked position in some parse in a derivation.
\end{xlist}
\end{exe}
We saw evidence that Accusative Case assignment is subject to a locality condition, and this is what the ``klose to'' clause in (\ref{ex4:csys}) refers to. We haven't yet arrived at a working definition of this locality condition, however.

To bring the inflectional feature system and the Case system together, I'll start by translating Case into the feature based system that we've expressed verbal inflection with. We'll let certain \obar{T}s, prepositions, and v come with a value for the Case feature, and we'll assume that every DP has an unvalued Case feature. We'll model Case assignment, then, as fixing values for the Case feature under \textsc{agree}. Taking this step will create the system in (\ref{ex4:newsys1})--(\ref{ex4:featass}).
\enlargethispage*{20pt}
\begin{exe}
\ex \label{ex4:newsys1}
 Let X and Y be terms that dominate features, and let a feature in X have a value and a feature in Y be unvalued. A feature in Y can get the value of a matching feature on X just in case Y \textsc{agree}s with X. X and Y \textsc{agree} iff:
\begin{xlist}
	\ex X c-commands Y, and
	\ex X is klose to Y.
\end{xlist}

	\ex X is klose to Y iff:
	\begin{xlist}
	\ex For Y a head, there is no \obar{Z} such that X c-commands \obar{Z} and \obar{Z}\\ c-commands~Y.
	\ex For Y a phrase, <something yet to be determined>
\end{xlist}

	\ex Stray Feature Filter\\
	Every valued feature must be morphologically expressed and every unvalued feature must get a value at S-structure.

	\ex If a lexical item has $\bullet$, then it must be pronounced adjoined to the heads whose features it is in an \textsc{agree} relation with at S-structure.

	\ex Feature assignments \label{ex4:featass}
	\begin{xlist}
	\ex An agreement feature is part of verbs and certain \obar{T}/\obar{Agr}s. It comes valued in~\obar{T}/\obar{Agr}.
	\ex A tense feature is part of certain \obar{T}s and verbs. It comes valued on some~\obar{T}s.
	\ex A unvalued Case feature is part of every argument DP.
	\ex \label{ex4:newsys1end} A valued Case feature is part of v and certain prepositions.
	\end{xlist}
\end{exe}
The following changes have been made. First, the terms that can enter into an \textsc{agree} relation have been widened to include phrases as well as heads. The locality condition imposed on the term that has the value to assign and the term whose feature gets assigned that value has been relativized so that it discriminates between the cases in which a head or phrase are involved. If a feature is in a head, then the Head Movement Constraint applies; if a features is part of a phrase, then the definition there is a different locality condition invoked, one whose particulars we have yet to figure out.

Next, the Case Filter has been stated as a condition on features. The requirement is now that unvalued features must get a value before the term that bears those features can be pronounced. This condition has been added to the Stray Feature filter, which requires that values get assigned to terms that can express them. The combined condition is renamed the Feature Valuation Requirement. Finally, valued Case features are stated to be part of our former roster of Case assigners, and unvalued Case features are stated to be part of the DPs that the Case filter holds of.

This reformulation leaves out the condition under which nominative Case assignment was formerly thought to apply. That condition had nominative Case assigned to the Specifier position of TP by certain agreement morphemes that reside in \obar{T}. It seems that it is the agreement that determines nominative Case because it appears that whether nominative Case is available depends on the kind of clause involved, and the ``kinds'' can be distinguished by their heads. The assumption that nominative Case is assigned to the Specifier of TPs, is assumed because it allows for an explanation of why argument DPs seem to be forced to move into that position when no other Case is available to them. Together, these two assumptions require that nominative Case be assigned by \obar{T} to their Specifier positions. To fold nominative Case assignment into this system requires we abandon one or the other of these assumptions.

In fact, both assumptions are suspect. As you might have noticed, the fact that Argument DPs must move to Specifier of TP is often redundantly captured by the EPP. We might conjecture, then, that nominative Case is assigned to some other position, and let Argument movement to Specifier of TP be driven by the need to satisfy the EPP. On the other hand, we've just seen that where nominative Case is available seems to depend on the nature of \obar{T}. There is a way of keeping both these ideas.

Consider what would be expected if the features in finite \obar{T} were to assign Case under the conditions that prepositions and vs do. In examples in which no NegP is present, we would expect the Specifier of vP to be the position that nominative Case is assigned to. A sentence like (\ref{ex4:nomsamp1})  would have the D-structure indicated.
\begin{exe}
	%
%
\marginnote[0pt]{Only the Case features are indicated in (\ref{ex4:nomsamp1}).}%
%
\ex \label{ex4:nomsamp1} Wayne ran. \hspace*{12pt}
\small
\begin{forest}
[TP[\ibar{T}
	[\obar{T} [past\\ case: \emph{nom}]]
	[vP
		[case: \\ DP [Wayne, roof]]
		[\ibar{v} 
			[\obar{v} [agent] ]
			[VP [run, roof] ]
		 ]
	 ]
]]
\end{forest}
\end{exe}
The DP bears an unvalued Case feature, a feature that must be valued before the Stray Feature Filter will be satisfied. We can invoke \textsc{agree} to remedy this.
\begin{exe}
\ex \label{ex4:nomsamp2}
\small
\begin{forest}
[\ibar{T}
	[\obar{T} [past\\ case: \emph{nom}, name=probe] ]
	[vP 
		[case: \emph{nom}\\DP, name=goal [Wayne, roof] ]
		[\ibar{v}
			[\obar{v} [agent] ]
			[VP [run, roof] ]
		 ]
	 ]
]
\draw[|-|, dotted] (goal) to[out=south west,in=south west] node[pos=0.5,xshift=-5mm]{\emph{agree}} (probe);
\end{forest}
\end{exe}
This representation satisfies the Stray Feature Filter (the former Case Filter), but it does not yet satisfy the EPP, which requires further movement of \emph{Wayne} to Specifier of TP. The result is (\ref{ex4:nomsamp3}) \vpageref[below]{ex4:nomsamp3}, which corresponds to the grammatical S-structure outcome for this sentence.
\begin{exe}
\ex \label{ex4:nomsamp3}
\small
\begin{forest}
[TP
	[case: \emph{nom}\\ DP [Wayne, name=binder, roof] ]
	[\ibar{T} 
		[\obar{T} [past\\ case: \emph{nom}] ]
		[vP
			[{\ }, name=trace ]
			[\ibar{v} 
				[\obar{v} [agent] ]
				[VP [run, roof] ]
			 ]
		 ]
	]
]
\draw[->,>=stealth', ultra thin, gray] (trace) to[out=-110,in=-70,looseness=1.4] node[pos=0.75,xshift=-10mm]{} (binder);
\end{forest}
\end{exe}

Under this interpretation of how nominative Case is assigned, Argument movement must be allowed to move arguments from positions in which they get Case marked. There are cases which suggest Argument movement should be disallowed from moving from Case marked positions; (\ref{ex4:nocase}) is one such example.
\begin{exe}
\ex[*]{\label{ex4:nocase} \linkto Wayne \under seemed to \from{\ \ } that pigs fly.}
\end{exe}
Adopting this view of nominative Case assignment therefore requires resolving this conflict.

One proposal is to allow movement from Case marked positions only when the argument moves to the Specifier of the head responsible for assigning that Case. Indeed, one proposal is to reduce the EPP to the $\bullet$ property. Imagine that the nominative feature associated with \obar{T} gives \obar{T} the $\bullet$ property. The $\bullet$ property would then require that \obar{T} be put together with the DP whose Case feature it values. The Likes-Attracts-Likes Constraint, however, prevents phrases from adjoining to heads, so it will prevent the DP from merging directly with \obar{T}. If we loosen the statement of the requirement that the $\bullet$ property invokes so that the Likes-Attracts-Likes Constraint is left to determine how the moving term adjoins to the head with the $\bullet$ property, then we could see the EPP as just another instance of the $\bullet$ property. The requirement invoked by the $\bullet$ property could be expressed as~(\ref{ex4:newstar}).
\begin{exe}
\ex \label{ex4:newstar} If a term, $\alpha$, has $\bullet$, then a projection of the \obar{X} that contains $\alpha$ must be pronounced as a sister to a projection of a \obar{Z} that \obar{X} \textsc{agree}s with.
\end{exe}
% //todo Fix this. The stray feature filter is also the reverse Case filter, and that will prevent this derivation.

This will explain why Argument movement is blocked in examples like (\ref{ex4:nocase}), but allowed in (\ref{ex4:nomsamp3}). Consider~(\ref{ex4:nocased}): the D-structure representation~(\ref{ex4:nocase}) gets.
\begin{exe}
\ex \label{ex4:nocased}
\small
\begin{forest}
[TP
	[\obar{T} [past\\ $\bullet$case: \emph{nom}]]
	[VP [\ibar{V} 
		[\ibar{V} 
			[\obar{V} [$\surd$seem] ]
			[CP [that pigs fly, roof] ]
		 ]
		[PP [\ibar{P}
			[\obar{P} [to\\ case: \emph{acc}] ]
			[case: \\DP [Wayne, roof] ]
		] ]
	]
]]
\end{forest}
\end{exe}
The Case feature associated with \emph{Wayne} can get a value only from the preposition \emph{to}, as this is the only term with a valued Case feature that is klose enough to it. This means that \obar{T} cannot assign its Case feature value to this, or any other, DP. Its valued nominative Case feature will therefore not get morphologically expressed, and this is a violation of the Stray Feature Filter. What makes nominative Case different from the others, then, is the EPP, which independently requires that Specifier of TP be occupied. A working%
%
\marginnote{Indeed, the common way of referring to $\bullet$ is as an ``EPP feature,'' though technically it's an attribute of a feature rather than a feature.}%
%
\ hypothesis is that the EPP is just $\bullet$, whose consequences in this domain is movement of the Nominative Case marked DP rather than movement of a head. Let us add, then,~(\ref{ex4:nomvl}).
\begin{exe}
\raggedright
	\ex \label{ex4:nomvl} A valued $\bullet$ Nominative feature resides within \obar{T}.
\end{exe}
Presently the requirement that $\bullet$ imposes is written so that it applies correctly to cases in which one head moves to another. We need to widen this so that its requirements can be satisfied by moving a phrase. This is what (\ref{ex4:bulreq}) does.
\begin{exe}
\raggedright
	\ex \label{ex4:bulreq} If a term, $\alpha$, has $\bullet$, then a projection of the \obar{X} containing $\alpha$ must be pronounced as a sister to a projection of a \obar{Z} that \obar{X} \textsc{agrees} with.
\end{exe}
To see how (\ref{ex4:bulreq}) works, consider the two scenarios in~(\ref{ex4:2scenes}).
\begin{exe}
\raggedright
	\ex \label{ex4:2scenes}
	\small
	\begin{forest}
	[\ibar{X}
		[\obar{X}, name=probe [$\bullet$ feature: \textit{value}]]
		[YP [\ibar{Y}
			[\obar{Y}, name=goal [feature: \textit{value}]]
			[ZP ]
		] ]
	]
	\draw[|-|, dotted] (goal) to[out=south west,in=south west] node[pos=0.5,xshift=-5mm]{\emph{agree}} (probe);
	\end{forest}
	\hspace*{8pt}
	\begin{forest}
	[\ibar{X}
		[\obar{X}, name=probe [$\bullet$ feature: \textit{value}]]
		[YP [\ibar{Y}
			[\obar{Y} ]
			[ZP
				[WP\\ feature: \textit{value}, name=goal ]
				[\ibar{Z} [\obar{Z}] ]
			 ]
		] ]
	]
	\draw[|-|, dotted] (goal) to[out=south west,in=south west] node[pos=0.5,xshift=-5mm]{\emph{agree}} (probe);
	\end{forest}
\end{exe}
For the requirement in (\ref{ex4:bulreq}) to be satisfied, movement must create a structure in which \obar{Y} or WP is a sister to a projection of \obar{Y}. The Likes-Attracts-Likes constraint allows that to happen by adjoining \obar{Y} to \obar{X} with Head Movement. But the Likes-Attracts-Likes constraint doesn't allow that to happen by adjoining WP to \obar{Y}. Instead, the Likes-Attracts-Likes compliant way of satisfying (\ref{ex4:bulreq}) is for WP to move to Specifier of YP. The $\bullet$, then, requires the structures in (\ref{ex4:2scenes}) to map onto the representations in~(\ref{ex4:2scenesm}).
\begin{exe}
\raggedright
	\ex \label{ex4:2scenesm}
	\footnotesize
	\begin{forest}
	[\ibar{X}
		[\obar{X}
			[\obar{Y} [feature: \textit{value}, name=binder]]
			[\obar{X} [$\bullet$ feature: \textit{value}]]
		 ]
		[YP [\ibar{Y}
			[{\ }, name=trace ]
			[ZP ]
		] ]
	]
	\draw[->,>=stealth', thin, gray] (trace) to[out=south west,in=south east,looseness=1] node[pos=0.75,xshift=-10mm]{} (binder);
	\end{forest}
	\hspace*{4pt}
	\begin{forest}
	[XP
		[WP\\ feature: \textit{value}, name=binder ]
		[\ibar{X}
			[\obar{X} [$\bullet$ feature: \textit{value}]]
			[YP [\ibar{Y}
				[\obar{Y} ]
				[ZP
					[{\ }, name=trace ]
					[\ibar{Z} [\obar{Z}] ]
				 ]
			] ]
		 ]
	]
	\draw[->,>=stealth', thin, gray] (trace) to[out=south west,in=south,looseness=1] node[pos=0.75,xshift=-10mm]{} (binder);
	\end{forest}
\end{exe}

Let's consider how this will apply to a situation in which nominative Case is applied. The sentence will have a representation like (\ref{ex4:wayne1}) upon agreement.
\begin{exe}
	\raggedright
	\ex \label{ex4:wayne1}
	\small
	\begin{forest}
	[\ibar{T}
		[\obar{T}, name=probe [past\\ $\bullet$case: \emph{nom}] ]
		[vP 
			[case: \emph{nom}\\ DP, name=goal [Wayne, roof] ]
			[\ibar{v}
				[\obar{v} [agent] ]
				[VP [run, roof] ]
			 ]
		 ]
	]
	\draw[|-|, dotted] (goal) to[out=south west,in=south west,looseness=3] node[pos=0.5,xshift=-5mm]{\emph{agree}} (probe);
	\end{forest}
\end{exe}
And though this isn't a licit S-structure, since the requirements attendant with the $\bullet$ aren't met, it can be made so by Argument Movement.
\begin{exe}
	\raggedright
	\ex \label{ex4:wayne2}
	\small
	\begin{forest}
	[TP
		[case: \emph{nom}\\DP [Wayne, name=binder, roof] ]
		[\ibar{T}
			[\obar{T} [past\\ $\bullet$case: \emph{nom}] ]
			[vP
				[{\ }, name=trace ]
				[\ibar{v}
					[\obar{v} [agent] ]
					[VP [run, roof] ]
				 ]
			 ]
		 ]
	]
	\draw[->,>=stealth', ultra thin, gray] (trace) to[out=-110,in=-80,looseness=1.6] node[pos=0.75,xshift=-10mm]{} (binder);
	\end{forest}
\end{exe}

One problem in this method of collapsing the EPP with the $\bullet$ attribute comes when examples like~(\ref{ex4:cpmove}) are considered.
\begin{exe}
	\ex \label{ex4:cpmove} \linkto [\CP that mice squeak] \under seems [\AP \from{\ \ } obvious].
\end{exe}
In~(\ref{ex4:cpmove}), the EPP is satisfied by letting Argument movement move the CP \emph{that mice squeak} from its $\theta$-marked position within the AP headed by \emph{obvious} into the root AgrP. For this example to fit the proposal, it would require that the finite \obar{Agr} pass the Nominative value to the Case feature on this CP. That is, we would have to let this CP be assigned Nominative Case. In our discussion of Stowell's proposals about the order of verbal complements, we saw that it doesn't seem that CPs are subject to the Case filter, despite the evidence adduced on behalf of this by Stowell. We'll return to this problem in~\ref{sec:subinv}.

\subsection{Low Subjects: English}

This reworking of the conditions under which nominative Case assignment happens claims that nominative Case is assigned to a position lower than where the DP bearing that Case surfaces. There are examples which can be construed as evidence for this claim. In these examples, the nominative Case marked DP surfaces in a position c-commanded by the Case-assigning \obar{Agr}, and the EPP is satisfied with \emph{there}. One of these examples is~(\ref{ex4:there1}).
\begin{exe}
\ex \label{ex4:there1} There should be a solution in this chapter.
\end{exe}
It doesn't seem that \emph{there} in this example has the same meaning that the locative \emph{there} has in~(\ref{ex4:thereloc}).
\begin{exe}
\ex \label{ex4:thereloc} A solution stands there.
\end{exe}
Indeed, the distribution of the \emph{there} that is found in (\ref{ex4:there1}) is what is expected if it is an expletive. Just as in the scenarios in which the expletive \emph{it} appears, we might conjecture that the EPP (or $\bullet$ property) is satisfied by virtue of \emph{there} standing in Specifier of AgrP. We know from other examples that the copular verb \emph{be} does not assign accusative Case, so we can also conjecture that the Case which \emph{a solution} bears is nominative. Under these assumptions, an S-structure parse for (\ref{ex4:there1}) might be~(\ref{ex4:there1d}).
\begin{exe}
\ex \label{ex4:there1d}
\small
\begin{forest}
[TP
	[DP [there, roof]]
	[\ibar{T} 
		[\obar{T} [should\\ $\bullet$case: \emph{nom}] ]
		[VP [\ibar{V}
			[\obar{V} [be] ]
			[PP
				[case: \emph{nom}\\DP [a solution, roof] ]
				[\ibar{P}
					[\obar{P} [in] ]
					[DP [this chapter, roof] ]
				 ]
			 ]
		] ]
	]
]
\end{forest}
\end{exe}

If this is the right interpretation of these examples, then it confirms the hypothesis that nominative Case assignment occurs under c-command, just as accusative Case assignment does. But it presents a host of problems for the particular proposals we have arrived at. If \emph{there} satisfies the EPP in this example, then collapsing the EPP and the $\bullet$ property becomes difficult. Further, the distance between \obar{T} and the DP it Case marks in (\ref{ex4:there1d}) does not fit to our definition of klose: there is clearly more than one phrase that dominates \emph{a solution} but not \obar{Agr}. To fit this analysis to our system would require a different definition of klose, then.

Moreover, examples of this type are governed by a set of constraints that remain little understood. The semantic type of the nominative Case marked DP matters, for instance. That argument must be an indefinite of a certain type, as the contrasts in (\ref{ex4:defthere}) illustrate.
\pagebreak
\begin{exe}
\ex \label{ex4:defthere}
\begin{xlist}
	\ex[]{There shouldn't be any solutions in this chapter.}
	\ex[]{There should be solutions in this chapter.}
	\ex[]{There should be no solutions in this chapter.}
	\ex[]{There should be many solutions in this chapter.}
	\ex[*]{There should be every solution in this chapter.}
	\ex[*]{There should be the solution in this chapter.}
	\ex[*]{There should be most solutions in this chapter.}
	\ex[*]{There should be all solutions in this chapter.}
	\ex[*]{There should be both solutions in this chapter.}
	\ex[*]{There should be each solution in this chapter.}
	\ex[*]{There should be your solution in this chapter.}
\end{xlist}
\end{exe}
And the kind of predicate involved also matters. While PPs support this construction, APs and DPs do not.
\begin{exe}
\ex \label{ex4:therepred}
\begin{xlist}
	\ex[*]{There is a man a fink.}
	\sn[]{\emph{compare:} A man is a fink.}
	\ex[*]{There is a man happy.}
	\sn[]{\emph{compare:} A man is happy.}
\end{xlist}
\end{exe}
VPs support the construction, but for many verbs, only when they are embedded under \emph{be}.
\begin{exe}
\ex \label{ex4:therepred2}
\begin{xlist}
	\ex[]{There was a woman given a book.}
	\ex[]{There was a woman drinking milk.}
	\ex[*]{There has a woman given me a book.}
	\ex[*]{There has a woman drunk milk.}
	\ex[*]{There a woman drank milk.}
\end{xlist}
\end{exe}
The exception are some (but not all) unaccusative verbs, as in~(\ref{ex4:therepred3}).
\begin{exe}
\ex \label{ex4:therepred3}
\begin{xlist}
	\ex[]{There arrived a woman.}
	\ex[]{There appeared an elf.}
\end{xlist}
\end{exe}
And in these cases, as can be seen, the nominative Case marked DP remains in its \mbox{$\theta$-marked} object position.\marginnote[-1.5in]{For some discussion of the constraints on the semantic type of the DP, see \citet{Milsark:1974}, \citet{Safir:1985,Safir:1987b}, \citet{Heim:1987} and \citet{Diesing:1992b}. For discussion of the constraints on the predicates, see \citet{Levin:1993}. There is a large, and heterogeneous, literature that addresses the problem of assimilating the way \emph{there} satisfies the EPP to a more general theory of the EPP. Among many others, see \citet{Chomsky:1995}, \citet{Lasnik:1992b, Lasnik:1995}, \citet{Moro:1997}, \citet{Chomsky:1998,Chomsky:2000,Chomsky:2001a}, \citet{Boskovic:2002, Boskovic:2002c}, \citet{Groat:1999,Groat:1995}, \citet{Hazout:2004a}, and \citet{Epstien:2005}. A good introduction to the language variation in expletives is \citet{Vikner:1995}.}

A better understanding of how to assimilate this construction to the general theory of Case assignment will probably appear once the source of these constraints is found.%\enlargethispage*{4pt}

\subsection{Low Subjects: Subject Inversion Languages} \label{sec:subinv}

There are languages in which Nominative Case-marked DPs surface in \mbox{S-structure} positions that, like those we've just seen in English, are c-commanded by \obar{T}. In these languages, interestingly, there is no expletive like \emph{there} involved. Spanish, Italian and Greek are such languages. Examples from Spanish and Greek are in~(\ref{ex4:noepp1}).%
%
\marginnote[-45pt]{My description of these languages will assume that, as in English apparently, the agreement and tense features are bundled into the \obar{T} position at D-structure. If, as in French, they are inserted into different positions at D-structure, then I should call \obar{Agr} the node that assigns Nominative Case.}%
%
\pagebreak
\begin{exe}
\ex \label{ex4:noepp1}
\begin{xlist}
	\ex \gll Leyo Juan el libro. \\ read John the book\\ `John read the book.'
	\ex \gll Pandreftike {o Petros} {tin Ilektra}.  \\ married Peter Ilektra\\ `Peter married Electra.'
\end{xlist}
\hspace*{\fill} \citep[(3c,d):~492]{Alexiadou:1998a}\hspace*{.5in}
\end{exe}
We might venture a parse for (\ref{ex4:noepp1}b) like that in~(\ref{ex4:noepp1t}).
\begin{exe}
\ex \label{ex4:noepp1t}
\small
\begin{forest}
[TP[\ibar{T}
	[\obar{T} 
		[\obar{V} [pandreftike, name=binder] ]
		[\obar{T} [case: nom] ]
	]
	[VP 
		[case: \emph{nom}\\ DP [o Petros, roof] ]
		[\ibar{V}
			[{\ }, name=trace ]
			[DP [tin Ilektra, roof] ]
		 ]
	]
]]
\draw[->,>=stealth', ultra thin, gray] (trace) to[out=-110,in=-80,looseness=1.5] node[pos=0.75,xshift=-10mm]{} (binder);
\end{forest}%
%
\marginnote[0pt]{I've suppressed ``\obar{v}'' in this parse.}%
%
\end{exe}
These examples don't pose the same problem for our definition of klose that the ones in English did, but in these cases there appears to be a problem for expressing the EPP with the $\bullet$ property. That no expletive is required in the Specifier of TP in these examples suggests that these are languages that have no EPP. Maybe that means there is no $\bullet$ on their nominative case.

Note that in (\ref{ex4:noepp1t}), the main verb has moved to \obar{T}, thereby putting it in a position that the linearization principles of Greek will cause it to precede the nominative Case marked DP. This is a very common trait of languages that allow their nominative DPs to remain low; as in Greek those DPs follow the verb (when they are not moved by other processes like wh-movement). In particular, the implications in (\ref{ex4:sicor}) seem to hold.
\begin{exe}
\ex \label{ex4:sicor} If a language has (\ref{ex4:sicorex})--(\ref{ex4:sicorexl}), then its nominative DPs can follow the finite verb.
\begin{xlist}
	\ex \label{ex4:sicorex} Allows nominative DPs to surface c-commanded by \obar{T}.
	\ex \label{ex4:sicorexl}Linearizes VPs and the heads that verbs move to so that they precede their complements
	\ex Has no expletive
\end{xlist}
\end{exe}
These languages are said to display ``subject inversion.'' \citet{Perlmutter:1971} and \citet{Rizzi:1982} proposed that subject inversion is always found in languages that don't have the EPP, and are verb initial. One influential attempt to explain this correlation is in \citet{Alexiadou:1998a,Alexiadou:1999a}. They argue that, in fact, the EPP does hold in these languages, but that it is satisfied by moving the verb instead of the nominative Case-marked~DP.

For that to make sense under an interpretation of the EPP that collapses it with the $\bullet$ property, we cannot express the EPP as a property of Case. Verbs do not carry morphology that expresses Case features, and so they will not form an \textsc{agree} relation with the Case feature associated with \obar{T}. If the EPP is the $\bullet$ attribute associated with \obar{T}'s Case feature, then our present formulation of what the $\bullet$ attribute requires would not be satisfied by moving the verb to \obar{T}. For the proposal in \citet{Alexiadou:1998a} to make sense, we will have to see the $\bullet$ attribute as being associated with some feature that both the verb and the nominative Case marked DP express. That feature, as it turns out, is agreement. We have already built into our system that verbs express the values of agreement features. But, as the label ``agreement'' hints at, those features are determined by the nominative Case marked DP. As the table in (\ref{ex4:numbper}) indicates, the number and person of the nominative Case marked DP in English controls the morphological form of the finite verb.
\begin{exe}
\ex \label{ex4:numbper} For present tense \emph{be}\\[4pt]
\begin{tabular}{lll}
& singular & plural\\
\nth{1} person & \emph{am} & \emph{are}\\
\nth{2} person & \emph{are} & \emph{are}\\
\nth{3} person & \emph{is} & \emph{are}\\
\end{tabular}
\end{exe}
Indeed, all other verbs merely distinguish the \nth{3} singular forms from all others. But the existence of the paradigm in (\ref{ex4:numbper}) indicates that both person and number features are present in the agreement morphology of English. We should therefore decompose the agreement feature into number and person features and let \textsc{agree} fix the values of these features in a way that involves not just \obar{T} and the verb, but the nominative Case-marked DP as well. We'll take up that task in the next section.

The proposal in \citet{Alexiadou:1998a}, then, requires that it be the agreement feature that has the $\bullet$ property responsible for manufacturing EPP effects. Schematically, then, the relevant verb initial languages will have a D-structure representation like that given in~(\ref{ex4:ualex}).
\begin{exe}
\ex \label{ex4:ualex}
\small
\begin{forest}
[TP
	[\obar{T}\\ $\bullet$person: \emph{v}\\ $\bullet$number: \emph{w}\\ tense: \emph{y}\\ case: \emph{nom} ]
	[VP
		[DP\\ case: \emph{nom}\\person: \emph{v}\\ number: \emph{w} ]
		[\ibar{V}
			[\obar{V}\\ person: \emph{v}\\ number:\emph{w}\\ tense: \emph{y} ]
			[XP]
		 ]
	 ]
]
\end{forest}
\end{exe}
In one class of language, the requirements of the $\bullet$ attribute are satisfied by Argument moving the DP into Specifier of TP. This produces the typical word-order in English sentences, illustrated by~(\ref{ex4:alexeng}).
\begin{exe}
\ex \label{ex4:alexeng}
\small
\begin{forest}
[TP
	[DP\\ case: \emph{nom}\\ person: \emph{v}\\ number: \emph{w}, name=binder]
	[\ibar{T}
		[\obar{T}\\ $\bullet$person: \emph{v}\\ $\bullet$number: \emph{w}\\ case: \emph{nom}\\ tense: \emph{y} ]
		[VP
			[{\ }, name=trace ]
			[\ibar{V}
				[\obar{V}\\ person: \emph{v}\\ number: \emph{w}\\ tense: \emph{y} ]
				[XP ]
			 ]
		 ]
	 ]
]
\draw[->,>=stealth', ultra thin, gray] (trace) to[out=south west,in=south,looseness=1.5] node[pos=0.75,xshift=-10mm]{} (binder);
\end{forest}
\end{exe}
In the other class of language, the verb's movement satisfies the $\bullet$ property on the person and number features, resulting in an S-structure like~(\ref{ex4:alexnon}).
\begin{exe}
\ex \label{ex4:alexnon}
\small
\begin{forest}
[TP[\ibar{T}
	[\obar{T} 
		[\obar{V}\\ person: \emph{v}\\ number: \emph{w}\\ tense: \emph{y}, name=binder ]
		[\obar{T}\\ $\bullet$person: \emph{v}\\ $\bullet$number: \emph{w}\\ tense: \emph{y} ]
	 ]
	[VP
		[DP\\ case:\\ person: \emph{v}\\ number: \emph{w} ]
		[\ibar{V}
			[{\ }, name=trace ]
			[XP ]
		 ]
	 ]
]]
\draw[->,>=stealth', ultra thin, gray] (trace) to[out=-110,in=-80,looseness=1.2] node[pos=0.75,xshift=-10mm]{} (binder);
\end{forest}
\end{exe}
As can be seen, the correlative properties of high verb and low nominative DP are captured.

The cases in English involving expletive \emph{there} seem to form a special case. Like the languages with the syntax in (\ref{ex4:alexnon}), the movement of the nominative Case marked DP is not required by the $\bullet$ property. But unlike the languages in which this is generally possible, the movement of verb does not satisfy the $\bullet$ property; instead, the presence of \emph{there} does. We will not examine how \emph{there} has this ability, nor will we try to unpack the causes for the constraints that govern this word-order in English. With these problems bracketed, however, we can see the \emph{there}-construction in English as a special instance of the syntax that gives subject inversion elsewhere.

If the EPP arises because of $\bullet$ being on agreement features, rather than Case features, we might have a solution to the problem raised at the end of 4.7. That problem is how the EPP is satisfied by the movement of a CP, as well as a DP. If $\bullet$ is on Nominative Case, then we should expect its requirements to be satisfied only by moving things that \textsc{agree} in Case with the Nominative feature. We have evidence that CPs are not subject to the Case filter, and so it's unlikely that they could host an unvalued Nominative feature. But it's conceivable that CPs have the person and number features that English agreement morphology houses. This would permit derivations like that in~(\ref{ex4:cpraise}).
\begin{exe}
\raggedright
	\ex \label{ex4:cpraise}
	\small
	\begin{forest}
	[TP [\ibar{T}
		[\obar{T}, name=probe [$\bullet$ person: nth{3}\\ $\bullet$ number: sing\\ tense: pres]]
		[VP [\ibar{V}
			[\obar{V} [seem]]
			[AP
				[CP\\ person: \nth{3}\\ number: sing, name=goal [that mice squeak, roof]]
				[\ibar{A} [obvious, roof]]
			 ]
		] ]
	]]
	\draw[|-|, dotted] (goal) to[out=south west,in=south west] node[pos=0.5,xshift=-5mm]{\emph{agree}} (probe);
	\end{forest}
	\begin{forest}
	[TP 
	[CP\\ person: \nth{3}\\ number: sing, name=goal [that mice squeak, roof, name=binder]]
	[\ibar{T}
		[\obar{T}, name=probe [$\bullet$ person: nth{3}\\ $\bullet$ number: sing\\ tense: pres]]
		[VP [\ibar{V}
			[\obar{V} [seem]]
			[AP
				[{\ }, name=trace]
				[\ibar{A} [obvious, roof]]
			 ]
		] ]
	]]
	\draw[->,>=stealth', thin, gray] (trace) to[out=south west,in=south,looseness=1] node[pos=0.75,xshift=-10mm]{} (binder);
	\end{forest}
\end{exe}

A fact that might be related to this is peculiar restriction on A-movement of CPs discussed in section 3.7. We discovered there evidence that CPs can A-move only if they stand in positions where argument DPs can. Thus, for example, though (\ref{ex4:cppara1}) are grammatical instances of the passive, (\ref{ex2:cppara2}) are ungrammatical.
\begin{exe}
\raggedright
	\ex \label{ex4:cppara1}
	\begin{xlist}
		\ex It was hoped that this will be the outcome.
		\ex It was insisted that this is true.
		\ex It was felt that this is likely.
	\end{xlist}
	\pagebreak
	\ex \label{ex2:cppara2}
	\begin{xlist}
		\ex[*]{That this will be the outcome was hoped.} 
		\ex[*]{That this is true was insisted.}
		\ex[*]{That this is likely was felt.}
	\end{xlist}
\end{exe}
The verbs \textit{hope}, \textit{insist}, and \textit{feel} do not c-select DP objects, as can be seen from the ungrammaticality of~(\ref{ex4:cppara3}).
\begin{exe}
\raggedright
	\ex \label{ex4:cppara3}
	\begin{xlist}
		\ex[*]{This outcome was hoped.}
		\ex[*]{This claim was insisted.}
		\ex[*]{This emotion was felt.}
	\end{xlist}
\end{exe}
Imagine that c-selection is not for DPs, but is instead for phrases that have person and number features associated with them. Because all DPs are the only phrases that can, and must, have these features, this will come close to having the same empirical footprint as c-selecting for DP. Finally, imagine that person number features can be associated with CPs, but that CPs aren't required to have these features. Framing c-selection in terms of person/number features means that the verbs \textit{hope}, \textit{insiste} and \textit{feel} don't allow their complements to be phrases with person/number features. The prevents them from combining with DPs (which must have person/number features) and with CPs unless those CPs don't have person/number features. But if a CP doesn't have person/number features, then it won't be able to stand in an \textsc{agree} relation with those unvalued features in \obar{T}. The only way to satisfy a $\bullet$ associated with these unvalued person/number features will be with an expletive. That is why (\ref{ex2:cppara1}) is grammatical, but (\ref{ex2:cppara2}) isn't. 

Whether the \textsc{agree} relation between \obar{T} and an argument DP or CP involves $\bullet$'d person/number features, or a $\bullet$'d Case features, we'll have to abandon our present formulation of ``klose,'' the locality condition on feature valuation when phrases are involved. We designed klose so that it would allow accusative Case assignment in (\ref{ex4:ecmklose}), but not in~(\ref{ex4:necmklose}).
\begin{exe}
\ex \label{ex4:ecmklose}
\small
\begin{forest}
[\ibar{v}
	[\obar{v}\\ case: \emph{acc}, name=probe ]
	[VP [\ibar{V}
		[\obar{V} [believe] ]
		[AP
			[DP\\ case:, name=goal ]
			[\ibar{A} [obvious, roof] ]
		 ]
	] ]
]
\draw[|-|, dotted] (goal) to[out=south west,in=south west,looseness=1.3] node[pos=0.5,xshift=-5mm]{\emph{agree}} (probe);
\end{forest}
\end{exe}
\begin{exe}
\ex \label{ex4:necmklose}
\small
\begin{forest}
[*\ibar{v}
	[\obar{v}\\ case: \emph{acc}, name=probe ]
	[VP [\ibar{V}
		[\obar{V} [believe] ]
		[AP 
			[TP 
				[DP\\ case:, name=goal ]
				[\ibar{T} 
					[\obar{T} [to] ]
					[VP [be happy, roof] ]
				 ]
			 ]
			[\ibar{A} [obvious, roof] ]
		 ]
	] ]
]
\draw[|-|, dotted] (goal) to[out=south west,in=south west] node[pos=0.5,xshift=-5mm]{\emph{agree}} (probe);
\end{forest}
\end{exe}
We now need it to allow (\ref{ex4:longagr}) as well.
\begin{exe}
\ex \label{ex4:longagr}
\small
\begin{forest}
[TP
	[{\ } ]
	[\ibar{T}
		[\obar{T}\\case: \emph{nom}, name=probe ]
		[VP [\ibar{V}
			[\obar{V} [have] ]
			[IP [\ibar{I}
				[\obar{I} [perfect] ]
				[vP
					[DP\\case:, name=goal [dogs, roof] ]
					[\ibar{v} [bark, roof] ]
				 ]
			] ]
		] ]
	 ]
]
\draw[|-|, dotted] (goal) to[out=south west,in=south west] node[pos=0.5,xshift=-5mm]{\emph{agree}} (probe);
\end{forest}
\end{exe}
We cannot define klose in such a way that it legislates against too many phrases between the Case assigner and the DP that gets assigned Case, as we did originally.

One direction that some have taken is to make reference to the geometries of the phrase markers. In the good examples, both (\ref{ex4:longagr}) and (\ref{ex4:ecmklose}), the DP is embedded in phrases that are either complements or \obar{X}s. In the bad example, by contrast, the DP is embedded within a phrase that stands in Specifier position. A wide array of movement, and other, relations seem to be blocked when they reach into a phrase in a Specifier position. For instance, neither \emph{wh} Movement nor Complex NP Shift are possible out of such phrases:
\pagebreak
\begin{exe}
\ex \label{ex4:subjcon}
\begin{xlist}
	\ex[*]{Who have you said that a book about -- arrived.}
	\ex[*]{That she eats -- is obvious lukewarm natto rice.}
	\sn \emph{compare:}
	\sn[]{She eats fervently lukewarm nattoo rice.}
\end{xlist}
\end{exe}
It appears that there may be something quite general about phrases in Specifier positions that makes them incapable of engaging in syntactic relations outside those phrases. \citet{Ross:1967}, who charted out many such cases, calls these regions ``islands.''\marginnote[0pt]{The particular island involved here goes by the names: ``Subject Condition'' and ``Generalized Left Branch Condition.'' For two popular approaches to defining this island, see \citet{Huang:1982}, \citet{Lasnik:1984} and \citet{Chomsky:1986}, for one, and \citet{Kayne:1983} and \citet{Pesetsky:1982} for the other.} One hypothesis about the ungrammaticality of (\ref{ex4:necmklose}), then, is that this has nothing to do with proprietary constraints on \textsc{agree} but is the result of a global constraint against establishing relations into Specifiers. Let us adopt that idea. This means that we have not found any locality condition on the assignment of features to phrases. We change \textsc{agree} to~(\ref{ex4:agreerev}).
\begin{exe}
\ex \label{ex4:agreerev} \textsc{agree} (provisional)\\
 Let X and Y be terms that dominate features, and let a feature in X have a value and those in Y be unvalued. A feature in Y can get the value of a matching feature in X just in case Y \textsc{agree}s with X. X and Y \textsc{agree} iff:
\begin{xlist}
	\ex X c-commands Y, and
	\ex If X and Y are heads, there is no \obar{Z} such that X c-commands \obar{Z} and \obar{Z} c-commands~Y.
\end{xlist}
\end{exe}

Let's now turn to unpack agreement into the person and number features we've seen them consist in, and work out how \textsc{agree} can match the values of these features on \obar{T}, a verb and a DP.
% \enlargethispage*{14pt}

\subsection{Case and Agreement}

Recognizing that the person and number features make up agreement, and that they are given the same values on both the finite verb and the nominative Case-marked DP, we want to distribute these features as indicated in~(\ref{ex4:ca1}).
\begin{exe}
\ex \label{ex4:ca1}
\small
\begin{forest}
[\ibar{T}
	[\obar{T}\\ person: \emph{v}\\ number: \emph{w}\\ case: \emph{nom}\\ tense: \emph{pres}]
	[VP 
		[DP\\ person: \emph{v}\\ number: \emph{w}\\ case: ]
		[\ibar{V}
			[\obar{V}\\ person: \emph{v}\\ number: \emph{w}\\ tense: ]
		 ]
	]
]
\end{forest}
\end{exe}
Let's consider now how \textsc{agree} would cause the values for the person feature (\emph{v} in (\ref{ex4:ca1})) and the number feature (\emph{w} in (\ref{ex4:ca1})) to be fixed.

It seems natural to think of these features as coming valued on the DP involved in the relation. These features have a role in the semantic interpretation of the DPs, though they don't obviously have any semantic role on \obar{T} or the verb. For this reason, we should take the D-structure representation to be as indicated by~(\ref{ex4:ca2}), which maps onto the sentence \emph{she runs}.
\begin{exe}
\ex \label{ex4:ca2}
\small
\begin{forest}
[TP
	[\obar{T}\\ case: \emph{nom}\\ person: \\ number:\\ tense: \emph{pres} ]
	[VP 
		[DP\\ case: \\ person: \emph{\nth{3}}\\ number: \emph{sing} [she, roof] ]
		[\ibar{V}
			[\obar{V}\\ person: \\ number:\\ tense: [run] ]
		 ]
	]
]
\end{forest}
\end{exe}
Under our present definition of \textsc{agree}, values are passed from one item, $\alpha$, to another, $\beta$, only if $\alpha$ c-commands $\beta$. This means that the person and number values on DP cannot be assigned to \obar{T} in (\ref{ex4:ca2}), though it can assign those features directly to the verb. This definition of \textsc{agree}, then, is not compatible with there being person and number features on \obar{T} if those features come valued on DPs. But if we do not put person and number features on \obar{T}, we will no longer be able to give these features the $\bullet$ attribute and employ the account for the subject inversion typology that we reviewed in the previous section.

In fact, there is another reason to think that \obar{T} has person and number features associated with it. There is a close connection between the features that figure in agreement---in English: person and number---and those that are responsible for Case. Features that can enter into agreement relations are called ``$\phi$ features,'' and in addition to person and number, they also include gender and, perhaps, word class markers like those that play a role in Bantu languages and elsewhere. It is very common cross-linguistically that the DPs which bear nominative Case also agree in $\phi$ features with the verb. Moreover, the distribution of agreement and nominative Case within, and across, languages is very similar. Within English we see this in the fact that those DPs which bear (overt) nominative Case do so in the very clauses where they trigger (overt) agreement. Nominative Case marked DPs are exclusively found as ``subjects'' of finite clauses, and this is where agreement triggering DPs are found as well. This is extremely common: in those languages that have both nominative Case assignment and agreement, it is the nominative Case marked DP that enters into agreement. This correlation is also reflected in those languages that have only nominative Case assignment or agreement. Just as in English, the nominative Case, or agreement, targets certain DPs in finite clauses but not in other reduced clauses.

This relationship is found with respect to accusative Case as well. There is no agreement associated with accusative Case-marked DPs in English, but in many languages there is. This so-called ``object agreement'' has a distribution that looks very much like that which accusative Case-marked DPs have in English: it depends on the choice of verb (``transitive'' verbs trigger it and ``intransitive'' verbs don't) and it targets DPs that are roughly in the positions that accusative Case marked DPs are in English.

What these correlations suggest is that Case and agreement are always controlled by the same term, and that when both Case and agreement are involved, they are bundled together by that one term. I'll describe these correlations with~(\ref{ex4:caprinc}).%
%
\marginnote[10pt]{This claims that if a language has both Case assignment and agreement, only one of the Case assigners will be forced to trigger agreement with the DP it assigns Case to. What we (probably) want is to force every Case assigner to control agreement with the DP it assigns Case to. We want to allow languages like English, in which one Case assigner (nominative) is tied to agreement but another (accusative) isn't, while still tying together the Case assignment and agreement phenomena.}%
%
\begin{exe}
\ex \label{ex4:caprinc} If a language has an unvalued Case feature or unvalued $\phi$ features, then there will be some head in which they are all bundled.
\end{exe}
$\phi$ features might well exist in languages where there is no agreement phenomenon, in which case they would only play a role in determining the meanings of the DPs they are in. In such languages, all the $\phi$ features would be valued. The first clause in (\ref{ex4:caprinc}) therefore restricts the constraint to those languages that have Case marking or agreement. It requires in such languages that there be a head which mediates the values of those features, thereby gluing them together. In English, this head is~\obar{T}.

For these reasons, then, we should leave person and number features associated with \obar{T}. As noted above, this means that unless we decide against letting the valued versions of the person and number features come with the DP, we have a problem for our formulation of \textsc{agree}. We cannot restrict \textsc{agree} so that it only passes values down to c-commanded terms. In (\ref{ex4:ca2}), repeated \vpageref[below]{ex4:ca2a}, we need to let the values for person and number on the DP be passed up to~\obar{T}.
\begin{exe}
\exr{ex4:ca2} \label{ex4:ca2a}
\small
\begin{forest}
[TP[\ibar{T}
	[\obar{T}\\ $\bullet$person:\\$\bullet$number:\\ case: \emph{nom}\\ tense: \emph{pres} ]
	[VP 
		[DP\\ person: \emph{\nth{3}}\\ number: \emph{sing}\\ case: ]
		[\ibar{V} [\obar{V}\\ person:\\number:\\tense: [run]]]
	]
]]
\end{forest}
\end{exe}
\pagebreak
To this end, \textsc{agree} becomes~(\ref{ex4:notlastagree}).
\begin{exe}
\ex \label{ex4:notlastagree} \textsc{agree} (provisional)\\
Let X and Y be terms that dominate features, and let a features in X have a value and a feature in Y be unvalued. A feature in X or Y can get the value of the matching feature in the other just in case X and Y \textsc{agree}. X and Y \textsc{agree} iff:
\begin{xlist}
	\ex one of X or Y c-commands the other, and
	\ex If Y is a head, there is no \obar{Z} such that one of X or Y c-commands \obar{Z} and \obar{Z} c-commands the other of X and~Y.
\end{xlist}
\end{exe}
This version of \textsc{agree} still requires that the two terms in an \textsc{agree} relation also be in a c-command relation. But it doesn't require that the values be passed from the c-commander to the c-commandee. This will have the desired effect, then, of letting the values for person and number features on a DP be assigned to a \obar{T} that this DP does not c-command. This will allow \textsc{agree} to convert~(\ref{ex4:ca2}) to~(\ref{ex4:ca3}).
\begin{exe}
\ex \label{ex4:ca3}
\small
\begin{forest}
[TP
	[\obar{T}\\ $\bullet$person: \emph{\nth{3}}\\ $\bullet$number: \emph{sing}\\ case: \emph{nom}\\ tense: \emph{pres} ]
	[VP
		[DP\\ case:\\ person: \emph{\nth{3}}\\ number: \emph{sing} ]
		[\ibar{V}
			[\obar{V}\\ person:\\number:\\tense: [run]]
		 ]
	 ]
]
\end{forest}
\end{exe}

The version of \textsc{agree} in (\ref{ex4:notlastagree}) would also let the values that DP has for the person and number features be fixed on \obar{V} by putting these terms into an \textsc{agree} relation with DP. This isn't desired, however. If the DP could directly agree with verbs, and other terms, then our desire to tie agreement to Case would be lost. We would expect that the values for $\phi$ features on a verb could be valued by a DP that is not getting its Case from \obar{T}. What we see is something different, however. It is the very DP that bears nominative Case whose $\phi$ features arise on the verb. We must prevent the DP from entering into an \textsc{agree} relation with \obar{V}. Its values for person and number must come from \obar{T}. One way of ensuring this is to require that the c-commanding term in an \textsc{agree} relation be a head. This would prevent the DP in (\ref{ex4:notlastagree}) from entering into an \textsc{agree} relation with the verb that follows. We therefore change \textsc{agree} once more.
\pagebreak
\begin{exe}
\ex \label{ex4:nxt2lastagree} \textsc{agree}\\
If X c-commands Y, then the unvalued features in X or Y can get the value of the matching features in the other just in case X and Y \textsc{agree}. X and Y \textsc{agree} iff:
\begin{xlist}
	\ex X is a head, and
	\ex If Y is a head, there is no \obar{Z} such that X c-commands \obar{Z} and \obar{Z} c-commands~\obar{Y}.
\end{xlist}
\end{exe}

\subsection{A locality condition on Agreement}

Now that the definition of \textsc{agree} has matured to the point where it applies correctly to the basic cases of agreement and Case assignment, we can investigate what locality condition, if any, there is on it. We started with the observation that Case assignment is constrained by a locality condition that we've now credited to a general ban against syntactic relations into phrases in Specifier positions.

For inflection feature assignment from \obar{T} to \obar{V}, we have relied on a locality condition that captures Lisa Travis's Head Movement Constraint. That condition prevents \textsc{agree} from holding from one head across another:
\begin{exe}
	\raggedright
	\ex \label{ex4:hmf}
	If \obar{X} \textsc{agree}s with \obar{Y}, then there can be no \obar{Z} that c-commands \obar{Y} and is c-commanded by~\obar{X}.
\end{exe}
This is what is responsible for blocking \textsc{agree}ment from \obar{T} across \obar{Neg} to a main verb, as in~(\ref{ex4:neghmill}).
\begin{exe}
	\raggedright
	\ex \label{ex4:neghmill}
	\small
	\begin{forest}
	[\ibar{T}
		[\obar{T}\\ person: \emph{\nth{3}}\\ number: \emph{sing}\\ case: \emph{nom}\\ tense: \emph{pres}]
		[NegP 
			[AdvP [not, roof] ]
			[\ibar{Neg}
				[\obar{Neg} [neg]]
				[VP [\ibar{V}
					[\obar{V} [seem\\ person:\\ number:\\ tense:] ]
					[CP [that pigs fly, roof] ]
				] ]
			 ]
		]
	]
	\end{forest}
\end{exe}
Before we'd developed the system of features and \textsc{agree} which now control head movement, we also saw that there was a locality condition that prevents deeply embedded \obar{I}s from moving to \obar{C}, as in~(\ref{ex4:hmcitc}).
\begin{exe}
	\raggedright
	\ex[*]{\label{ex4:hmcitc} Been Mary should have eating nattoo.}
\end{exe}
This is a scenario in which the \obar{C} heading this sentence has a silent question morpheme in it, and it is the need for that morpheme's stem that drives movement of \obar{I} to \obar{C}. Expressing this idea in terms of features, we would now say that this \obar{C} has a question feature in it that has the $\bullet$ attribute. In a grammatical instance of \obar{I} to \obar{C} movement, this will invoke derivations like that indicated in~(\ref{ex4:itocf}).
\begin{exe}
	\raggedright
	\ex \label{ex4:itocf}
	\small
	\begin{forest}
	[CP[\ibar{C}
		[\obar{C}
			[\obar{T} [should\\ force: \emph{Q}, name=binder] ]
			[\obar{C}\\ force: $\bullet$\emph{Q} ]
		 ]
		[TP 
			[DP [she, roof]]
			[\ibar{T}
				[{\ }, name=trace ]
				[vP [run, roof] ]
			 ]
		 ]
	]]
	\draw[->,>=stealth', ultra thin, gray] (trace) to[out=-110,in=-80,looseness=1] node[pos=0.75,xshift=-10mm]{} (binder);
	\end{forest}
	= Should she run?
\end{exe}
I've called the feature-type that Q is a value of, ``force.'' Unvalued force features are foundz on inflectional heads, we might imagine, and therefore on the \obar{T}, \emph{should}, in~(\ref{ex4:itocf}). We can translate our analysis of \obar{I} to \obar{C} movement in German, and the other languages displaying verb second word order, in a similar way. The CPs heading root declarative clauses in German can be headed by a \obar{C} that contains a declarative force feature with a $\bullet$ attribute. That will produce derivations like that indicated in~(\ref{ex4:v2feat}).
\enlargethispage*{35pt}
\begin{exe}
	\raggedright
	\ex \label{ex4:v2feat}
	\small
	\begin{forest}
	[CP
		[XP ]
		[\ibar{C}
			[\obar{C} 
				[\obar{T}, name=binder
					[\obar{V} [$\surd$haben] ]
					[\obar{T}\\ force: \emph{decl} [past]  ]
				 ]
				[\obar{C}\\ $\bullet$force: \emph{decl} ]
			]
			[TP
				[DP [sie, roof] ]
				[\ibar{T}
					[VP [\ibar{V}
						[IP [das Buch gelesen, roof] ]
						[{\ }, name=tracev ]
					] ]
					[{\ }, name=trace ]
				 ]
			 ]
		 ]
	]
	\draw[->,>=stealth', ultra thin, gray] (tracev) to[out=south east,in=south,looseness=1] node[pos=0.75,xshift=10mm]{} (trace);
	\draw[->,>=stealth', ultra thin, gray] (trace) to[out=south west,in=east,looseness=1] node[pos=0.75,xshift=-10mm]{} (binder);
	\end{forest}%
	%
	\marginnote[0pt]{I have only given the relevant features on \obar{T} and \obar{V} here. I've suppressed tense and agreement features. XP represents some phrase which has been topicalized.}%
	%
	\gll XP hat sie das Buch gelesen. \\ XP has she the book read \\ `She has read the book.'
\end{exe}
In these context in which a feature in \obar{C} \textsc{agree}s with some inflectional head, we want to make sure that it correctly targets the highest of those inflectional heads. If it could target lower ones, then we should expect to find the ungrammatical (\ref{ex4:hmcitc}). The D-structure for (\ref{ex4:hmcitc}) would be~(\ref{ex4:hmcitcd}).
\begin{exe}
	\raggedright
	\ex \label{ex4:hmcitcd}
	\small
	\begin{forest}
	[CP[\ibar{C}
		[\obar{C}\\ force: $\bullet$\emph{Q}]
		[TP 
			[DP [she, roof]]
			[\ibar{T}
				[\obar{T} [should] ]
				[VP [\ibar{V}
					[\obar{V} [have] ]
					[IP [\ibar{I}
						[\obar{I} [perfect\\ force:] ]
						[VP [\ibar{V}
							[\obar{V} [be] ]
							[IP [eating nattoo, roof] ]
						] ]
					] ]
				] ]
			 ]
		]
	]]
	\end{forest}%
	%
	\marginnote[0pt]{Note that there is no unvalued force feature in \obar{T}. We must allow the morphemes that are lexically inserted into \obar{T} to come without force features, for otherwise it wouldn't be possible to use, say, \emph{should} in anything but a question.}%
	%
\end{exe}
If \textsc{agree} could allow the force feature in \obar{C} to value the force feature in \obar{I}, then the $\bullet$ attribute on Q will force \emph{be} to move to \emph{en}, and the resulting inflected verb to move to \obar{C}, creating the ungrammatical (\ref{ex4:hmcitc}). This can be prevented if \textsc{agree} cannot hold from \obar{C} to \obar{I} in this scenario. The locality condition that we devise for \textsc{agree} should have these consequences.

In general terms, we want this locality condition to prevent some \obar{X} from \textsc{agree}ing past another \obar{Y} in the way that the Head Movement Constraint (i.e. (\ref{ex4:hmf})) describes. But we don't want that locality condition to hold when \obar{X} is \textsc{agree}ing with a phrase, as when $\phi$ or case features are involved. In those instances, \textsc{agree} can cross several heads, as we saw in (\ref{ex4:longagr}), repeated here.
\begin{exe}
\exr{ex4:longagr}
\small
\begin{forest}
[TP[\ibar{T}
	[\obar{T}\\ case: \emph{nom}, name=probe ]
	[VP [\ibar{V}
		[\obar{V} [have] ]
		[IP [\ibar{I}
			[\obar{I} [perfect] ]
			[vP 
				[DP\\ case:, name=goal ]
				[\ibar{v} [bark, roof] ]
			 ]
		] ]
	] ]
]]
\draw[|-|, dotted] (goal) to[out=south west,in=south west] node[pos=0.5,xshift=-5mm]{\emph{agree}} (probe);
\end{forest}
\end{exe}

There is a popular way of framing the locality condition so that it has these consequences that builds upon work by Luigi Rizzi.\cite{Rizzi:1990} The central observation behind Rizzi's idea is that what distinguishes the cases where an intervening head matters from those where it doesn't is whether the intervening head could have satisfied the \textsc{agree}ment relation or not. \textsc{agree}ment from \obar{C} to \obar{I} in (\ref{ex4:hmcitcd}) is blocked because \obar{C}'s force feature could have agreed with a term in \obar{T}. That is why the grammatical sentence associated with this D-structure is~(\ref{ex4:hmcitca}).
\begin{exe}
	\raggedright
	\ex \label{ex4:hmcitca}
	Should she have been eating nattoo?
\end{exe}
And \textsc{agree}ment between \obar{T} and the main verb across \obar{Neg} in (\ref{ex4:neghmill}) is blocked because \obar{T} can \textsc{agree} with \obar{Neg}, as indicated by the case of Head Movement involving negated auxiliary verbs:
\begin{exe}
	\raggedright
	\ex \label{ex4:negauxd}
	\small
	\begin{forest}
	[TP
		[DP [she, roof] ]
		[\ibar{T}
			[\obar{T} 
				[\obar{Neg}, name=binder 
					[\obar{V} [$\surd$have\\ person: \emph{\nth{3}}\\ number: \emph{sing}\\ tense: \emph{pres}] ]
					[\obar{Neg} [neg] ]
				]
				[\obar{T}\\ $\bullet$person: \emph{\nth{3}}\\ $\bullet$number: \emph{sing}\\ tense: \emph{pres} ]
			 ]
			[NegP [\ibar{Neg}
				[{\ }, name=trace ]
				[VP [eaten nattoo, roof] ]
			] ]
		 ]
	]
	\draw[->,>=stealth', ultra thin, gray] (trace) to[out=south west,in=south,looseness=1.3] node[pos=0.75,xshift=-10mm]{} (binder);
	\end{forest}
\end{exe}
By contrast, in (\ref{ex4:longagr}), by contrast, it's the Case feature which is driving the \textsc{agree} relation, and Case features are not found on heads. They are found only on argument DPs. The intervening heads do not matter in this instance to the \textsc{agree} relation, therefore, as it applies to Case. Rizzi's idea was to relativize the locality condition to the particular kind of \textsc{agree} relation involved. He called his condition ``Relativized Minimality.''

Rizzi framed his condition as a constraint on movement since the cases he examined were all ones in which movement was triggered by the \textsc{agree} relation. Translated into a condition on \textsc{agree}, it is often called ``closest attract.'' I will frame the condition in the informal way in~(\ref{ex4:rm}).
\begin{exe}
	\raggedright
	\ex \label{ex4:rm} Closest Attract\\
	If \textsc{agree} holds between X and Y valuing $f$, then there can be no Z, such that \textsc{agree} for $f$ could have held between Z and X and~Y.
\end{exe}

\section{Summary}

Here is a brief recap of the elements we have introduced to our grammar in this chapter.
\begin{exe}
\ex \label{ex4:recap}
\begin{xlist}
	\ex The EPP\\
	Let $\alpha$ and $\beta$ be features in an \textsc{agree} relation, and let $A$ immediately dominate $\alpha$ and $B$ immediately dominate $\beta$. If $\alpha$ or $\beta$ have $\bullet$, then:
	\begin{xlist}
		\ex there must be a projection of $A$ that dominates $B$, or
		\ex there must be a projection of $B$ that dominates $A$.
	\end{xlist}
	\ex If a language has an unvalued Case feature or unvalued $\phi$ features, then there will be some head in which they are bundled.
	% \pagebreak
	\ex \textsc{agree}\\
	If X c-commands Y, then an unvalued feature in X or Y can get the value of the matching features in the other just in case X and Y \textsc{agree}. X and Y \textsc{agree} iff:
\begin{xlist}
	\ex X is a head, and
	\ex There is no Z that c-commands~Y, and X could have \textsc{agree}d with that feature in~Z.
\end{xlist}
\ex Stray Feature Filter\\[3pt]
The values of a feature must be morphologically expressed and every feature must get a value by S-structure.
\end{xlist}
\end{exe}

